\part{Phase and status history}\label{cr:part:phase-and-status-history}

\chapter{Phase history}\label{cr:ch:phase-history}

This chapter holds the completion record of each phase.

% Converted from docs/readme/PHASE_HISTORY.md by tools/docs/md2tex.py; edit this file, not the original.
\section{Phase History}\label{cr:status:phase-history}

\begin{record}

Detailed completion records for each project phase. For the high-level
map and current status, see the root \src{README.md} and
\cref{cr:status:phase-status} (\cref{cr:status:phase-status}).

\end{record}

\subsection{\texorpdfstring{Phase 0 Complete \ok{}
(2026-01-18)}{Phase 0 Complete  (2026-01-18)}}\label{cr:status:phase-history:phase-0-complete-white-check-mark-2026-01-18}

\begin{record}

\begin{itemize}
\tightlist
\item
  \ok{} All 7 architectural specifications finalized and approved
\item
  \ok{} Python reference models implemented and validated (66/66 tests
  passing)
\item
  \ok{} cocotb test infrastructure ready
\end{itemize}

\end{record}

\subsection{\texorpdfstring{Phase 1 Complete \ok{}
(2026-02-13)}{Phase 1 Complete  (2026-02-13)}}\label{cr:status:phase-history:phase-1-complete-white-check-mark-2026-02-13}

\begin{record}

\begin{itemize}
\tightlist
\item
  \ok{} All 8 RTL modules implemented (\textasciitilde1,900 lines)
\item
  \ok{} All 9/9 verification exit criteria met
\item
  \ok{} 37/37 RV32I instructions verified
\item
  \ok{} 10,000 random instructions, 0 failures
\item
  \ok{} AXI protocol compliance verified (11/11 tests)
\item
  \ok{} Debug interface complete (6/6 tests)
\item
  \ok{} Full coverage: 37/37 instructions, 8/8 FSM states
\item
  \ok{} Archived to \src{micro_p/} directory
\end{itemize}

\end{record}

\subsection{\texorpdfstring{Phase 2 Complete \ok{}
(2026-03-08)}{Phase 2 Complete  (2026-03-08)}}\label{cr:status:phase-history:phase-2-complete-white-check-mark-2026-03-08}

\begin{record}

\begin{itemize}
\tightlist
\item
  \ok{} Architecture approved (5-stage pipeline + interrupts + CSR)
\item
  \ok{} All 14 RTL modules implemented
\item
  \ok{} Comprehensive verification: 111/111 tests passing (all 7 suites)
\item
  \ok{} Random regression: 50,000 instructions (500 seeds × 100), 0
  failures
\item
  \ok{} IRQ latency: ≤3 cycles from assertion to
  trap\_\allowbreak{}taken
\item
  \ok{} Backend: 75 MHz achieved on Sky130 130nm
\end{itemize}

\end{record}

\subsection{\texorpdfstring{Phase 3 Complete \ok{}
(2026-05-21)}{Phase 3 Complete  (2026-05-21)}}\label{cr:status:phase-history:phase-3-complete-white-check-mark-2026-05-21}

\begin{record}

\begin{itemize}
\tightlist
\item
  \ok{} L1 I-Cache (4 KB, direct-mapped) + L1 D-Cache (4 KB, write-back
  + write-allocate)
\item
  \ok{} FENCE.I: 1-cycle I-cache invalidation
\item
  \ok{} Cache arbiter: D\$ priority over I\$ (D-write \textgreater{}
  D-read \textgreater{} I-read)
\item
  \ok{} Verification: I-cache 7/7, D-cache 8/8, integration 5/5;
  \textbf{139/139 full regression}
\item
  \ok{} PPA sign-off (ASAP7 7nm, Run 43, 2026-05-20): \textbf{1418 MHz /
  27.27 mW / 3,844 µm²}, 0 DRC, 0 antenna
\end{itemize}

\end{record}

\subsection{\texorpdfstring{Phase 4 Complete \ok{}
(2026-05-27)}{Phase 4 Complete  (2026-05-27)}}\label{cr:status:phase-history:phase-4-complete-white-check-mark-2026-05-27}

\begin{record}

\begin{itemize}
\tightlist
\item
  \ok{} 9/9 GPU-Lite RTL modules (\src{rtl/gpu/}): top, command queue,
  warp scheduler, compute unit, vector regfile, vector ALU, memory unit,
  coalescer, shared memory
\item
  \ok{} SIMT execution: 8-lane warps, round-robin scheduler,
  single-level divergence stack, 16 KB / 32-bank shared memory
\item
  \ok{} Verification: GPU unit + kernel tests green; CPU re-gate 140/140
  + 100k random; 1,000-kernel random regression, 0
  deadlocks/\allowbreak{}mismatches
\item
  \ok{} PPA sign-off (ASAP7 7nm, Run \code{RUN_2026-05-28_06-29-48}):
  \textbf{571 MHz (1.75 ns) / 262 mW / 115,600 µm² die / 60,500 µm²
  stdcell}, Setup WS +197.3 ps, Hold WS +16.3 ps (0 viol), 0 antenna
\end{itemize}

\end{record}

\subsection{\texorpdfstring{Phase 5 Complete \ok{}
(2026-06-24)}{Phase 5 Complete  (2026-06-24)}}\label{cr:status:phase-history:phase-5-complete-white-check-mark-2026-06-24}

\begin{record}

\begin{itemize}
\tightlist
\item
  \ok{} M1--M12 complete: AXI4 crossbar and AXI-Lite/\allowbreak{}APB
  fabric; Phase 3 refill FSMs upgraded to AXI4 bursts; UART, SPI, timer,
  IRQ controller and DMA; behavioral AXI4 SRAM controller;
  \code{soc_top} integration; CSR performance counters; SoC verification
  (boot, DMA/\allowbreak{}UART/\allowbreak{}SPI loopback, SW coherency,
  CPU-GPU IRQ); M10 L2 decision NO-GO; M11 backend flow; M12 sign-off
  (\src{CLAUDE.md} items 1--10)
\item
  \ok{} Behavioral PMU (2026-08-01) and the 2-domain multi-clock SoC
  with CDC (2026-08-09) (\src{CLAUDE.md} items 11--12)
\item
  \warn{} The M11 sign-off figures (571 MHz / 62.9 mW, fabric-only,
  excluding both hard macros) are unvalidated, and the routing claims
  are withdrawn. The current-RTL measurement is setup WNS −727.37 ps,
  781 violators, hold clean (bead \bead{0d0}). See
  \cref{cr:run:phase5}.
\end{itemize}

\end{record}

\subsection{\texorpdfstring{Phase 6 Complete \ok{}
(2026-10-05)}{Phase 6 Complete  (2026-10-05)}}\label{cr:status:phase-history:phase-6-complete-white-check-mark-2026-10-05}

\begin{record}

\begin{itemize}
\tightlist
\item
  \ok{} 6a peripherals, all APB4 slaves: GPIO (\code{ckc}, RTL
  2026-09-20, Sky130-hardened 2026-09-26), PWM (\code{f7vs.6},
  2026-09-28), WDT (\code{f7vs.7}, 2026-09-29), TRNG (\code{f7vs.8},
  2026-10-01), I2C (\code{f7vs.9}, 2026-10-03, \prref{208}). APB indices
  7--11, IRQ bits 5--9.
\item
  \ok{} 6b CRYPTO (\code{f7vs.10}, 2026-10-03): AES-128 (ECB and CTR)
  and SHA-256 at APB index 12 / IRQ bit 10; \code{test_crypto} 41/41.
  Not production cryptography (\src{CLAUDE.md} item 14).
\item
  \ok{} 6c INT8 NPU (\code{f7vs.11}, 2026-10-04): 4×4 weight-stationary
  array at APB index 13 / IRQ bit 11; \code{test_npu} 39/39,
  \code{test_soc_npu} 1/1.
\item
  \ok{} Gate A synthesis probes per item (CRYPTO 277 875 µm², NPU 203
  179 µm² stdcell) and Gate B, the batched Sky130 harden
  (\code{f7vs.15}, \prref{214}, \code{RUN_2026-10-05_06-50-11}): setup met
  at all 9 corners (worst max\_ss +0.579 ns), hold clean, Netgen LVS
  clear, routing DRC 0.
\item
  ASAP7 is out of scope for all of Phase 6 (beads \code{ma7},
  \code{lxv}, \code{2kn}). Sources:
  \src{docs/PHASE6_IP_EXPANSION_PLAN.md} §8, \cref{cr:status:phase-status}.
\end{itemize}

\end{record}

\subsection{In Progress: DFT and FPGA (status 2026-10-15)}\label{cr:status:phase-history:in-progress-dft-and-fpga-status-2026-10-15}

\begin{record}

\begin{itemize}
\tightlist
\item
  \textbf{DFT: scan, MBIST and JTAG for the Sky130 SoC} (epic
  \code{j41m}, \ghref{244}). Stage 0 architecture (\prref{250},
  \src{docs/design/DFT_ARCHITECTURE.md}); Stage 1a scan-ready SoC fabric
  (\prref{258}; bead \bead{j41m.2} still open); Stage 1b IEEE 1149.1 JTAG
  TAP (\prref{270}, \code{j41m.3}); TRST reset-sync check (\prref{280}) and CI
  \code{dft-scan-check} for both PDKs (\prref{281}). Stage 2 scan insertion
  not started; MBIST not built.
\item
  \textbf{FPGA prototyping} (epic \code{hlrm}, \ghref{245}). Stage 0
  without hardware: FPGA-define memory arms and yosys resource counts
  (\prref{279}), SRAM backdoor helper with \code{FPGA=1} regression 762
  passed (\prref{282}), board recommendation Digilent Nexys Video XC7A200T
  and ECP5 fmax lower bound (\prref{283},
  \src{docs/FPGA_BOARD_SELECTION.md}). Board-gated stages wait for a
  board.
\end{itemize}

\end{record}

\subsection{Implemented Features}\label{cr:status:phase-history:implemented-features}

\subsubsection{\texorpdfstring{Phase 1 (Complete)
\ok{}}{Phase 1 (Complete) }}\label{cr:status:phase-history:phase-1-complete-white-check-mark}

\begin{record}

\begin{itemize}
\tightlist
\item
  \textbf{ISA}: Complete RISC-V RV32I base integer instruction set (37
  instructions)
\item
  \textbf{Architecture}: Single-cycle execution with AXI stalls
\item
  \textbf{Memory Interface}: AXI4-Lite Master (unified
  instruction/\allowbreak{}data bus)
\item
  \textbf{Debug Interface}: APB3 Slave with full capabilities

  \begin{itemize}
  \tightlist
  \item
    Halt/\allowbreak{}Resume/\allowbreak{}Single-step execution
  \item
    Register file read/\allowbreak{}write access
  \item
    Program counter manipulation
  \item
    Hardware breakpoints (2 breakpoints)
  \end{itemize}
\end{itemize}

\end{record}

\subsubsection{\texorpdfstring{Phase 2 \ok{}}{Phase 2 }}\label{cr:status:phase-history:phase-2-white-check-mark}

\begin{record}

\begin{itemize}
\tightlist
\item
  \textbf{ISA}: RV32I + Zicsr (CSR instructions:
  CSRRW/\allowbreak{}S/\allowbreak{}C/\allowbreak{}I variants)
\item
  \textbf{Architecture}: 5-stage in-order pipeline
  (IF/\allowbreak{}ID/\allowbreak{}EX/\allowbreak{}MEM/\allowbreak{}WB)
\item
  \textbf{Interrupt Support}: RISC-V M-mode (timer + external
  interrupts)
\item
  \textbf{Hazard Handling}: Detection, stalling, EX/MEM/WB forwarding
\item
  \textbf{Target Performance}: 1 IPC at 200 MHz (2x Phase 1 frequency)
\item
  \textbf{AXI Arbitration}: Priority arbiter for IF vs MEM requests
\item
  \textbf{Debug Interface}: Enhanced with pipeline drain support
\end{itemize}

\end{record}

\chapter{Status history}\label{cr:ch:status-history}

This chapter holds the earlier project status statements.

% Converted from docs/PHASE_STATUS.md by tools/docs/md2tex.py; edit this file, not the original.
\section{Project Phase Status}\label{cr:status:phase-status}

\begin{record}

Last updated: 2026-10-15

\end{record}

\subsection{Current Phase}\label{cr:status:phase-status:current-phase}

\begin{record}

\textbf{Phase 5: SoC Integration} - \ok{} COMPLETE (2026-06-24) ---
M1--M12 done. M10 L2 gate → NO-GO
(\src{docs/M10_L2_DECISION_ANALYSIS.md}); M11 ASAP7 SoC P\&R signed off
\textbf{571 MHz / 62.9 mW / 520×520 µm / 65.6 \% util / 0 DRC / 0
antenna} (sv2v frontend; \cref{cr:run:phase5}). Indicative
ASAP7 (predictive PDK). \warn{} \textbf{Timing/\allowbreak{}power now
also unvalidated} (bead \bead{8f3}/\code{0p6}, 2026-09-16): produced by
\code{OpenROAD.STAMidPNR-3} without an in-session re-route, a defect
proven to give false-clean zero-wire STA; run 14\textquotesingle s
artifacts are wiped and cannot be re-checked. \ok{} \textbf{Update
2026-09-18 (bead \bead{0d0})}: the SoC on \emph{current} RTL has since
been re-measured three times on the fixed flow
(\code{STA_POSTGRT_INSESSION_GRT=1}, pinned OpenROAD 26Q2,
annotation-gated). Current honest point, superseding the earlier 258.5
mW / −2084 ps snapshot: \textbf{setup WNS −727.37 ps / TNS −60,665 ps /
781 violators, hold CLEAN, 283.6 mW total / 89.6 mW fabric-only}
(\code{RUN_2026-09-17_20-41-11}). Timing is still NOT closed at 1750 ps,
and run 14\textquotesingle s own numbers remain unvalidated ---
\cref{cr:run:phase5} §"Power figure caveat".

\textbf{Phase 6: IP Expansion} - \ok{} COMPLETE (2026-10-05) --- 6a
(GPIO, PWM, WDT, TRNG, I2C), 6b (CRYPTO) and 6c (NPU) RTL complete; epic
\code{f7vs} closed (15/15 children). The batched Sky130 SoC harden (Gate
B, \code{f7vs.15}, \prref{214}) closed timing at all 9 corners on
2026-10-05. ASAP7 stays out of scope for Phase 6 (beads
\code{ma7}/\code{lxv}/\code{2kn}). Details in the Phase 6 section below.

\textbf{DFT: scan + MBIST + JTAG for the Sky130 SoC} (epic \code{j41m},
\ghref{244}) - \wip{} IN PROGRESS --- Stage 0 architecture \ok{}
(\src{docs/design/DFT_ARCHITECTURE.md}, \prref{250}); Stage 1a scan-ready
SoC fabric merged (\prref{258}; bead \bead{j41m.2} still open); Stage 1b
IEEE 1149.1 JTAG TAP \ok{} (\prref{270}, \code{j41m.3}). Stage 2 scan
insertion not started; MBIST not built (Stage 4 plan only).

\textbf{FPGA prototyping} (epic \code{hlrm}, \ghref{245}) - \wip{} IN
PROGRESS --- Stage 0 \ok{} without hardware:
\src{docs/FPGA_STAGE0_EVALUATION.md} (\prref{279}) and
\src{docs/FPGA_BOARD_SELECTION.md} (\prref{283}; recommendation Digilent
Nexys Video XC7A200T). Open-tool ECP5 fmax lower bound:
\textasciitilde48--52 MHz CPU-only on LFE5U-85F (a lower bound, not an
estimate of the Artix-7). Board-gated stages not started. A Vivado
Artix-7 run is in progress on a separate branch with no results yet.

\textbf{Phase 7: Mixed-Signal PLL Clock Generator} - \wip{} IN PROGRESS
--- M-a/\allowbreak{}M-b1/\allowbreak{}M-b2/\allowbreak{}M-c \ok{}
complete (2026-06-20); M-d = documentation (this update). Dual-PDK
charge-pump PLL via analog-design agents (ASAP7 indicative + Sky130 real
CP-block DRC/LVS), AMS RNM integrated as SoC clock source. See
\src{docs/PHASE7_MIXED_SIGNAL_PLL_PLAN.md}. (Phase 7 ran alongside the
in-progress Phase 5 PD/\allowbreak{}sign-off tail.)

\begin{itemize}
\tightlist
\item
  M1--M8 \ok{}: AXI4 crossbar + AXI-Lite ring, cache burst upgrade,
  peripherals
  (UART/\allowbreak{}SPI/\allowbreak{}timer/\allowbreak{}IRQ), DMA,
  behavioral SRAM, perf counters, SoC top.
\item
  M9 \ok{} SoC verification: boot 100/100; DMA+UART+SPI loopback; SW
  coherency (D\$ flush→GPU→D\$ inval); CPU-GPU IRQ integration; DUT-side
  boot SRAM check. \code{soc_all} 73/73 at M9 (now 120/120 --- see the
  PMU entry below); 1M+ cycle stress (1,079,867 cyc, 0 fail). Two RTL
  bugs found+fixed: D-cache MMIO caching (\code{go9}) and
  axi4\_\allowbreak{}crossbar AR/AW handshake+arbitration (\code{7fs}).
\end{itemize}

\textbf{Phase 6a: GPIO controller (bead \bead{ckc})} - \ok{} COMPLETE
(2026-09-26; RTL 2026-09-20, Sky130 harden 2026-09-26). First Phase 6
peripheral. New \src{rtl/periph/gpio_controller.sv} --- APB4 slave on
\code{apb4_register_bank}, 32 pins.

\begin{itemize}
\tightlist
\item
  \textbf{Bus}: APB4, per \src{docs/PERIPHERAL_BUS_EVALUATION.md}
  ("Phase 6 peripherals become trivial APB drop-ins"). The "AXI4-Lite"
  wording in \cref{cr:status:roadmap} and \src{CLAUDE.md} was stale and has
  been corrected.
\item
  \textbf{Integration}: APB slave index 7 at
  \texttt{0x2000\_\allowbreak{}A000–AFFF}; \code{APB_N_SLAVES} 7 → 8;
  \code{AXIL_APB_LIMIT} and \code{PERIPH_LIMIT} both extended
  \code{0x2000_9FFF} → \code{0x2000_AFFF} --- the same two-edit pattern
  used for PLL, PMU and PLL2. New IRQ source at
  \code{interrupt_controller} bit 5 (\code{N_SOURCES} 5 → 6).
\item
  \textbf{Pins}: unidirectional triplet \code{gpio_out_o[31:0]} /
  \code{gpio_oe_o[31:0]} / \code{gpio_in_i[31:0]} on \code{soc_top}.
  There is \textbf{no tristate anywhere in this RTL tree} and the Sky130
  SoC hardens as a core macro with no pad ring, so the bidirectional
  merge is deliberately a pad-ring integration concern. \code{gpio_in_i}
  is asynchronous and is synchronised per pin by \code{N_PINS} separate
  \texttt{cdc\_\allowbreak{}sync\ \#(.\allowbreak{}WIDTH(1))} instances
  --- not one wide instance, because that primitive\textquotesingle s
  scope warning forbids multi-bit binary buses.
\item
  \textbf{Interrupts}: per-pin level/\allowbreak{}edge with polarity.
  \code{GPIO_IRQ_STAT} is \textbf{edge-sticky / level-live} (ARM PL061
  semantics): edge-mode bits latch until a \code{GPIO_IRQ_CLR} write,
  with a set beating a same-cycle clear so an edge is never lost to a
  racing clear; level-mode bits track their condition live and ignore
  \code{GPIO_IRQ_CLR} entirely. A fully sticky raw register was rejected
  because the reset defaults (level, active-low, pins low) make every
  pin\textquotesingle s condition true, latching all-ones out of reset.
  \code{GPIO_IRQ_EN} masks the output only. \code{irq_o} is level-held,
  which the 2-FF CPU-domain crossing at \code{soc_top.sv:623-627}
  requires.
\item
  \textbf{Verification}: \code{test_gpio} \textbf{16/16} (reset defaults
  incl. the deliberate all-ones \code{GPIO_IRQ_STAT} at
  level/\allowbreak{}active-low defaults, RW round-trip, output/OE
  mirroring, the cycle-exact 3-edge input-sync latency,
  rising/\allowbreak{}falling edge capture, level tracking, edge
  stickiness + clear, set-beats-same-cycle-clear,
  clear-no-effect-on-level, \code{pstrb} partial-word clear,
  \code{GPIO_IRQ_EN} masking the output only, \code{irq_o} held 50
  cycles, \code{GPIO_IRQ_CLR} reads-as-0, out-of-range access).
  \code{soc_all} \textbf{217/217} across 27 suites, up from 201.
  Extending the APB window broke
  \code{test_axil_interconnect.test_decerr_unmapped}, which probed
  \code{0x2000_A000} expecting DECERR --- now the GPIO slot;
  \code{BAD_HIGH} bumped to \code{0x2000_B000}, the maintenance step
  that test\textquotesingle s own comment prescribes for each window
  growth.
\item
  \textbf{Hardened 2026-09-26} (\code{RUN_2026-09-26_00-07-59}, bead
  \bead{00ef}): \textbf{Netgen LVS PASSED} ("Circuits match uniquely", 0
  errors), routing DRC 0, \textbf{setup AND hold both clean at all 9
  corners} (+7.08 ns / +0.1413 ns), 48.25 mW, 38.98 \% util, die
  unchanged 6700 × 3100 µm. Magic DRC 9,081 is the only remaining
  deferred error (bead \bead{45a} waiver). First Sky130 SoC hardening
  since 2026-07-30 --- \warn{} NOT single-variable vs. that baseline (33
  RTL commits between, 13 touching \code{soc_top.sv}). \textbf{The ss
  setup failure that bead \bead{ujv} was deferred on is gone} on current
  RTL. KLayout DRC deliberately skipped (8.5 GB OOM on this 15 GB host,
  and it sits before LVS in the flow) --- deferred, not waived. Hold
  went \textbf{71 → 3 → 0}: resizer hold margins 0.05 → 0.3 closed 68,
  and
  \texttt{set\_\allowbreak{}false\_\allowbreak{}path\ -\allowbreak{}from\ [get\_\allowbreak{}ports\ gpio\_\allowbreak{}in\_\allowbreak{}i]}
  closed the 3 survivors (all one path, \code{gpio_in_i[31]} → a
  \code{cdc_sync} stage-1 flop). Detail:
  \src{docs/SKY130_REAL_DRC_LVS_EVALUATION.md}.
\end{itemize}

\textbf{Sky130 PD quality campaign, 2026-09-26} --- four runs after the
GPIO harden, each single-variable:

{\small\setlength{\tabcolsep}{4pt}

{
\begin{longtable}[]{@{}
  >{\raggedright\arraybackslash}p{\dimexpr(\linewidth-6\tabcolsep)*55/1000\relax}
  >{\raggedright\arraybackslash}p{\dimexpr(\linewidth-6\tabcolsep)*419/1000\relax}
  >{\raggedright\arraybackslash}p{\dimexpr(\linewidth-6\tabcolsep)*316/1000\relax}
  >{\raggedright\arraybackslash}p{\dimexpr(\linewidth-6\tabcolsep)*190/1000\relax}@{}}
\toprule\noalign{}
\begin{minipage}[b]{\linewidth}\raggedright
Bead
\end{minipage} & \begin{minipage}[b]{\linewidth}\raggedright
Change
\end{minipage} & \begin{minipage}[b]{\linewidth}\raggedright
Result
\end{minipage} & \begin{minipage}[b]{\linewidth}\raggedright
Cost
\end{minipage} \\
\midrule\noalign{}
\endhead
\bottomrule\noalign{}
\endlastfoot
\code{00ef} & \code{set_false_path} on \code{gpio_in_i} +
\code{cpu_rst_n_i} & hold 3 → \textbf{0} at all 9 corners & --- \\
\code{e45j} & design-repair margins 10→30 / 20→40 \% & max-slew 8623 →
\textbf{2380}, cap 346 → \textbf{70} & antenna 191 → 263 pins \\
\code{58q} & \code{GRT_ANTENNA_ITERS} 8→20, \code{MARGIN} 25→50 &
antenna 263 → \textbf{182} pins, 182 → \textbf{114} nets & max-slew 2380
→ 3199 \\
\end{longtable}
}

}

\warn{} \textbf{The slew and antenna knobs pull against each other} ---
repair buffers cut slew and raise antenna; antenna diodes cut antenna
and raise slew, each diode being another load pin. Both current points
are \textbf{single-axis optima}; the Pareto frontier is demonstrated but
has never been swept jointly.

Two structural findings from that campaign, both recorded rather than
fixed:

\begin{itemize}
\tightlist
\item
  \textbf{No post-RCX repair stage exists in LibreLane.} Every
  repair/\allowbreak{}resizer step is pre-route on GRT-\emph{estimated}
  parasitics, so real-parasitic erosion appears first at post-route STA
  and last at a point where nothing can act on it (slew/cap 0 at steps
  33/40/41 → thousands at step 51). The same gap explains the hold
  regression. Adding a stage needs a shared-LibreLane edit ---
  evaluated, and recommended \emph{not yet}, in
  \src{docs/LIBRELANE_PATCHING_EVALUATION.md}.
\item
  \textbf{\texttt{max\ fanout\ violations\ =\allowbreak{}\ 0} is
  vacuous} --- \code{sky130_soc.sdc} has no \code{set_max_fanout} at
  all, so a 211-fanout net reports clean. Do not cite that metric until
  a real constraint exists.
\end{itemize}

\textbf{LibreLane install state captured 2026-09-26}
(\src{pnr/librelane_patches/}). The shared install carried 15 modified
files (+641/\allowbreak{}−89) as \emph{uncommitted} edits, and the
pre-existing patch archive lived under the \textbf{gitignored}
\src{memory/} --- so no PD figure in this repo was reproducible from a
clean checkout, and the exact tool state behind every number existed on
one disk, unbacked-up. Now tracked, with sha256s and a verified
(\texttt{git\ apply\ -\allowbreak{}-\allowbreak{}check\ -\allowbreak{}-\allowbreak{}reverse})
full diff. \warn{} One patch in that set,
\code{drt.tcl}\textquotesingle s catch-and-continue around
\code{detailed_route}, is why every ASAP7 "0 DRC / 0 antenna" claim was
withdrawn (bead \bead{xy6}); re-examining it is the next step.

\textbf{Bead \code{ujv} re-examined 2026-09-26 --- its ss setup failure
NO LONGER REPRODUCES.} At an unchanged 25.0 ns period, max\_ss went
\textbf{−4.328 ns FAIL → +8.097 ns PASS} (nom\_ss −2.885 → +9.232,
min\_ss −1.120 → +10.069), a +12.4 ns swing. Its parking condition ("do
not chase ss until hold is green at all 9 corners") is also now
satisfied. \warn{} But the \textbf{root cause is bypassed, not fixed}:
\code{ujv} blamed the CPU macro\textquotesingle s 8.155 ns ss clk→Q, and
\code{u_cpu} now appears \textbf{zero times} in
max\_ss\textquotesingle s 1000 worst paths while the macro views are
unchanged since 2026-07-27 --- so that weakness is latent and would
resurface at a shorter period. \textbf{Opportunity:} +8.1 ns of ss
margin now exists where −4.3 ns did, so real frequency headroom above 40
MHz is likely and has never been swept. The ROADMAP\textquotesingle s 75
MHz Sky130 figure for GPIO remains a projection. Sky130 SoC headroom at
the last sign-off was 37.9 \% utilisation on a 6700 × 3100 µm die, with
the tightest margin max\_tt setup +0.339 ns.

\textbf{Pre-Phase-6 \#5: behavioral PMU (GH epic \#98)} - \ok{} COMPLETE
(2026-08-01). New \src{rtl/soc/pmu.sv} --- APB4 slave (reusing
\code{apb4_register_bank}) + a per-domain sequencer FSM encoding the 4
PST states from \src{pnr/constraints/phase5_soc.upf}
(NORMAL/\allowbreak{}CPU\_\allowbreak{}OFF/\allowbreak{}GPU\_\allowbreak{}OFF/\allowbreak{}IDLE).

\begin{itemize}
\tightlist
\item
  \textbf{Sequencing order is fixed by the UPF, not chosen}: power-down
  \texttt{ret\_\allowbreak{}save\ →\ iso\_\allowbreak{}en\ →\ gate\ clock\ →\ assert\ reset};
  power-up
  \texttt{release\ reset\ →\ ungate\ clock\ →\ ret\_\allowbreak{}restore\ →\ de-\allowbreak{}assert\ iso\_\allowbreak{}en}.
  Each domain walks an explicit 9-state chain, one action per cycle, so
  the order is structurally guaranteed and directly assertable.
\item
  \textbf{Integration} (\#100): no free APB slot existed --- every 4 KB
  slot from \code{0x2000_2000} to \code{0x2000_7FFF} was taken --- so
  \code{PERIPH_LIMIT} and \code{AXIL_APB_LIMIT} were both extended to
  \code{0x2000_8FFF}, with PMU at
  \texttt{0x2000\_\allowbreak{}8000–8FFF} as APB slave 5. Two
  \code{rv32i_clock_gate} cells now gate \code{core_clk} into
  \code{u_cpu}/\code{u_gpu} --- the first SoC-level clock gates.
  Isolation clamps sit in \code{soc_top} outside the gated domains, per
  the UPF\textquotesingle s \texttt{-\allowbreak{}location\ parent}.
\item
  \textbf{Scope limit}: this is a \emph{sequencing} controller.
  Retention save/\allowbreak{}restore are emitted in the correct order
  for verification to assert but are a functional no-op --- the UPF
  (rewritten 2026-10-17, bead \bead{h3lm}) declares no retention and no
  power switches. No real supply removal, power switches, or retention
  silicon; \code{pg_ctrl} is deliberately not emitted.
\item
  \textbf{Verification} (\#101): \code{soc_all} \textbf{120/120} (was
  82/82), \code{test_pmu} 10/10 with per-cycle ordering assertions in
  both directions. Reset defaults leave both domains
  on/\allowbreak{}un-isolated/\allowbreak{}un-reset, so the SoC is
  functionally identical to pre-PMU for firmware that never writes CTRL.
\end{itemize}

\textbf{Pre-Phase-6 \#4: 2-domain multi-clock SoC + CDC (GH epic \#90)}
- \ok{} COMPLETE (2026-08-09). The SoC is split into a CPU domain and a
GPU+bus+peripheral domain, giving exactly \textbf{one} CDC boundary (CPU
↔ crossbar M0).

\begin{itemize}
\tightlist
\item
  \textbf{PD result} (\#96, run 23 \code{RUN_2026-08-09_14-36-16}):
  \code{cpu_clk} 780 ps at \textbf{+9.8 ps} WS, TNS 0, 0 violators →
  \textbf{1282 MHz MET}; \code{sys_clk} 1750 ps at \textbf{+74.8 ps} WS,
  TNS 0, 0 violators → \textbf{571 MHz MET}. 0 routing DRC, 0 antenna, 0
  slew/cap, 51.9 mW, 66.9 \% util, die unchanged 520 × 520 µm,
  post-route CDC budget check PASS. \textbf{The CPU runs at 2.24× the
  fabric clock and is no longer de-rated onto the 1750 ps cycle} --- the
  point of the epic. Run 14 remains the accepted single-clock M11
  sign-off and is byte-reproducible; the multi-clock flow runs off
  \src{pnr/asap7/soc/config_multiclock.json} +
  \code{constraints/phase5_soc_multiclock.sdc}.
\item
  \textbf{RTL} (\#91/\allowbreak{}\#92/\allowbreak{}\#93):
  \src{rtl/soc/async_axi_fifo.sv} (dual-clock AXI4 bridge on new
  \src{rtl/soc/cdc/} primitives), a second \code{pll_subsystem} +
  \code{APB_PLL2} slot, and \src{rtl/soc/apb_cdc_bridge.sv}.
  \code{soc_all} 159/159; \code{soc_multiclock} 4/4 at a \textbf{7 ns /
  3 ns coprime ratio} --- the first real CDC coverage in this repo, and
  it found two RTL bugs invisible to every 1:1 suite (both a \emph{reset
  value that advertises availability}).
\item
  \textbf{Two PD root causes, both now fixed and documented}: the CPU
  macro Liberty exported \textasciitilde600--800 ps combinational in→out
  arcs on its APB outputs (fixed by making all three outputs pure
  registered), and the macro\textquotesingle s boundary pins
  free-floated between regenerations --- \textbf{393 of 401 moved},
  worth 535 ps of setup --- fixed by \src{pnr/asap7/cpu/pin_order.cfg}
  (\src{pnr/asap7/cpu/README_pin_order.md}).
\item
  \warn{} \textbf{Do not quote 62.9 → 51.9 mW as a multi-clock power
  win.} \code{report_power} attributes 0.00 W to the macros, and the
  drop coincides with the first-ever ICG insertion
  (\code{USE_ICG_CELL}). ASAP7 IR drop remains unobtainable ---
  \code{analyze_power_grid} fails \code{PSM-0069} on both rails from the
  M1-only tap-cell connectivity artifact. Both tracked as beads.
\item
  \warn{} \textbf{The run-23 timing figures above (both domains MET) are
  also unvalidated} (bead \bead{8f3}/\code{0p6}): produced by the stock,
  non-re-routing \code{OpenROAD.STAMidPNR-3} step, proven elsewhere to
  report false-clean zero-wire STA. Run 23\textquotesingle s artifacts
  are wiped; not re-verifiable directly. \ok{} The honest
  \emph{current-RTL} point (2026-09-18, bead \bead{0d0}) is setup WNS
  −727.37 ps / 781 violators / hold clean / 283.6 mW at
  \code{RUN_2026-09-17_20-41-11} --- a different, still-unclosed design
  point, not a re-run of run 23. See \cref{cr:run:phase5}
  §"Power figure caveat".
\item
  Full record: \cref{cr:run:phase5} (\cref{cr:run:phase5})
  and \src{docs/3PLL_CDC_EVALUATION.md}.
\end{itemize}

\textbf{Sky130 real DRC/LVS sign-off (GH epic \#102)} --- Stages 1--2
\ok{} complete (2026-07-31); Stages 3--4 \hold{} host-gated.

\begin{itemize}
\tightlist
\item
  \textbf{Stage 1} (CPU standalone, \ghref{103}): KLayout DRC 0, Netgen LVS
  MATCH, setup +0.366 ns @ 75 MHz. Magic DRC 27.7 M waived --- 100 \%
  inside SRAM macro footprints, PDK artifact (bead \bead{45a}).
\item
  \textbf{Stage 2} (CPU macro + peripherals SoC, no GPU, \ghref{104}): LVS
  PASSED, routing DRC 0, PDN 0, hold clean, 40 MHz, nom\_tt 37.63 mW,
  226,952 stdcells, die 6700 × 3100 µm. Setup fails at ss (bead
  \bead{ujv}) and antenna leaves a 140/120 residual (bead \bead{58q})
  --- both documented, deferred.
\item
  \warn{} \textbf{Timing scope: typical-corner (nom\_tt) only.} The SRAM
  macro is characterized at TT alone, so the nine reported corners are
  nine labels backed by one macro model. The physical gates
  (LVS/\allowbreak{}DRC/\allowbreak{}PDN) are not affected and stand as
  real results. Closing ss honestly needs per-corner SRAM SPICE
  characterization, measured at \textasciitilde80--95 h on a host that
  reboots every 2--8 h → bead \bead{o1i}, deferred host-blocked.
\item
  \textbf{Stages 3--4} (GPU standalone + full SoC incl. GPU, \ghref{105}/\allowbreak{}\#106): not attempted; need ≥32 GB RAM +
  \textasciitilde500 GB scratch. GPU stays ASAP7-indicative-only.
\item
  Full record: \src{docs/SKY130_REAL_DRC_LVS_EVALUATION.md} §7.
\end{itemize}

\textbf{Previous Phases}:

\begin{itemize}
\tightlist
\item
  Phase 4 (GPU-Lite SIMT Compute Engine) - \ok{} COMPLETE (2026-05-27)
  --- all GPU tests green, ASAP7 ≥500 MHz sign-off
\item
  Phase 3 (Memory System \& Caches) - \ok{} COMPLETE (2026-05-21) ---
  all 20 cache tests passing, 139/139 total
\item
  Phase 2 (Pipelined CPU) - \ok{} COMPLETE (2026-03-08) --- 75 MHz on
  Sky130 130nm
\item
  Phase 1 (Minimal RV32I Core) - \ok{} COMPLETE (2026-02-13)
\item
  Phase 0 (Foundations) - \ok{} COMPLETE (2026-01-18)
\end{itemize}

\end{record}

\subsection{Phase Progress}\label{cr:status:phase-status:phase-progress}

\subsubsection{\texorpdfstring{Phase 0: Foundations
\ok{}}{Phase 0: Foundations }}\label{cr:status:phase-status:phase-0-foundations-white-check-mark}

\begin{record}

\textbf{Status}: Specifications and implementation complete

\textbf{Completed}:

\begin{itemize}
\tightlist
\item
  \ok{} RV32I subset defined
  (PHASE0\_\allowbreak{}ARCHITECTURE\_\allowbreak{}SPEC.md)
\item
  \ok{} Reset behavior specified
\item
  \ok{} Trap behavior specified
\item
  \ok{} Memory ordering rules specified
\item
  \ok{} Commit semantics specified
\item
  \ok{} Project structure created (rtl/, tb/, sim/, docs/)
\item
  \ok{} Interface specifications (RTL\_\allowbreak{}DEFINITION.md) -
  AXI4-Lite \& APB3 protocols defined
\item
  \ok{} GPU architecture specification
  (PHASE4\_\allowbreak{}GPU\_\allowbreak{}ARCHITECTURE\_\allowbreak{}SPEC.md)
  - complete
\item
  \ok{} Reference model specification
  (REFERENCE\_\allowbreak{}MODEL\_\allowbreak{}SPEC.md) - complete
\item
  \ok{} Memory map specification (MEMORY\_\allowbreak{}MAP.md) -
  complete with 4 KB alignment
\item
  \ok{} Phase-aligned verification plan
  (VERIFICATION\_\allowbreak{}PLAN.md) - structured by phases
\item
  \ok{} Phase 1 CPU specification
  (PHASE1\_\allowbreak{}ARCHITECTURE\_\allowbreak{}SPEC.md) - complete
\item
  \ok{} Python reference model implementation - 66/66 tests passing
  (memory, RV32I, GPU)
\item
  \ok{} cocotb test infrastructure setup - complete with BFMs,
  scoreboard, utilities, documentation
\end{itemize}

\textbf{Remaining}:

\begin{itemize}
\tightlist
\item
  None - Phase 0 complete!
\end{itemize}

\textbf{Exit Criteria Status}:

\begin{itemize}
\tightlist
\item
  \ok{} Written ISA + microarchitecture spec
  (PHASE0\_\allowbreak{}ARCHITECTURE\_\allowbreak{}SPEC.md exists)
\item
  \ok{} No RTL yet (confirmed - directories empty)
\item
  \ok{} Supporting specifications complete (all 7 specs finalized)
\item
  \ok{} Python reference model matches specification (66/66 tests
  passing)
\item
  \ok{} Test infrastructure ready (cocotb infrastructure complete)
\item
  \ok{} Final human specification review complete (approved 2026-01-18)
\end{itemize}

\textbf{Target Completion}: 2026-01-31

\end{record}

\subsubsection{\texorpdfstring{Phase 1: Minimal RV32I Core \ok{}
VERIFICATION
COMPLETE}{Phase 1: Minimal RV32I Core  VERIFICATION COMPLETE}}\label{cr:status:phase-status:phase-1-minimal-rv32i-core-white-check-mark-verification-complete}

\begin{record}

\textbf{Status}: RTL implementation complete, all 9/9 verification exit
criteria met (2026-02-13)

\textbf{Prerequisites}: \ok{} Phase 0 exit criteria met

\textbf{RTL Modules} (8/8 complete, \textasciitilde1,900 lines total):

{
\begin{longtable}[]{@{}
  >{\raggedright\arraybackslash}p{\dimexpr(\linewidth-4\tabcolsep)*374/1000\relax}
  >{\raggedright\arraybackslash}p{\dimexpr(\linewidth-4\tabcolsep)*484/1000\relax}
  >{\raggedright\arraybackslash}p{\dimexpr(\linewidth-4\tabcolsep)*121/1000\relax}@{}}
\toprule\noalign{}
\begin{minipage}[b]{\linewidth}\raggedright
Module
\end{minipage} & \begin{minipage}[b]{\linewidth}\raggedright
File
\end{minipage} & \begin{minipage}[b]{\linewidth}\raggedright
Status
\end{minipage} \\
\midrule\noalign{}
\endhead
\bottomrule\noalign{}
\endlastfoot
CPU top-level (AXI4-Lite + APB3) & \src{rtl/cpu/rv32i_cpu_top.sv} &
\ok{} Complete \\
CPU core wrapper & \src{rtl/cpu/core/rv32i_core.sv} & \ok{} Complete \\
Control FSM & \src{rtl/cpu/core/rv32i_control.sv} & \ok{} Complete \\
Instruction decoder & \src{rtl/cpu/core/rv32i_decode.sv} & \ok{}
Complete \\
ALU & \src{rtl/cpu/core/rv32i_alu.sv} & \ok{} Complete \\
Register file & \src{rtl/cpu/core/rv32i_regfile.sv} & \ok{} Complete \\
Immediate generator & \src{rtl/cpu/core/rv32i_imm_gen.sv} & \ok{}
Complete \\
Branch comparator & \src{rtl/cpu/core/rv32i_branch_comp.sv} & \ok{}
Complete \\
\end{longtable}
}

\textbf{Verification Exit Criteria} (9/9 met):

{
\begin{longtable}[]{@{}
  >{\raggedright\arraybackslash}p{\dimexpr(\linewidth-6\tabcolsep)*56/1000\relax}
  >{\raggedright\arraybackslash}p{\dimexpr(\linewidth-6\tabcolsep)*243/1000\relax}
  >{\raggedright\arraybackslash}p{\dimexpr(\linewidth-6\tabcolsep)*159/1000\relax}
  >{\raggedright\arraybackslash}p{\dimexpr(\linewidth-6\tabcolsep)*523/1000\relax}@{}}
\toprule\noalign{}
\begin{minipage}[b]{\linewidth}\raggedright
\#
\end{minipage} & \begin{minipage}[b]{\linewidth}\raggedright
Criterion
\end{minipage} & \begin{minipage}[b]{\linewidth}\raggedright
Target
\end{minipage} & \begin{minipage}[b]{\linewidth}\raggedright
Status
\end{minipage} \\
\midrule\noalign{}
\endhead
\bottomrule\noalign{}
\endlastfoot
1 & Smoke tests passing & 6/6 & \ok{} MET --- 6/6 with scoreboard \\
2 & Scoreboard mismatches & 0 & \ok{} MET --- 0 mismatches \\
3 & Instruction coverage & 37/37 (100\%) & \ok{} MET --- all 37 RV32I
instructions \\
4 & Random instruction tests & 10,000+, 0 fail & \ok{} MET --- 10,000
instructions, 0 failures \\
5 & AXI protocol tests & 100\% pass & \ok{} MET --- 11/11 tests
passing \\
6 & Debug interface tests & 100\% pass & \ok{} MET --- 6/6 tests
(single-step, BP0, BP1, GPR/PC write) \\
7 & Code coverage & \textgreater95\% & \ok{} MET --- Verilator annotated
reports, \texttt{make\ coverage} \\
8 & State coverage & 8/8 (100\%) & \ok{} MET --- 8/8 FSM states
covered \\
9 & Failing random seeds & 0 & \ok{} MET --- 0 failing seeds (100/100
pass) \\
\end{longtable}
}

\textbf{Key RTL Fixes Applied}:

\begin{itemize}
\tightlist
\item
  Branch/\allowbreak{}jump timing fix (registered decision flag in
  rv32i\_\allowbreak{}control.sv)
\item
  Load data latching fix (mem\_\allowbreak{}rdata\_\allowbreak{}raw
  register in rv32i\_\allowbreak{}core.sv)
\item
  Register file combinational reads (rv32i\_\allowbreak{}regfile.sv)
\end{itemize}

\textbf{Verification Milestones}:

\begin{itemize}
\tightlist
\item
  2026-01-24: Task 1 complete --- Scoreboard integration (6/6 smoke
  tests)
\item
  2026-01-26: Task 2 complete --- ISA compliance tests (37/37 passing)
\item
  2026-01-28: Task 3 complete --- Random instruction generator (10,000
  instructions, 0 failures)
\item
  2026-02-07: Task 4 complete --- AXI4-Lite protocol tests (11/11
  implemented)
\item
  2026-02-13: Tasks 5, 6, 7 complete --- Debug interface (6/6), Coverage
  (100\%), Docs
\end{itemize}

\end{record}

\subsubsection{\texorpdfstring{Phase 2: Pipelined CPU \ok{}
COMPLETE}{Phase 2: Pipelined CPU  COMPLETE}}\label{cr:status:phase-status:phase-2-pipelined-cpu-white-check-mark-complete}

\begin{record}

\textbf{Status}: \ok{} COMPLETE (2026-03-08) --- 75 MHz achieved on
Sky130 130nm

\textbf{Prerequisites}: \ok{} Phase 1 exit criteria met (2026-02-13)

\textbf{Architecture Spec}:
\src{docs/design/PHASE2_ARCHITECTURE_SPEC.md} --- APPROVED (2026-02-14),
all 7 open questions resolved

\textbf{RTL Modules} (14/14 complete):

{
\begin{longtable}[]{@{}
  >{\raggedright\arraybackslash}p{\dimexpr(\linewidth-4\tabcolsep)*330/1000\relax}
  >{\raggedright\arraybackslash}p{\dimexpr(\linewidth-4\tabcolsep)*543/1000\relax}
  >{\raggedright\arraybackslash}p{\dimexpr(\linewidth-4\tabcolsep)*107/1000\relax}@{}}
\toprule\noalign{}
\begin{minipage}[b]{\linewidth}\raggedright
Module
\end{minipage} & \begin{minipage}[b]{\linewidth}\raggedright
File
\end{minipage} & \begin{minipage}[b]{\linewidth}\raggedright
Status
\end{minipage} \\
\midrule\noalign{}
\endhead
\bottomrule\noalign{}
\endlastfoot
CPU top-level (AXI4-Lite + APB3) & \src{rtl/cpu/rv32i_cpu_top.sv} &
\ok{} Complete \\
CPU core wrapper & \src{rtl/cpu/core/rv32i_core.sv} & \ok{} Complete \\
IF pipeline stage & \src{rtl/cpu/core/pipeline/rv32i_pipeline_if.sv} &
\ok{} Complete \\
ID pipeline stage & \src{rtl/cpu/core/pipeline/rv32i_pipeline_id.sv} &
\ok{} Complete \\
EX pipeline stage & \src{rtl/cpu/core/pipeline/rv32i_pipeline_ex.sv} &
\ok{} Complete \\
MEM pipeline stage & \src{rtl/cpu/core/pipeline/rv32i_pipeline_mem.sv} &
\ok{} Complete \\
WB pipeline stage & \src{rtl/cpu/core/pipeline/rv32i_pipeline_wb.sv} &
\ok{} Complete \\
Pipeline package & \src{rtl/cpu/core/rv32i_pipeline_pkg.sv} & \ok{}
Complete \\
Hazard unit & \src{rtl/cpu/core/rv32i_hazard_unit.sv} & \ok{}
Complete \\
Forwarding unit & \src{rtl/cpu/core/rv32i_forwarding_unit.sv} & \ok{}
Complete \\
CSR file & \src{rtl/cpu/core/rv32i_csr_file.sv} & \ok{} Complete \\
Interrupt controller & \src{rtl/cpu/core/rv32i_interrupt_ctrl.sv} &
\ok{} Complete \\
AXI arbiter & \src{rtl/cpu/rv32i_axi_arbiter.sv} & \ok{} Complete \\
Decode (updated for CSR) & \src{rtl/cpu/core/rv32i_decode.sv} & \ok{}
Complete \\
\end{longtable}
}

\textbf{Reused Phase 1 Modules} (4/4):

\begin{itemize}
\tightlist
\item
  ALU (\code{rv32i_alu.sv})
\item
  Register file (\code{rv32i_regfile.sv})
\item
  Immediate generator (\code{rv32i_imm_gen.sv})
\item
  Branch comparator (\code{rv32i_branch_comp.sv})
\end{itemize}

\textbf{Architecture Decisions Implemented (2026-02-14)}:

\begin{itemize}
\tightlist
\item
  OQ-1: \ok{} Modified in-place (Phase 1 archived to \src{micro_p/})
\item
  OQ-2: \ok{} External interrupt (MEIP) \textgreater{} Timer interrupt
  (MTIP)
\item
  OQ-3: \ok{} EBREAK sets \code{mcause=3} + \code{mepc} + triggers debug
  halt
\item
  OQ-4: \ok{} CSR write takes priority over same-cycle interrupt check
  in EX
\item
  OQ-5: \ok{} Debug halt drains pipeline immediately (no wait for MRET)
\item
  OQ-6: \ok{} In-flight AXI transaction completes; flushed responses
  discarded
\item
  OQ-7: \ok{} Add pipeline stage first → ASAP7 second → relax frequency
  last resort
\end{itemize}

\textbf{Verification Progress} (2026-02-27):

\begin{itemize}
\tightlist
\item
  \ok{} RTL implementation complete (2026-02-16)
\item
  \ok{} Initial testing complete (115/115 regression tests passing)
\item
  \ok{} Comprehensive verification complete (111/111 tests, all 7 suites
  pass):

  \begin{itemize}
  \tightlist
  \item
    smoke\_\allowbreak{}uvm: 4/4 PASS
  \item
    isa\_uvm: 54/54 PASS (37/37 RV32I instructions + CSR)
  \item
    pipeline\_\allowbreak{}hazards: 16/16 PASS (RAW/\allowbreak{}control
    hazards, forwarding, store-store, JAL rd)
  \item
    interrupts: 12/12 PASS (timer/\allowbreak{}ext IRQ, MIE gating,
    MRET, CSR insns, latency ≤3 cyc)
  \item
    debug: 6/6 PASS (single-step, BP0/BP1, GPR write, PC write, reg
    reads)
  \item
    axi\_\allowbreak{}protocol: 12/12 PASS (back-pressure, error
    injection, arbiter)
  \item
    fault\_\allowbreak{}injection: 7/7 PASS (misaligned, illegal, AXI
    fetch error)
  \end{itemize}
\item
  \ok{} Random regression: 500 seeds × 100 instructions = \textbf{50,000
  instructions, 0 failures}
\item
  \ok{} Backend flow: \textbf{75 MHz achieved on Sky130 130nm} (200 MHz
  target not met due to PDK limitations)

  \begin{itemize}
  \tightlist
  \item
    SDC constraints: \src{pnr/constraints/phase2_cpu.sdc}
  \item
    UPF power intent: \src{pnr/constraints/phase2_cpu.upf}
  \end{itemize}
\item
  \ok{} ASAP7 backend flow: \textbf{1418 MHz achieved at Run 43
  (2026-05-20)} --- 27.27 mW, 3 844 µm² stdcell, 0
  DRC/\allowbreak{}antenna/\allowbreak{}timing violations \warn{}
  unvalidated, STA-zero-wire (bead \bead{8f3}/\code{0p6}); artifacts
  wiped, not re-checkable --- \cref{cr:run:asap7}

  \begin{itemize}
  \tightlist
  \item
    Run directory: \code{pnr/asap7/runs/RUN_2026-05-20_06-27-10/}
  \item
    Config: \code{pnr/asap7/config.json} (CLOCK\_\allowbreak{}PERIOD
    0.705, CTS clustering 8/10)
  \item
    Constraints: \code{pnr/asap7/constraints/asap7.sdc}
  \item
    Full per-run history: \cref{cr:run:asap7}
  \end{itemize}
\end{itemize}

\end{record}

\subsubsection{\texorpdfstring{Phase 3: Memory System \& Caches \ok{}
COMPLETE}{Phase 3: Memory System \& Caches  COMPLETE}}\label{cr:status:phase-status:phase-3-memory-system-caches-white-check-mark-complete}

\begin{record}

\textbf{Status}: \ok{} COMPLETE (2026-05-21) --- all 20 Phase 3 tests
passing, 139/139 total tests clean

\textbf{Prerequisites}: \ok{} Phase 2 exit criteria met (2026-03-08)

\textbf{Architecture Spec}:
\src{docs/design/PHASE3_ARCHITECTURE_SPEC.md} --- APPROVED (2026-03-08),
all 6 open questions resolved

\textbf{RTL Modules} (10/10 complete):

{
\begin{longtable}[]{@{}
  >{\raggedright\arraybackslash}p{\dimexpr(\linewidth-4\tabcolsep)*330/1000\relax}
  >{\raggedright\arraybackslash}p{\dimexpr(\linewidth-4\tabcolsep)*543/1000\relax}
  >{\raggedright\arraybackslash}p{\dimexpr(\linewidth-4\tabcolsep)*107/1000\relax}@{}}
\toprule\noalign{}
\begin{minipage}[b]{\linewidth}\raggedright
Module
\end{minipage} & \begin{minipage}[b]{\linewidth}\raggedright
File
\end{minipage} & \begin{minipage}[b]{\linewidth}\raggedright
Status
\end{minipage} \\
\midrule\noalign{}
\endhead
\bottomrule\noalign{}
\endlastfoot
Cache package & \src{rtl/mem/rv32i_cache_pkg.sv} & \ok{} Complete \\
I-Cache & \src{rtl/mem/rv32i_icache.sv} & \ok{} Complete \\
D-Cache & \src{rtl/mem/rv32i_dcache.sv} & \ok{} Complete \\
Cache arbiter & \src{rtl/mem/rv32i_cache_arbiter.sv} & \ok{} Complete \\
IF stage (cache IF) & \src{rtl/cpu/core/pipeline/rv32i_pipeline_if.sv} &
\ok{} Complete \\
MEM stage (cache IF + FENCE.I) &
\src{rtl/cpu/core/pipeline/rv32i_pipeline_mem.sv} & \ok{} Complete \\
Hazard unit (renamed stalls) & \src{rtl/cpu/core/rv32i_hazard_unit.sv} &
\ok{} Complete \\
Core (cache integration) & \src{rtl/cpu/core/rv32i_core.sv} & \ok{}
Complete \\
Pipeline package (fence\_i field) &
\src{rtl/cpu/core/rv32i_pipeline_pkg.sv} & \ok{} Complete \\
Decoder (FENCE.I) & \src{rtl/cpu/core/rv32i_decode.sv} & \ok{}
Complete \\
\end{longtable}
}

\textbf{Architecture Decisions Implemented (2026-03-08)}:

\begin{itemize}
\tightlist
\item
  OQ-1: \ok{} Direct-mapped (1-way) associativity
\item
  OQ-2: \ok{} 16-byte cache line (4 words, 4 AXI transactions per
  refill)
\item
  OQ-3: \ok{} Write-back + write-allocate for D-cache
\item
  OQ-4: \ok{} No AXI burst --- 4 separate AXI4-Lite transactions per
  refill
\item
  OQ-5: \ok{} 75 MHz target on Sky130 (matches Phase 2 achieved
  frequency)
\item
  OQ-6: \ok{} Blocking cache (stall pipeline on every miss)
\end{itemize}

\textbf{Bug Fixes Landed (2026-05-21)}:

\begin{itemize}
\tightlist
\item
  \ok{} PDN \code{pdn_asap7.tcl} --- GND nets removed from
  \code{-secondary_power} (API bug \#35)
\item
  \ok{} I-cache AXI cancel race --- \code{cancel_ar_q} +
  \code{cancel_wait_r_q} flags added (bug \#36): three timing races
  fixed: R-same-cycle-as-cancel, AR-accepted-same-cycle-as-cancel, AXI
  A3.2.1 arvalid-no-retract compliance
\item
  \ok{} D-cache CS\_REFILL --- audit comment added confirming no
  equivalent cancel race
\end{itemize}

\textbf{Verification Results} (2026-05-21):

{
\begin{longtable}[]{@{}lll@{}}
\toprule\noalign{}
Suite & Tests & Result \\
\midrule\noalign{}
\endhead
\bottomrule\noalign{}
\endlastfoot
I-Cache unit tests (\texttt{make\ icache}) & 7/7 & \ok{} PASS \\
D-Cache unit tests (\texttt{make\ dcache}) & 8/8 & \ok{} PASS \\
Cache integration (\texttt{make\ cache\_\allowbreak{}integration}) & 5/5
& \ok{} PASS \\
Phase 2 full regression (\texttt{make\ test}) & 119/119 & \ok{} PASS \\
\textbf{Total} & \textbf{139/139} & \textbf{\ok{} ALL PASS} \\
\end{longtable}
}

\textbf{Achieved frequency}: 75 MHz on Sky130 130nm (ASAP7 runs 1--43
logged in \cref{cr:run:asap7-cpu})

\end{record}

\subsubsection{\texorpdfstring{Phase 4: GPU-Lite Compute Engine \ok{}
COMPLETE}{Phase 4: GPU-Lite Compute Engine  COMPLETE}}\label{cr:status:phase-status:phase-4-gpu-lite-compute-engine-white-check-mark-complete}

\begin{record}

\textbf{Status}: \ok{} COMPLETE (2026-05-27) --- all GPU tests green,
1,000-kernel random regression pass, ASAP7 PD sign-off

\begin{quote}
\textbf{Note 2026-10-10 (bead \bead{47lf}, \ghref{254})}: the Phase 4 "all
GPU tests green" sign-off did not exercise \code{BLOCK_X} above 56 ---
no multi-warp test compared RTL against a model, and \code{gpu_top}
silently ran only 7 of 8 warps for \code{BLOCK_X} 57..64 (warp 7 never
executed). Fixed in \src{rtl/gpu} (\code{WARP_CNT_W}); 64-thread
operation is verified from that date (\code{kernel_block_warps.py}, all
64 lanes vs \code{GpuRefModel}). The ASAP7 \code{gpu_top} macro views
were hardened from the pre-fix RTL and are not re-hardened (see
\cref{cr:run:asap7-gpu}).

\textbf{Note 2026-10-11 (bead \bead{6hy6})}: the three \code{BLOCK_X}
cases left by \code{47lf} are now defined and fixed: \code{BLOCK_X} 0 or
\textgreater{} 64 raises GPU ERROR at launch (reusing the \code{q6w0}
ERROR path), the block limit is 64 threads (the spec text saying 256 was
wrong), and a partial last warp starts with only its valid lanes active.
See \cref{cr:fix:fixes-index:17-gpu-block-x-launch-contract-zero-hang-64-truncation-partial-warp-lanes-bead-6hy6}.
\end{quote}

\textbf{Prerequisites}: \ok{} Phase 3 exit criteria met (2026-05-21)

\textbf{RTL Modules} (9/9 complete):

{
\begin{longtable}[]{@{}
  >{\raggedright\arraybackslash}p{\dimexpr(\linewidth-4\tabcolsep)*291/1000\relax}
  >{\raggedright\arraybackslash}p{\dimexpr(\linewidth-4\tabcolsep)*543/1000\relax}
  >{\raggedright\arraybackslash}p{\dimexpr(\linewidth-4\tabcolsep)*146/1000\relax}@{}}
\toprule\noalign{}
\begin{minipage}[b]{\linewidth}\raggedright
Module
\end{minipage} & \begin{minipage}[b]{\linewidth}\raggedright
File
\end{minipage} & \begin{minipage}[b]{\linewidth}\raggedright
Status
\end{minipage} \\
\midrule\noalign{}
\endhead
\bottomrule\noalign{}
\endlastfoot
GPU top & \src{rtl/gpu/gpu_top.sv} & \ok{} Complete \\
Command queue & \src{rtl/gpu/gpu_command_queue.sv} & \ok{} Complete \\
Warp scheduler & \src{rtl/gpu/warp_scheduler.sv} & \ok{} Complete \\
Compute unit & \src{rtl/gpu/gpu_compute_unit.sv} & \ok{} Complete \\
Vector register file & \src{rtl/gpu/vector_register_file.sv} & \ok{}
Complete \\
Vector ALU & \src{rtl/gpu/vector_alu.sv} & \ok{} Complete \\
Memory unit & \src{rtl/gpu/gpu_memory_unit.sv} & \ok{} Complete \\
Memory coalescer & \src{rtl/gpu/memory_coalescer.sv} & \ok{} Complete \\
Shared memory & \src{rtl/gpu/shared_memory.sv} & \ok{} Complete \\
\end{longtable}
}

\textbf{Verification Results} (2026-05-23):

{
\begin{longtable}[]{@{}
  >{\raggedright\arraybackslash}p{\dimexpr(\linewidth-4\tabcolsep)*630/1000\relax}
  >{\raggedright\arraybackslash}p{\dimexpr(\linewidth-4\tabcolsep)*257/1000\relax}
  >{\raggedright\arraybackslash}p{\dimexpr(\linewidth-4\tabcolsep)*93/1000\relax}@{}}
\toprule\noalign{}
\begin{minipage}[b]{\linewidth}\raggedright
Suite
\end{minipage} & \begin{minipage}[b]{\linewidth}\raggedright
Tests
\end{minipage} & \begin{minipage}[b]{\linewidth}\raggedright
Result
\end{minipage} \\
\midrule\noalign{}
\endhead
\bottomrule\noalign{}
\endlastfoot
GPU unit tests (\texttt{make\ gpu\_\allowbreak{}unit}) & all & \ok{}
PASS \\
GPU kernel tests (\texttt{make\ gpu\_\allowbreak{}kernels}) & all &
\ok{} PASS \\
CPU-GPU handoff (\code{test_cpu_gpu_handoff.py}) & 1/1 & \ok{} PASS \\
CPU re-gate (\texttt{make\ test} +
\texttt{make\ random\_\allowbreak{}uvm}) & 140/140 + 100k instr & \ok{}
PASS \\
GPU random regression (\texttt{make\ gpu\_\allowbreak{}random}) & 1,000
kernels & \ok{} PASS \\
\end{longtable}
}

\textbf{Physical Design} (2026-05-28):

\begin{itemize}
\tightlist
\item
  \ok{} ASAP7 sign-off: \textbf{571 MHz (1.75 ns) / 262 mW / 115,600 µm²
  die / 60,500 µm² stdcell / 70\% util} \warn{} unvalidated,
  STA-zero-wire (bead \bead{8f3}/\code{0p6}); GPU run artifacts wiped,
  re-validation deferred (host cost estimate: several hours / 9--14+
  GiB) --- \cref{cr:run:asap7-gpu}
\item
  \ok{} Setup WS +197.3 ps (0 violations), Hold WS +16.3 ps (0
  violations), slew/\allowbreak{}cap/\allowbreak{}fanout 0, antenna 0
\item
  \ok{} Run: \code{pnr/asap7/gpu/runs/RUN_2026-05-28_06-29-48/}
  (supersedes 500 MHz \code{RUN_2026-05-27_11-16-37})
\item
  \ok{} Constraints: \src{pnr/asap7/gpu/constraints/asap7_gpu.sdc}; full
  history: \cref{cr:run:asap7-gpu}
\end{itemize}

\textbf{Phase 4 sign-off frequency is 571 MHz} --- the
\code{CLOCK_PERIOD} 2.0→1.75 ns stretch push closed clean
(\code{RUN_2026-05-28_06-29-48}), upgrading the GPU signoff from the
prior 500 MHz \code{RUN_2026-05-27_11-16-37} (now superseded). Hold is
clean at +16.3 ps / 0 viol at the final stage; the prior 500 MHz
run\textquotesingle s \code{final/metrics.json} showed hold −212 ps /
368 viol.

\textbf{Signoff caveats} (deferred to Phase-5 SoC PD, same as the prior
run): PDN connectivity not closed (\code{PSM-0069} / \code{PDN-0179},
9.56 M grid viol --- identical to the 500 MHz run); 325 \code{DRT-0074}
on top-level I/O ports only (0 internal-net DRC); timing is post-GRT
estimated (\code{STAPostPNR} + \code{RCX} gated off --- same methodology
as prior run). A confirmation re-run with post-PnR STA + RCX enabled is
recommended before locking the number.

\textbf{Step-37 runtime note}: \code{repair_design_postgrt} is the
bottleneck (\textasciitilde18.9 h single-threaded --- runs GRT twice + a
multi-thousand-iteration repair loop on \textasciitilde485 K instances).
Mitigation: \code{DRT_THREADS} 4→12 (\src{pnr/asap7/gpu/config.json}).

\textbf{Macro views}: signoff DB exported for Phase-5 SoC via
\texttt{make\ macro-\allowbreak{}views-\allowbreak{}asap7\ BLOCK=\allowbreak{}gpu}
→
\texttt{pnr/\allowbreak{}asap7/\allowbreak{}gpu/\allowbreak{}macro/\allowbreak{}\{gpu\_\allowbreak{}top.\allowbreak{}lef,\allowbreak{}\ *.\allowbreak{}lib,\allowbreak{}\ gpu\_\allowbreak{}top.\allowbreak{}nl.\allowbreak{}v.\allowbreak{}gz\}}
(netlist gzipped to clear GitHub\textquotesingle s 100 MB limit;
\texttt{gunzip\ -\allowbreak{}k} to restore).

\end{record}

\subsubsection{\texorpdfstring{Phase 5: SoC Integration
\ok{}}{Phase 5: SoC Integration }}\label{cr:status:phase-status:phase-5-soc-integration-white-check-mark}

\begin{record}

\textbf{Status}: COMPLETE (2026-06-24) --- M1--M12 done. M11 ASAP7 SoC
P\&R signed off 571 MHz / 62.9 mW / 520×520 µm / 65.6 \% util / 0 DRC /
0 antenna \warn{} timing/\allowbreak{}power unvalidated
post-bead-\code{8f3} (see top of this file)
(\cref{cr:run:phase5}). Milestone detail in
\src{docs/PHASE5_SOC_INTEGRATION_PLAN.md} (golden spec)

{\small\setlength{\tabcolsep}{4pt}

{
\begin{longtable}[]{@{}
  >{\raggedright\arraybackslash}p{\dimexpr(\linewidth-4\tabcolsep)*80/1000\relax}
  >{\raggedright\arraybackslash}p{\dimexpr(\linewidth-4\tabcolsep)*450/1000\relax}
  >{\raggedright\arraybackslash}p{\dimexpr(\linewidth-4\tabcolsep)*450/1000\relax}@{}}
\toprule\noalign{}
\begin{minipage}[b]{\linewidth}\raggedright
Milestone
\end{minipage} & \begin{minipage}[b]{\linewidth}\raggedright
Scope
\end{minipage} & \begin{minipage}[b]{\linewidth}\raggedright
Status
\end{minipage} \\
\midrule\noalign{}
\endhead
\bottomrule\noalign{}
\endlastfoot
M1 & AXI4/\allowbreak{}AXI-Lite shared package & \ok{} 2026-05-31
(lint-clean) \\
M2 & Cache refill FSMs → AXI4 burst & \ok{} 2026-05-31 (146/146
tests) \\
M3 & AXI4 crossbar + AXI-Lite interconnect & \ok{} 2026-05-31 (16/16
tests) \\
M4 & UART / SPI / timer / IRQ controller & \ok{} 2026-06-01 \\
M5 & DMA engine & \ok{} 2026-06-01 (6/6 tests) \\
M6 & Behavioral SRAM controller & \ok{} 2026-06-01 (7/7 tests) \\
M7 & Performance counters (CSR + AXI-Lite GPU stats) & \ok{} 2026-06-02
(CPU re-sign-off 1282 MHz) \\
M8 & SoC top integration (\src{rtl/soc/soc_top.sv}) & \ok{} 2026-06-02
(lint-clean) \\
M9 & SoC verification --- foundation slice (boot 100/100) & \ok{}
2026-06-03 (\prref{65}) \\
M9 & Fast-follows: CPU→GPU kernel launch, SW coherency,
DMA/\allowbreak{}peripheral loopback, SRAM readback, 1M-cycle random,
benchmarks & \ok{} 2026-06-24 (1,079,867 cyc, 0 fail) \\
M10 & L2 cache decision gate (needs M9 benchmarks) & \ok{} 2026-06-21
--- \textbf{NO-GO} (\src{docs/M10_L2_DECISION_ANALYSIS.md}) \\
M11 & SoC P\&R + STA (\code{phase5_soc.sdc}/\code{.upf}) & \ok{}
2026-06-24 \warn{} figures unvalidated, bead \bead{0p6} \\
M12 & Sign-off + documentation & \ok{} 2026-06-24 \\
\end{longtable}
}

}

\end{record}

\subsubsection{\texorpdfstring{Phase 6: IP Expansion \ok{} COMPLETE
(2026-10-05)}{Phase 6: IP Expansion  COMPLETE (2026-10-05)}}\label{cr:status:phase-status:phase-6-ip-expansion-white-check-mark-complete-2026-10-05}

\begin{record}

\textbf{Golden spec}: \src{docs/PHASE6_IP_EXPANSION_PLAN.md} (added
2026-09-28, bead \bead{f7vs.1}). \textbf{Tracking}: bead epic
\code{claude_verilog_test-f7vs}. \textbf{Definition of done, per item}:
RTL + cocotb verification + SoC integration. Physical design is a
separately batched gate (Gate A synthesis probe per item; Gate B one
Sky130 harden), not per-item.

{\footnotesize\setlength{\tabcolsep}{3pt}

{
\begin{longtable}[]{@{}
  >{\raggedright\arraybackslash}p{\dimexpr(\linewidth-8\tabcolsep)*49/1000\relax}
  >{\raggedright\arraybackslash}p{\dimexpr(\linewidth-8\tabcolsep)*363/1000\relax}
  >{\raggedright\arraybackslash}p{\dimexpr(\linewidth-8\tabcolsep)*115/1000\relax}
  >{\raggedright\arraybackslash}p{\dimexpr(\linewidth-8\tabcolsep)*79/1000\relax}
  >{\raggedright\arraybackslash}p{\dimexpr(\linewidth-8\tabcolsep)*375/1000\relax}@{}}
\toprule\noalign{}
\begin{minipage}[b]{\linewidth}\raggedright
Item
\end{minipage} & \begin{minipage}[b]{\linewidth}\raggedright
Scope
\end{minipage} & \begin{minipage}[b]{\linewidth}\raggedright
APB idx / IRQ bit
\end{minipage} & \begin{minipage}[b]{\linewidth}\raggedright
Bead
\end{minipage} & \begin{minipage}[b]{\linewidth}\raggedright
Status
\end{minipage} \\
\midrule\noalign{}
\endhead
\bottomrule\noalign{}
\endlastfoot
6a-1 & GPIO controller, 32 pins & 7 / 5 & \code{ckc}, \code{00ef} &
\ok{} 2026-09-26 (Sky130-hardened; LVS PASS, DRC 0, setup+hold clean
9/9) \\
G1 & Golden spec + ROADMAP bus/\allowbreak{}numbering corrections & ---
& \code{f7vs.1} & \ok{} 2026-09-28 \\
G2 & APB window + IRQ source pre-allocation & --- & \code{f7vs.2} &
\ok{} 2026-09-28 (\prref{197}) \\
G3 & \code{tb_soc_top.sv} ports + SDC exceptions, one forced clean & ---
& \code{f7vs.3} & \ok{} 2026-09-28 (\prref{197}) \\
G4 & File-list checker + first \code{apb_interconnect} unit suite & ---
& \code{f7vs.4} & \ok{} 2026-09-28 (\prref{197}) \\
G5 & This section + de-staled Next Actions & --- & \code{f7vs.5} & \ok{}
2026-09-28 (\prref{197}) \\
6a-2 & PWM controller, 4 channels & 8 / 6 & \code{f7vs.6} & \ok{}
2026-09-28 (RTL + 16/16 L1 + soc\_pwm L2; \code{soc_all_ci} 279) \\
6a-3 & Watchdog timer & 9 / 7 & \code{f7vs.7} & \ok{} 2026-09-29 (RTL +
22/22 L1 + soc\_wdt L2 incl. RST\_EN reset path; \code{soc_all_ci}
303) \\
6a-4 & TRNG (portable LFSR entropy; RO source Sky130-only) & 10 / 8 &
\code{f7vs.8} & \ok{} 2026-10-01 (RTL bit-exact vs Python model + 22/22
L1 + soc\_trng L2; \code{soc_all_ci} 326) \\
6a-5 & I2C master controller & 11 / 9 & \code{f7vs.9} & \ok{} 2026-10-03
(RTL + bit engine + \code{i2c_slave} BFM + 46/46 L1 + soc\_i2c L2; \prref{208}; \code{soc_all_ci} 394) \\
6b & CRYPTO --- AES-128 (ECB+CTR) + SHA-256 & 12 / 10 & \code{f7vs.10} &
\ok{} 2026-10-03 (RTL 3 files + 41/41 L1 + soc\_\allowbreak{}crypto L2 +
\textbf{Gate A run}: 277 875 µm² ≈ 28 \% of SoC stdcell area --- first
recorded as 3.59 \% against a stdcell+macro denominator, corrected
2026-10-04 --- +14.04 ns ss @ 25 ns) \\
6c & INT8 NPU, 4×4 weight-stationary, 4 KB weight SRAM & 13 / 11 &
\code{f7vs.11} & \ok{} 2026-10-04 (RTL 3 files + 39/39 L1 + soc\_npu L2
+ mutation 65 mutants: 55 killed, 10 provably equivalent, 0 gaps +
\textbf{Gate A}: 203 179 µm² stdcell ≈ 20 \% of SoC stdcell + 4 KB
macro, +1.84 ns ss @ 25 ns) \\
Gate B & Batched Sky130 SoC harden with all Phase 6 IP & --- &
\code{f7vs.15} & \ok{} 2026-10-05 (\code{RUN_2026-10-05_06-50-11}, \prref{214}): setup met at all 9 corners (worst max\_ss +0.579 ns), hold 0
violators, Netgen LVS clear, routing DRC 0, 81.0 mW total, 44.8 \% util;
antenna 210 nets / 305 pins accepted by user decision; \code{e45j}
slew/cap recorded, not a Phase 6 exit criterion. Source:
\src{docs/PHASE6_IP_EXPANSION_PLAN.md} §8, run 2 \\
\end{longtable}
}

}

Three documentation defects were corrected when the golden spec landed:
AES/SHA is \textbf{APB4}, not AXI4-Lite (\code{ROADMAP.md:431} (\cref{cr:status:roadmap}) was a
survivor of the pre-2026-09-20 text); the AES throughput figures
describe a \textbf{deferred AXI4-master version}, not what 6b lands; and
the TRNG \textbf{RTL is portable} --- only its ring-oscillator entropy
source is Sky130-exclusive.

\warn{} \textbf{ASAP7 is out of scope for all of Phase 6},
triple-blocked on beads \code{ma7} (proven Synlig \code{OPT_MUXTREE}
miscompile in the CPU/GPU macro netlists), \code{lxv}
(routing-congestion regression) and \code{2kn} (needs a ≥32 GB host;
this one has \textasciitilde15 GB). Every peripheral still goes into the
ASAP7 sv2v file list, enforced by
\src{tools/verif/check_periph_filelists.py}, so unblocking is a re-run
rather than a porting exercise.

\end{record}

\subsubsection{\texorpdfstring{Phase 7: Mixed-Signal PLL Clock Generator
\wip{}}{Phase 7: Mixed-Signal PLL Clock Generator }}\label{cr:status:phase-status:phase-7-mixed-signal-pll-clock-generator-construction}

\begin{record}

\textbf{Status}:
M-a/\allowbreak{}M-b1/\allowbreak{}M-b2/\allowbreak{}M-c \ok{} complete
(2026-06-20); M-d = documentation (this update). Golden spec:
\src{docs/PHASE7_MIXED_SIGNAL_PLL_PLAN.md}.

Dual-PDK charge-pump integer-N PLL (ring VCO) designed end-to-end via
the analog-design orchestrator agents, integrated into the SoC as the
clock source via a PDK-agnostic real-number model (RNM). Doubled as a
validation run of the analog-design agent suite (real
ngspice/\allowbreak{}magic/\allowbreak{}netgen ran; closed-loop meta fix
loop + adversarial physical verification both added value --- the latter
caught a false DRC pass).

{\footnotesize\setlength{\tabcolsep}{3pt}

{
\begin{longtable}[]{@{}
  >{\raggedright\arraybackslash}p{\dimexpr(\linewidth-4\tabcolsep)*83/1000\relax}
  >{\raggedright\arraybackslash}p{\dimexpr(\linewidth-4\tabcolsep)*467/1000\relax}
  >{\raggedright\arraybackslash}p{\dimexpr(\linewidth-4\tabcolsep)*430/1000\relax}@{}}
\toprule\noalign{}
\begin{minipage}[b]{\linewidth}\raggedright
Milestone
\end{minipage} & \begin{minipage}[b]{\linewidth}\raggedright
Scope
\end{minipage} & \begin{minipage}[b]{\linewidth}\raggedright
Status
\end{minipage} \\
\midrule\noalign{}
\endhead
\bottomrule\noalign{}
\endlastfoot
M-a & Analog infra (\code{analog} devshell) + dual-PDK PLL architecture
+ modeling & \ok{} \\
M-b1 & ASAP7 variant: circuit + ngspice --- 100 MHz→1.282 GHz, N=13, 0.7
V & \ok{} (electrical \textbf{indicative}: BSIM4 substitute, no
BSIM-CMG) \\
M-b2 & Sky130 variant: real \code{sky130_fd_pr} circuit/\allowbreak{}sim
+ magic layout --- 10→100 MHz, N=10, 1.8 V; \textbf{CP-block DRC=0 +
netgen LVS MATCH} & \ok{} (CP block;
VCO/\allowbreak{}LF/\allowbreak{}BIAS regen \hold{}) \\
M-c & SoC integration (\src{rtl/soc/pll/}): clock seam, AXI-Lite slave
@0x2000\_\allowbreak{}7000 (\code{AXIL_N_SLAVES} 6→7), 2 RTL bugs fixed,
cosim 7/7 + regression & \ok{} \\
M-d & Phase 7 documentation
(plan/\allowbreak{}roadmap/\allowbreak{}status/\allowbreak{}CLAUDE.md) &
\ok{} 2026-06-20 \\
\end{longtable}
}

}

\textbf{Real vs indicative}: Sky130 CP block carries \textbf{real}
DRC+LVS sign-off; ASAP7 electrical is \textbf{indicative} (no BSIM-CMG
FinFET models in open ngspice). \textbf{Follow-ups}: Sky130 full-chip
DRC/LVS (mechanical VCO/\allowbreak{}LF/\allowbreak{}BIAS PDK-generator
regen) and the RNM-mode AXI CDC synchroniser. Two RTL bugs found+fixed:
PLL bootstrap deadlock (regs were on the gated core domain → moved to
ref \code{clk_i}, default-enable \code{CONTROL[0]}) and
\code{PERIPH_LIMIT} decode hole (0x2000\_\allowbreak{}6FFF →
0x2000\_\allowbreak{}7FFF).

\end{record}

\subsection{Recent Project Changes}\label{cr:status:phase-status:recent-project-changes}

\subsubsection{2026-06-20: Phase 7 Mixed-Signal PLL --- M-a..M-c
complete, M-d docs}\label{cr:status:phase-status:2026-06-20-phase-7-mixed-signal-pll-m-am-c-complete-m-d-docs}

\begin{record}

\begin{itemize}
\tightlist
\item
  \ok{} Dual-PDK charge-pump integer-N PLL (ring VCO) designed
  end-to-end via the analog-design orchestrator agents.
\item
  \ok{} ASAP7 variant (indicative): 100 MHz→1.282 GHz, N=13, 0.7 V;
  electrical via calibrated planar BSIM4 substitute (no BSIM-CMG in open
  ngspice).
\item
  \ok{} Sky130 variant (real): 10→100 MHz, N=10, 1.8 V, real
  \code{sky130_fd_pr}; Kvco 530 MHz/V; CP UP/DN mismatch 4.78\%→0.37\%
  (independent bias legs); behavioral lock 100.000 MHz / 1.38 µs; real
  magic layout + \textbf{CP-block DRC=0 + netgen LVS MATCH} (4 pfet + 1
  nfet extracted). Adversarial PV caught a false DRC pass (hand-drawn
  geometry 50× too small → 0 FETs).
\item
  \ok{} SoC integration (\src{rtl/soc/pll/}): clock seam
  (\code{clk_i}→PLL ref→\code{core_clk},
  \texttt{core\_\allowbreak{}rst\_\allowbreak{}n\ =\allowbreak{}\ rst\_\allowbreak{}n\_\allowbreak{}i\ \&\ pll\_\allowbreak{}locked}),
  AXI-Lite slave @0x2000\_\allowbreak{}7000 (\code{AXIL_N_SLAVES} 6→7).
  Cosim 7/7 (\code{test_pll_lock} 3/3 + \code{test_pll_regs} 4/4) +
  boot/\allowbreak{}periph regression clean.
\item
  \ok{} Two RTL bugs found+fixed: PLL bootstrap deadlock (regs on gated
  core domain → ref \code{clk_i} + default-enable \code{CONTROL[0]});
  \code{PERIPH_LIMIT}
  0x2000\_\allowbreak{}6FFF→0x2000\_\allowbreak{}7FFF.
\item
  \hold{} Follow-ups: Sky130 full-chip DRC/LVS
  (VCO/\allowbreak{}LF/\allowbreak{}BIAS PDK-generator regen) and
  RNM-mode AXI CDC synchroniser.
\end{itemize}

\end{record}

\subsubsection{\texorpdfstring{2026-05-21: Phase 3 COMPLETE
\ok{}}{2026-05-21: Phase 3 COMPLETE }}\label{cr:status:phase-status:2026-05-21-phase-3-complete-white-check-mark}

\begin{record}

\textbf{All exit criteria met --- 139/139 tests passing}:

\begin{itemize}
\tightlist
\item
  \ok{} 10/10 Phase 3 RTL modules implemented (cache\_\allowbreak{}pkg,
  icache, dcache, cache\_\allowbreak{}arbiter, 6 modified pipeline
  files)
\item
  \ok{} Bug fixes: PDN secondary\_\allowbreak{}power API (\#35), I-cache
  AXI cancel race with 3 sub-cases (\#36)
\item
  \ok{} I-cache unit tests: 7/7 PASS (\texttt{make\ icache}) --- hit,
  miss, conflict, FENCE.I, latency, boundary
\item
  \ok{} D-cache unit tests: 8/8 PASS (\texttt{make\ dcache}) ---
  read/\allowbreak{}write hit/miss, dirty eviction, strobes, latency
\item
  \ok{} Cache integration tests: 5/5 PASS
  (\texttt{make\ cache\_\allowbreak{}integration}) --- locality,
  load-after-store, FENCE.I self-modifying code, conflict stress, warmup
  IPC
\item
  \ok{} Phase 2 full regression: 119/119 PASS (no regressions from Phase
  3 RTL changes)
\item
  \ok{} Makefile fix: \code{PYTHON3=/usr/bin/python3} (system Python
  3.10 matches \code{PYTHONHOME=/usr}; nix Python 3.11 was mismatched)
\item
  \ok{} Test timing fix: 4 tests updated to use
  \code{_fetch()}/\code{_read()}/\code{_write()} helpers (5-state SRAM
  pipeline takes 3 cycles for a hit from CS\_DONE, not 1)
\item
  \textbf{TOTAL: 139/139 tests (119 Phase 2 + 20 Phase 3), 0 failures}
\item
  \textbf{Branch}: \code{bug-fix-35-36} → ready to merge to \code{main}
\end{itemize}

\end{record}

\subsubsection{2026-03-08: Phase 3 RTL Implementation Started}\label{cr:status:phase-status:2026-03-08-phase-3-rtl-implementation-started}

\begin{record}

\textbf{Phase 3 kicked off}:

\begin{itemize}
\tightlist
\item
  \ok{} PHASE3\_\allowbreak{}ARCHITECTURE\_\allowbreak{}SPEC.md approved
  --- all 6 open questions resolved
\item
  \ok{} Phase 2 marked complete --- 75 MHz achieved on Sky130 130nm
\item
  \wip{} Phase 3 RTL implementation in progress (cache package, I-cache,
  D-cache, arbiter)
\item
  \wip{} Python cache reference model development started
  (\src{tb/models/cache_model.py})
\end{itemize}

\textbf{Architecture highlights}:

\begin{itemize}
\tightlist
\item
  I-Cache: 4 KB direct-mapped, 16-byte lines, 256 sets; read-only;
  FENCE.I invalidation
\item
  D-Cache: 4 KB direct-mapped, 16-byte lines, 256 sets; write-back +
  write-allocate
\item
  External interface: 4 sequential AXI4-Lite transactions per refill (no
  burst)
\item
  Cache arbiter: D-cache priority over I-cache (replaces Phase 2 AXI
  arbiter)
\item
  Target frequency: 75 MHz on Sky130 130nm
\end{itemize}

\end{record}

\subsubsection{2026-02-27: Phase 2 Comprehensive Verification COMPLETE}\label{cr:status:phase-status:2026-02-27-phase-2-comprehensive-verification-complete}

\begin{record}

\textbf{All 7 test suites passing, 50,000 random instructions verified}:

\begin{itemize}
\tightlist
\item
  \ok{} smoke\_\allowbreak{}uvm: 4/4 PASS
\item
  \ok{} isa\_uvm: 54/54 PASS --- full RV32I ISA + Zicsr (CSR)
  instructions
\item
  \ok{} pipeline\_\allowbreak{}hazards: 16/16 PASS ---
  RAW/\allowbreak{}control hazards, EX/MEM/WB forwarding
\item
  \ok{} interrupts: 12/12 PASS --- timer/ext IRQ delivery, MIE gating,
  MRET, IRQ latency ≤3 cycles
\item
  \ok{} debug: 6/6 PASS --- single-step, breakpoints, GPR/PC write,
  register reads
\item
  \ok{} axi\_\allowbreak{}protocol: 12/12 PASS --- back-pressure, error
  injection, protocol compliance
\item
  \ok{} fault\_\allowbreak{}injection: 7/7 PASS --- misaligned access,
  illegal instruction, AXI fetch error
\item
  \ok{} Random regression: 500 seeds × 100 instructions = \textbf{50,000
  instructions, 0 failures}
\item
  \textbf{TOTAL: 111/111 Phase 2 tests passing}
\end{itemize}

\textbf{Infrastructure improvements}:

\begin{itemize}
\tightlist
\item
  Created \src{tb/cocotb/cpu/phase2_test_utils.py} --- shared APBDebug +
  setup helpers
\item
  Added Phase 2 Makefile targets (\code{pipeline_hazards},
  \code{interrupts}, \code{axi_protocol}, \code{fault_injection},
  \code{phase2_all})
\item
  Added 2 new pipeline hazard tests (\code{test_back_to_back_stores},
  \code{test_jal_rd_dependency})
\item
  Fixed \code{test_csr_mstatus_mie_gate}: DBG\_\allowbreak{}MSTATUS (APB
  0x200) is READ-ONLY per RTL design
\end{itemize}

\end{record}

\subsubsection{\texorpdfstring{2026-02-16: Phase 2 RTL Implementation
COMPLETE \ok{}}{2026-02-16: Phase 2 RTL Implementation COMPLETE }}\label{cr:status:phase-status:2026-02-16-phase-2-rtl-implementation-complete-tada}

\begin{record}

\textbf{All 14 RTL modules implemented}:

\begin{itemize}
\tightlist
\item
  \ok{} 5-stage pipeline modules (IF, ID, EX, MEM, WB)
\item
  \ok{} Hazard detection and forwarding units
\item
  \ok{} CSR file and interrupt controller
\item
  \ok{} AXI arbiter for IF/MEM priority
\item
  \ok{} Phase 1 modules reused (ALU, regfile, imm\_gen,
  branch\_\allowbreak{}comp)
\item
  \ok{} Initial regression: 115/115 tests passing
\item
  \wip{} Comprehensive verification in progress
\end{itemize}

\end{record}

\subsubsection{\texorpdfstring{2026-02-14: Phase 2 Architecture Approved
\ok{}}{2026-02-14: Phase 2 Architecture Approved }}\label{cr:status:phase-status:2026-02-14-phase-2-architecture-approved-tada}

\begin{record}

\textbf{Architecture specification finalized}:

\begin{itemize}
\tightlist
\item
  \ok{} All 7 open questions resolved
\item
  \ok{} 5-stage pipeline design approved
\item
  \ok{} Interrupt support (M-mode, timer + external)
\item
  \ok{} CSR instructions
  (CSRRW/\allowbreak{}S/\allowbreak{}C/\allowbreak{}I variants)
\item
  \ok{} Hazard handling strategy defined
\item
  \ok{} Debug interface updated for pipeline drain
\item
  \ok{} RTL implementation authorized
\end{itemize}

\end{record}

\subsubsection{\texorpdfstring{2026-02-13: Phase 1 Verification COMPLETE
\ok{}}{2026-02-13: Phase 1 Verification COMPLETE }}\label{cr:status:phase-status:2026-02-13-phase-1-verification-complete-tada}

\begin{record}

\textbf{All 9/9 exit criteria met}:

\begin{itemize}
\tightlist
\item
  \ok{} Tasks 5, 6, 7 complete --- debug interface tests (6/6), coverage
  (37/37 instructions, 8/8 states), documentation
\item
  \ok{} Phase 1 RTL implementation confirmed complete (8/8 modules,
  \textasciitilde1,900 lines)
\item
  \ok{} All verification suites passing with 0 failures
\item
  \ok{} Ready to begin Phase 2 (5-stage pipelined CPU)
\end{itemize}

\end{record}

\subsubsection{\texorpdfstring{2026-02-07: Task 4 Complete --- AXI
Protocol Tests
\ok{}}{2026-02-07: Task 4 Complete --- AXI Protocol Tests }}\label{cr:status:phase-status:2026-02-07-task-4-complete-axi-protocol-tests-tada}

\begin{record}

\begin{itemize}
\tightlist
\item
  \ok{} 11/11 AXI4-Lite protocol tests implemented and passing
\item
  \ok{} Back-pressure, error injection, and protocol compliance
  categories covered
\item
  \ok{} \src{tb/cocotb/cpu/axi_models.py} (380 lines) +
  \code{test_axi_protocol.py} (860 lines)
\end{itemize}

\end{record}

\subsubsection{\texorpdfstring{2026-01-28: Task 3 Complete --- Random
Instruction Tests
\ok{}}{2026-01-28: Task 3 Complete --- Random Instruction Tests }}\label{cr:status:phase-status:2026-01-28-task-3-complete-random-instruction-tests-tada}

\begin{record}

\begin{itemize}
\tightlist
\item
  \ok{} 10,000 random instructions (100 seeds × 100 instructions), 0
  failures
\item
  \ok{} \src{tb/generators/rv32i_instr_gen.py} implemented
\item
  \ok{} \code{tb/cocotb/cpu/test_random_instructions.py} with multi-seed
  support
\end{itemize}

\end{record}

\subsubsection{\texorpdfstring{2026-01-26: Task 2 Complete --- ISA
Compliance Tests
\ok{}}{2026-01-26: Task 2 Complete --- ISA Compliance Tests }}\label{cr:status:phase-status:2026-01-26-task-2-complete-isa-compliance-tests-tada}

\begin{record}

\begin{itemize}
\tightlist
\item
  \ok{} All 37/37 RV32I instructions tested and passing
\item
  \ok{} Major RTL bugs fixed (branch/\allowbreak{}jump timing, load data
  latching, register file reads)
\end{itemize}

\end{record}

\subsubsection{\texorpdfstring{2026-01-18: Phase 0 APPROVED - Ready for
Phase 1 \ok{}}{2026-01-18: Phase 0 APPROVED - Ready for Phase 1 }}\label{cr:status:phase-status:2026-01-18-phase-0-approved-ready-for-phase-1-tada}

\begin{record}

\textbf{Phase 0 Exit Criteria Met}:

\begin{itemize}
\tightlist
\item
  \ok{} All 7 specifications reviewed and approved by human
\item
  \ok{} Python reference models validated (66/66 tests passing)
\item
  \ok{} cocotb test infrastructure reviewed and approved
\item
  \ok{} Project ready to transition to Phase 1 RTL implementation
\end{itemize}

\textbf{Authorization}: Phase 1 RTL development may now begin per
PHASE1\_\allowbreak{}ARCHITECTURE\_\allowbreak{}SPEC.md

\end{record}

\subsubsection{\texorpdfstring{2026-01-18: Phase 0 Implementation
Complete \ok{}}{2026-01-18: Phase 0 Implementation Complete }}\label{cr:status:phase-status:2026-01-18-phase-0-implementation-complete-white-check-mark}

\begin{record}

\textbf{Python Reference Models}:

\begin{itemize}
\tightlist
\item
  \ok{} \src{tb/models/memory_model.py} - Sparse memory model with
  alignment checking (157 lines, 21 tests)
\item
  \ok{} \src{tb/models/rv32i_model.py} - Instruction-accurate RV32I CPU
  model (450+ lines, 33 tests)
\item
  \ok{} \src{tb/models/gpu_kernel_model.py} - SIMT GPU execution model
  (450+ lines, 12 tests)
\item
  \ok{} All 66 unit tests passing
\end{itemize}

\textbf{cocotb Test Infrastructure}:

\begin{itemize}
\tightlist
\item
  \ok{} \src{tb/cocotb/bfm/axi4lite_master.py} - AXI4-Lite master BFM
  (200+ lines)
\item
  \ok{} \src{tb/cocotb/bfm/apb3_master.py} - APB3 master BFM with debug
  interface (250+ lines)
\item
  \ok{} \src{tb/cocotb/common/scoreboard.py} - RTL vs reference model
  comparison (130+ lines)
\item
  \ok{} \src{tb/cocotb/common/clock_reset.py} - Clock and reset
  utilities
\item
  \ok{} \src{tb/cocotb/cpu/test_example_counter.py} - Example test (3/3
  tests passing)
\item
  \ok{} \src{tb/cocotb/cpu/test_smoke.py} - CPU smoke tests (6/6 tests
  passing)
\item
  \ok{} \code{tb/cocotb/cpu/test_isa_compliance.py} - ISA compliance
  tests (33/37 passing)
\item
  \ok{} Complete documentation (README.md,
  COCOTB\_\allowbreak{}SETUP\_\allowbreak{}SUMMARY.md)
\end{itemize}

\textbf{Issues Resolved}:

\begin{itemize}
\tightlist
\item
  Fixed Makefile clean target conflicts
\item
  Updated to cocotb 2.0 API (logging changes)
\item
  Fixed test timing issues in counter disable test
\end{itemize}

\textbf{Phase 0 Status}: \ok{} COMPLETE - All specifications, reference
models, and infrastructure approved (2026-01-18)

\end{record}

\subsubsection{2026-01-17: Phase 0 Documentation Complete}\label{cr:status:phase-status:2026-01-17-phase-0-documentation-complete}

\begin{record}

\begin{itemize}
\tightlist
\item
  \ok{} All 7 specification documents finalized
\item
  \ok{} MEMORY\_\allowbreak{}MAP.md aligned to 4 KB minimum regions
\item
  \ok{} All known documentation issues resolved
\item
  Ready to proceed with Python reference model implementation
\item
  Phase 0 specifications ready for human review
\end{itemize}

\end{record}

\subsubsection{2026-01-17: Specification Alignment}\label{cr:status:phase-status:2026-01-17-specification-alignment}

\begin{record}

\begin{itemize}
\tightlist
\item
  Identified documentation gaps
\item
  Created phase status tracking
\item
  Aligning all specifications with current project state
\end{itemize}

\end{record}

\subsubsection{2026-01-16: Project Restart}\label{cr:status:phase-status:2026-01-16-project-restart}

\begin{record}

\begin{itemize}
\tightlist
\item
  Removed previous RTL implementation
\item
  Starting fresh with specification-driven approach
\item
  Commit:
  \texttt{1b4bb3b\ [Code]\ Remove\ database\ to\ restart\ project}
\end{itemize}

\end{record}

\subsection{Known Issues}\label{cr:status:phase-status:known-issues}

\subsubsection{\texorpdfstring{\ok{} Resolved
(2026-01-17)}{ Resolved (2026-01-17)}}\label{cr:status:phase-status:white-check-mark-resolved-2026-01-17}

\begin{record}

All previous specification issues have been resolved:

\begin{enumerate}
\def\labelenumi{\arabic{enumi}.}
\tightlist
\item
  \ok{} \textbf{CLAUDE.md RTL references} - All RTL references properly
  labeled as "Planned Architecture (Phase 1)"
\item
  \ok{} \textbf{RTL\_\allowbreak{}DEFINITION.md protocol details} - Now
  includes AXI4-Lite (ARM IHI 0022E) and APB3 (ARM IHI 0024C)
  specifications
\item
  \ok{} \textbf{GPU specification} -
  PHASE4\_\allowbreak{}GPU\_\allowbreak{}ARCHITECTURE\_\allowbreak{}SPEC.md
  created (520 lines)
\item
  \ok{} \textbf{Reference model specification} -
  REFERENCE\_\allowbreak{}MODEL\_\allowbreak{}SPEC.md created (596
  lines)
\item
  \ok{} \textbf{Memory map} - MEMORY\_\allowbreak{}MAP.md created (360
  lines) with 4 KB minimum alignment
\item
  \ok{} \textbf{Verification plan phase alignment} - Restructured by
  phases (Phase 0-5 sections)
\item
  \ok{} \textbf{Interrupt support clarity} - Clearly stated as Phase 2+
  in \cref{cr:status:roadmap}
\end{enumerate}

\end{record}

\subsubsection{\texorpdfstring{\wip{} Current Issues}{ Current Issues}}\label{cr:status:phase-status:arrows-counterclockwise-current-issues}

\begin{record}

\textbf{None} - All Phase 0 documentation and implementation complete

\end{record}

\subsection{Next Actions}\label{cr:status:phase-status:next-actions}

\subsubsection{\texorpdfstring{Immediate --- DFT (\code{j41m}) and FPGA
(\code{hlrm})}{Immediate --- DFT () and FPGA ()}}\label{cr:status:phase-status:immediate-dft-j41m-and-fpga-hlrm}

\begin{record}

\textbf{Phases 0--6 are complete.} Phase 6 closed 2026-10-05 (Gate B).
Phase 7 M-a..M-c complete 2026-06-20.

\textbf{Current priority}: finish DFT Stage 1a (bead \bead{j41m.2},
still open; \code{j41m.9} is its GPU scan-reset test gap), then Stage 2
scan insertion (\code{j41m.8}, \code{j41m.4}--\code{j41m.6}). Open human
decisions in \src{docs/design/DFT_ARCHITECTURE.md}: Q-TRST and Q-STUB,
and the Sky130 CPU-macro re-harden, which depends on bead \bead{dud4}.

FPGA stages beyond Stage 0 are gated on a board choice (\code{hlrm}).
The Vivado Artix-7 run is in progress on a separate branch.

Golden specs: \src{docs/design/DFT_ARCHITECTURE.md},
\src{docs/FPGA_BOARD_SELECTION.md},
\src{docs/PHASE6_IP_EXPANSION_PLAN.md} (Phase 6, closed).

\textbf{Regression gates in CI (read 2026-10-15)}:

\begin{itemize}
\tightlist
\item
  \code{cocotb.yml} soc job:
  \texttt{PASS\_\allowbreak{}FLOOR:\allowbreak{}\ '759'} (its comment
  records 762 measured on \prref{276}). CPU job:
  \texttt{PASS\_\allowbreak{}FLOOR:\allowbreak{}\ '100'}.
\item
  \code{soc_coverage.yml} (nightly, cron \texttt{30\ 2\ *\ *\ *}):
  \texttt{PASS\_\allowbreak{}FLOOR:\allowbreak{}\ '759'},
  \texttt{CPU\_\allowbreak{}PASS\_\allowbreak{}FLOOR:\allowbreak{}\ '262'}.
  Line gate 95 \% on \src{rtl/soc}, \src{rtl/periph}, \src{rtl/npu} plus
  the control-signal toggle gate (bead \bead{s1cg}, \prref{246}); line gate
  95 \% also on \src{rtl/cpu}, \src{rtl/mem}, \src{rtl/gpu} (bead
  \bead{a5ze}, \prref{262}). Not a per-PR gate.
\item
  \code{qa-checks.yml}: struct-member select gate, blocking (\prref{275},
  bead \bead{ainf}); beads reopen gate.
\item
  \code{rtl-checks.yml}: \code{dft-scan-check} and
  \code{dft-scan-check-asap7} (\prref{281}, bead \bead{j41m.12}).
\end{itemize}

\end{record}

\subsubsection{Deferred / host-blocked (not Phase 6 work)}\label{cr:status:phase-status:deferred-host-blocked-not-phase-6-work}

\begin{record}

{\small\setlength{\tabcolsep}{4pt}

{
\begin{longtable}[]{@{}
  >{\raggedright\arraybackslash}p{\dimexpr(\linewidth-4\tabcolsep)*126/1000\relax}
  >{\raggedright\arraybackslash}p{\dimexpr(\linewidth-4\tabcolsep)*460/1000\relax}
  >{\raggedright\arraybackslash}p{\dimexpr(\linewidth-4\tabcolsep)*394/1000\relax}@{}}
\toprule\noalign{}
\begin{minipage}[b]{\linewidth}\raggedright
Bead
\end{minipage} & \begin{minipage}[b]{\linewidth}\raggedright
Item
\end{minipage} & \begin{minipage}[b]{\linewidth}\raggedright
Blocker
\end{minipage} \\
\midrule\noalign{}
\endhead
\bottomrule\noalign{}
\endlastfoot
\code{ma7} & ASAP7 CPU/GPU macros are Synlig-built with a proven branch
miscompile & ≥32 GB host (\code{2kn}) \\
\code{lxv} & ASAP7 CPU sv2v re-harden: routing-congestion regression &
≥32 GB host (\code{2kn}) \\
\code{8qn4} & ASAP7 SoC has never been run with GPIO & chained to
\code{ma7}/\code{lxv} \\
\code{e45j} & Sky130 post-RCX max-slew/\allowbreak{}max-cap; no post-RCX
repair stage exists in LibreLane & re-measured at Gate B; \textbf{not} a
Phase 6 exit criterion \\
\code{o1i} & Per-corner SRAM SPICE characterization
(\textasciitilde80-95 h) & host stability \\
\ghref{105}/\#106 & Sky130 GPU stages 3-4 & ≥32 GB RAM + \textasciitilde500
GB scratch \\
\code{hlrm} Stages 1+ & FPGA bring-up, firmware validation,
dual-oscillator CDC & a board choice (human decision) \\
\end{longtable}
}

}

\end{record}

\subsection{Documentation Structure}\label{cr:status:phase-status:documentation-structure}

\begin{record}

\begin{textblock}
docs/
├── ROADMAP.md                        # High-level project plan
├── PHASE_STATUS.md                   # This file - current status
├── design/
│   ├── PHASE0_ARCHITECTURE_SPEC.md   # Phase 0 CPU specification
│   ├── PHASE1_ARCHITECTURE_SPEC.md   # Phase 1 CPU specification (verified 2026-02-13)
│   ├── PHASE4_GPU_ARCHITECTURE_SPEC.md # Phase 4 GPU specification (frozen, ✅ complete)
│   ├── DESIGN_EXPECTATION.md         # High-level design goals
│   ├── RTL_DEFINITION.md             # Interface definitions
│   ├── GPU_MODEL.md                  # GPU execution model overview
│   ├── MEMORY_MAP.md                 # Address space allocation (complete 2026-01-17)
│   └── REFERENCE_MODEL_SPEC.md       # Python model specification (complete 2026-01-17)
└── verification/
    └── VERIFICATION_PLAN.md          # Verification strategy
\end{textblock}

\end{record}

\subsection{AI/Human Responsibilities for Phase 0 Remaining Tasks}\label{cr:status:phase-status:aihuman-responsibilities-for-phase-0-remaining-tasks}

\subsubsection{Python Reference Model Implementation}\label{cr:status:phase-status:python-reference-model-implementation}

\begin{record}

\textbf{AI may assist with}:

\begin{itemize}
\tightlist
\item
  Class structure and boilerplate code
\item
  Simple instruction implementations (ADD, SUB, AND, OR, XOR, SLL, SRL,
  SRA)
\item
  Register file and memory model scaffolding
\item
  Unit test generation and formatting
\end{itemize}

\textbf{Human must}:

\begin{itemize}
\tightlist
\item
  Implement complex instructions (BRANCH, LOAD, STORE, JAL, JALR)
\item
  Design and approve control flow logic
\item
  Verify instruction semantics match RISC-V specification
\item
  Cross-validate against spike simulator
\item
  Final code review and approval
\end{itemize}

\end{record}

\subsubsection{cocotb Infrastructure Setup}\label{cr:status:phase-status:cocotb-infrastructure-setup}

\begin{record}

\textbf{AI may assist with}:

\begin{itemize}
\tightlist
\item
  cocotb configuration files (Makefile, pyproject.toml)
\item
  AXI4-Lite and APB3 driver scaffolding
\item
  Test harness boilerplate
\item
  Monitor and scoreboard templates
\end{itemize}

\textbf{Human must}:

\begin{itemize}
\tightlist
\item
  Review test strategy alignment with VERIFICATION\_\allowbreak{}PLAN.md
\item
  Approve infrastructure design decisions
\item
  Validate driver implementations against protocol specs
\end{itemize}

\end{record}

\subsubsection{Final Specification Review}\label{cr:status:phase-status:final-specification-review}

\begin{record}

\textbf{HUMAN-ONLY}:

\begin{itemize}
\tightlist
\item
  Review all specifications for consistency and completeness
\item
  Approve Phase 0 completion
\item
  Authorize transition to Phase 1 (RTL implementation)
\end{itemize}

\end{record}

\subsection{Key Decisions}\label{cr:status:phase-status:key-decisions}

\subsubsection{Architecture Decisions}\label{cr:status:phase-status:architecture-decisions}

\begin{record}

\begin{itemize}
\tightlist
\item
  \textbf{ISA}: RV32I subset (no CSR, no MMU, no compressed)
\item
  \textbf{Initial pipeline}: Single-cycle (Phase 1), 5-stage pipeline
  (Phase 2)
\item
  \textbf{Memory interface}: AXI4-Lite for main memory, APB3 for debug
\item
  \textbf{Debug strategy}: APB3 slave interface with
  halt/\allowbreak{}resume/\allowbreak{}step
\item
  \textbf{Verification}: Python reference model + cocotb + pyuvm
\end{itemize}

\end{record}

\subsubsection{Process Decisions}\label{cr:status:phase-status:process-decisions}

\begin{record}

\begin{itemize}
\tightlist
\item
  \textbf{Specification-first}: All specs finalized before RTL
\item
  \textbf{No architecture drift}: Phase N cannot start until Phase N-1
  exits
\item
  \textbf{AI boundaries}: AI assists with RTL/\allowbreak{}tests, humans
  own architecture
\item
  \textbf{Verification requirement}: Reference model must match RTL
\end{itemize}

\end{record}

\subsection{Contact}\label{cr:status:phase-status:contact}

\begin{record}

For questions about project status or phase transitions, refer to
\cref{cr:status:roadmap} or the latest git commits.

\end{record}

% Converted from docs/ROADMAP.md by tools/docs/md2tex.py; edit this file, not the original.
\section{Project Roadmap}\label{cr:status:roadmap}

\subsection{Project Expectations}\label{cr:status:roadmap:project-expectations}

\begin{record}

\begin{enumerate}
\def\labelenumi{\arabic{enumi}.}
\tightlist
\item
  Create simple micro-processor with RV32I subset
\item
  Upgrade micro-processor to become pipelined CPU
\item
  Develop memory system \& caches
\item
  Create a simple GPU compute engine
\item
  Integrate all IPs into SoC with some peripherals
\end{enumerate}

\end{record}

\subsection{OpenROAD Back-End Integration Strategy}\label{cr:status:roadmap:openroad-back-end-integration-strategy}

\begin{record}

Starting from Phase 1, \textbf{every IP block} will go through the
complete OpenROAD physical design flow to ensure silicon-readiness. This
approach provides:

\begin{enumerate}
\def\labelenumi{\arabic{enumi}.}
\tightlist
\item
  \textbf{Early timing validation}: Identify critical paths before
  tape-out
\item
  \textbf{Power budgeting}: Understand power consumption per block
\item
  \textbf{Area estimation}: Guide floorplanning decisions
\item
  \textbf{Reality check}: Ensure RTL is physically implementable
\end{enumerate}

\end{record}

\subsubsection{Back-End Flow Overview}\label{cr:status:roadmap:back-end-flow-overview}

\begin{record}

For each new IP (CPU core, cache, GPU, peripherals):

\begin{textblock}
RTL Design
    ↓
Synthesis (Yosys/OpenROAD)
    ↓
Floorplan (OpenROAD)
    ↓
Placement (OpenROAD)
    ↓
Clock Tree Synthesis (TritonCTS)
    ↓
Routing (TritonRoute)
    ↓
Parasitic Extraction (OpenRCX)
    ↓
Static Timing Analysis (OpenSTA)
    ↓
Power Analysis (OpenROAD)
    ↓
Physical Verification (KLayout/Magic)
    ↓
Gate-Level Simulation (Verilator + UPF)
\end{textblock}

\end{record}

\subsubsection{Power-Aware Verification}\label{cr:status:roadmap:power-aware-verification}

\begin{record}

\textbf{UPF (Unified Power Format)} is used to specify power intent:

\begin{itemize}
\tightlist
\item
  Power domains and supply networks
\item
  Isolation strategies for domain boundaries
\item
  Retention strategies for state preservation
\item
  Power state tables (PST) for mode transitions
\end{itemize}

\textbf{Power-aware simulation} validates:

\begin{itemize}
\tightlist
\item
  Isolation cell functionality
\item
  Level shifter correctness
\item
  Retention register behavior
\item
  Power sequencing
\end{itemize}

\end{record}

\subsubsection{Timing Verification}\label{cr:status:roadmap:timing-verification}

\begin{record}

\textbf{SDC (Synopsys Design Constraints)} specifies timing
requirements:

\begin{itemize}
\tightlist
\item
  Clock definitions (period, uncertainty, jitter)
\item
  Input/\allowbreak{}output delays relative to clock
\item
  False paths (paths that don\textquotesingle t need timing checks)
\item
  Multi-cycle paths (paths with \textgreater1 cycle budget)
\item
  Clock domain crossing constraints
\end{itemize}

\textbf{Static Timing Analysis (STA)} verifies:

\begin{itemize}
\tightlist
\item
  Setup timing (data arrives before clock edge)
\item
  Hold timing (data stable after clock edge)
\item
  Clock skew and latency
\item
  Timing across all process corners (fast/\allowbreak{}slow)
\end{itemize}

\end{record}

\subsubsection{Target Technologies}\label{cr:status:roadmap:target-technologies}

\begin{record}

\begin{itemize}
\tightlist
\item
  \textbf{Sky130}: Open-source 130nm PDK (SkyWater + Google)
\item
  \textbf{ASAP7}: Predictive 7nm FinFET PDK (Arizona State University)
\end{itemize}

Both support the OpenROAD flow and are suitable for
academic/\allowbreak{}research projects.

\end{record}

\subsection{Phase}\label{cr:status:roadmap:phase}

\subsubsection{\texorpdfstring{Phase 0 --- Foundations \ok{} COMPLETE
(2026-01-18)}{Phase 0 --- Foundations  COMPLETE (2026-01-18)}}\label{cr:status:roadmap:phase-0-foundations-white-check-mark-complete-2026-01-18}

\begin{record}

\begin{itemize}
\tightlist
\item
  Goals

  \begin{itemize}
  \tightlist
  \item
    Define what correctness means
  \item
    Prevent architecture drift later
  \end{itemize}
\item
  Human-led (must)

  \begin{itemize}
  \tightlist
  \item
    Define RV32I subset (no CSR, no MMU, no compressed)
  \item
    Define:

    \begin{itemize}
    \tightlist
    \item
      Reset behavior
    \item
      Trap behavior
    \item
      Memory ordering rules
    \end{itemize}
  \item
    Decide pipeline strategy (single-cycle vs multi-cycle)
  \end{itemize}
\item
  AI-assisted

  \begin{itemize}
  \tightlist
  \item
    Generate:

    \begin{itemize}
    \tightlist
    \item
      Project skeleton
    \item
      Lint rules
    \item
      Basic SV coding guidelines
    \end{itemize}
  \item
    Draft documentation templates
  \end{itemize}
\end{itemize}

\ok{} Exit criteria MET

\begin{itemize}
\tightlist
\item
  \ok{} Written ISA + microarchitecture spec
  (PHASE0\_\allowbreak{}ARCHITECTURE\_\allowbreak{}SPEC.md)
\item
  \ok{} No RTL (as planned)
\item
  \ok{} All 7 specifications approved
\item
  \ok{} Python reference models validated (66/66 tests passing)
\item
  \ok{} cocotb test infrastructure ready
\end{itemize}

\end{record}

\subsubsection{\texorpdfstring{Phase 1 --- Minimal RV32I core \ok{}
COMPLETE
(2026-02-13)}{Phase 1 --- Minimal RV32I core  COMPLETE (2026-02-13)}}\label{cr:status:roadmap:phase-1-minimal-rv32i-core-white-check-mark-complete-2026-02-13}

\begin{record}

\begin{itemize}
\tightlist
\item
  Architecture

  \begin{itemize}
  \tightlist
  \item
    Single-issue
  \item
    In-order
  \item
    Single-cycle with AXI stalls
  \item
    Blocking loads/\allowbreak{}stores
  \item
    \textbf{No interrupts} (deferred to Phase 2)
  \item
    AXI4-Lite master (unified instruction/\allowbreak{}data bus)
  \item
    APB3 slave (debug interface with
    halt/\allowbreak{}resume/\allowbreak{}step/\allowbreak{}breakpoints)
  \end{itemize}
\end{itemize}

\textbf{RTL (SystemVerilog)} \ok{} ALL COMPLETE

\begin{itemize}
\tightlist
\item
  \ok{} 8 modules implemented (\textasciitilde1,900 lines)

  \begin{itemize}
  \tightlist
  \item
    \ok{} ALU module (rv32i\_\allowbreak{}alu.sv)
  \item
    \ok{} Register file (rv32i\_\allowbreak{}regfile.sv)
  \item
    \ok{} Instruction decoder (rv32i\_\allowbreak{}decode.sv)
  \item
    \ok{} Immediate generator
    (rv32i\_\allowbreak{}imm\_\allowbreak{}gen.sv)
  \item
    \ok{} Branch comparator
    (rv32i\_\allowbreak{}branch\_\allowbreak{}comp.sv)
  \item
    \ok{} Control FSM (rv32i\_\allowbreak{}control.sv)
  \item
    \ok{} Core wrapper (rv32i\_\allowbreak{}core.sv)
  \item
    \ok{} CPU top-level (rv32i\_\allowbreak{}cpu\_\allowbreak{}top.sv)
  \end{itemize}
\end{itemize}

\textbf{Verification (Python-first)} \ok{} ALL COMPLETE

\begin{itemize}
\tightlist
\item
  \ok{} cocotb test infrastructure

  \begin{itemize}
  \tightlist
  \item
    \ok{} AXI4-Lite memory models
  \item
    \ok{} APB3 debug interface drivers
  \item
    \ok{} Scoreboard validation
  \end{itemize}
\item
  \ok{} Comprehensive test suites

  \begin{itemize}
  \tightlist
  \item
    \ok{} 6/6 smoke tests passing
  \item
    \ok{} 37/37 ISA compliance tests
  \item
    \ok{} 10,000 random instructions (0 failures)
  \item
    \ok{} 11/11 AXI protocol tests
  \item
    \ok{} 6/6 debug interface tests
  \end{itemize}
\item
  \ok{} Coverage

  \begin{itemize}
  \tightlist
  \item
    \ok{} 37/37 instructions (100\%)
  \item
    \ok{} 8/8 FSM states (100\%)
  \end{itemize}
\end{itemize}

\textbf{Physical Design \& Power (Optional)}

\begin{itemize}
\tightlist
\item
  OpenROAD flow scripts ready
\item
  Phase 1 constraints available (phase1\_\allowbreak{}cpu.sdc,
  phase1\_\allowbreak{}cpu.upf)
\item
  Can be run for Phase 1 baseline metrics
\end{itemize}

\ok{} Exit criteria MET (9/9)

\begin{itemize}
\tightlist
\item
  \ok{} Smoke tests passing (6/6 with scoreboard)
\item
  \ok{} Scoreboard mismatches: 0
\item
  \ok{} Instruction coverage: 37/37 RV32I instructions (100\%)
\item
  \ok{} Random instruction tests: 10,000 instructions, 0 failures
\item
  \ok{} AXI protocol tests: 11/11 passing
\item
  \ok{} Debug interface tests: 6/6 passing (single-step, BP0, BP1,
  GPR/PC write)
\item
  \ok{} Code coverage: \textgreater95\% (Verilator annotated reports)
\item
  \ok{} State coverage: 8/8 FSM states (100\%)
\item
  \ok{} Failing random seeds: 0 (100/100 seeds pass)
\item
  \ok{} Archived to \src{micro_p/} directory
\end{itemize}

\end{record}

\subsubsection{\texorpdfstring{Phase 2 --- Pipelined CPU \ok{} COMPLETE
(2026-03-08)}{Phase 2 --- Pipelined CPU  COMPLETE (2026-03-08)}}\label{cr:status:roadmap:phase-2-pipelined-cpu-white-check-mark-complete-2026-03-08}

\begin{record}

\begin{itemize}
\tightlist
\item
  Architecture

  \begin{itemize}
  \tightlist
  \item
    5-stage pipeline
    (IF/\allowbreak{}ID/\allowbreak{}EX/\allowbreak{}MEM/\allowbreak{}WB)
  \item
    In-order execution
  \item
    Hazard detection + stalling + forwarding
  \item
    \textbf{Interrupt support added} (M-mode: timer + external)
  \item
    \textbf{CSR instructions}
    (CSRRW/\allowbreak{}S/\allowbreak{}C/\allowbreak{}I variants)
  \item
    Same AXI4-Lite and APB3 interfaces as Phase 1
  \item
    AXI arbiter for IF/MEM priority
  \item
    Achieved: 75 MHz on Sky130 130nm (200 MHz target not met due to PDK
    limitations)
  \end{itemize}
\end{itemize}

\textbf{RTL} \ok{} ALL COMPLETE

\begin{itemize}
\tightlist
\item
  \ok{} 14 modules implemented (\textasciitilde3,000 lines)

  \begin{itemize}
  \tightlist
  \item
    \ok{} 5 pipeline stage modules (IF, ID, EX, MEM, WB)
  \item
    \ok{} Pipeline package
    (rv32i\_\allowbreak{}pipeline\_\allowbreak{}pkg.sv)
  \item
    \ok{} Hazard detection unit
    (rv32i\_\allowbreak{}hazard\_\allowbreak{}unit.sv)
  \item
    \ok{} Forwarding unit
    (rv32i\_\allowbreak{}forwarding\_\allowbreak{}unit.sv)
  \item
    \ok{} CSR file (rv32i\_\allowbreak{}csr\_\allowbreak{}file.sv)
  \item
    \ok{} Interrupt controller
    (rv32i\_\allowbreak{}interrupt\_\allowbreak{}ctrl.sv)
  \item
    \ok{} AXI arbiter (rv32i\_\allowbreak{}axi\_\allowbreak{}arbiter.sv)
  \item
    \ok{} Updated decoder for CSR instructions
  \item
    \ok{} Reused: ALU, regfile, imm\_gen, branch\_\allowbreak{}comp
  \end{itemize}
\end{itemize}

\textbf{Verification} \ok{} ALL COMPLETE

\begin{itemize}
\tightlist
\item
  \ok{} Pipeline hazard tests --- 16/16
  (RAW/\allowbreak{}WAR/\allowbreak{}WAW, forwarding)
\item
  \ok{} Interrupt and CSR tests --- 12/12
\item
  \ok{} Debug interface tests --- 6/6 (pipeline drain)
\item
  \ok{} Random instruction regression --- 50,000 instructions, 0
  failures
\item
  \ok{} AXI arbiter protocol tests --- 12/12
\item
  \ok{} Fault injection tests --- 7/7
\item
  \ok{} Total: 111/111 tests passing
\end{itemize}

\textbf{Physical Design \& Power} \ok{} COMPLETE

\begin{itemize}
\tightlist
\item
  \ok{} Phase 2 constraints: \src{pnr/constraints/phase2_cpu.sdc},
  \src{pnr/constraints/phase2_cpu.upf}
\item
  \ok{} Backend flow executed on Sky130 130nm
\item
  \ok{} \textbf{75 MHz achieved} (200 MHz target not met --- Sky130 PDK
  limitation)
\end{itemize}

\notemark{} AI contribution \textasciitilde60\%

\ok{} Exit criteria MET

\begin{itemize}
\tightlist
\item
  \ok{} Zero false commits (scoreboard validation, 50k random
  instructions)
\item
  \ok{} Pipeline hazard tests passing (16/16)
\item
  \ok{} Interrupt tests passing (12/12)
\item
  \ok{} All 111/111 Phase 2 tests passing
\item
  \ok{} Timing closure at 75 MHz on Sky130
\end{itemize}

\end{record}

\subsubsection{\texorpdfstring{Phase 3 --- Memory system \& caches \ok{}
COMPLETE
(2026-05-21)}{Phase 3 --- Memory system \& caches  COMPLETE (2026-05-21)}}\label{cr:status:roadmap:phase-3-memory-system-caches-white-check-mark-complete-2026-05-21}

\begin{record}

\begin{itemize}
\tightlist
\item
  Architecture

  \begin{itemize}
  \tightlist
  \item
    I-Cache: 4 KB direct-mapped, 16-byte lines, 256 sets, read-only,
    FENCE.I invalidation
  \item
    D-Cache: 4 KB direct-mapped, 16-byte lines, 256 sets, write-back +
    write-allocate
  \item
    External interface: AXI4-Lite (4 separate transactions per refill
    --- no burst)
  \item
    Cache arbiter: D-cache priority over I-cache (replaces Phase 2 AXI
    arbiter)
  \item
    Architecture spec: \src{docs/design/PHASE3_ARCHITECTURE_SPEC.md} ---
    APPROVED (2026-03-08)
  \end{itemize}
\end{itemize}

\textbf{RTL} \ok{} COMPLETE (\textasciitilde1,424 lines)

\begin{itemize}
\tightlist
\item
  \ok{} New modules in \src{rtl/mem/}:

  \begin{itemize}
  \tightlist
  \item
    \ok{} \code{rv32i_cache_pkg.sv} --- parameters, types (4 KB, 16-byte
    lines, direct-mapped)
  \item
    \ok{} \code{rv32i_icache.sv} --- I-cache with
    IDLE/\allowbreak{}REFILL/\allowbreak{}DONE FSM, FENCE.I (575 lines)
  \item
    \ok{} \code{rv32i_dcache.sv} --- D-cache with
    IDLE/\allowbreak{}TAG\_\allowbreak{}CHECK/\allowbreak{}WRITEBACK/\allowbreak{}REFILL/\allowbreak{}DONE
    FSM, write-allocate (618 lines)
  \item
    \ok{} \code{rv32i_cache_arbiter.sv} --- D\$ priority AXI arbiter
    (183 lines)
  \end{itemize}
\item
  \ok{} Modified CPU modules:

  \begin{itemize}
  \tightlist
  \item
    \ok{} \code{rv32i_decode.sv} --- FENCE.I decode (opcode=0x0F,
    funct3=001)
  \item
    \ok{} \code{rv32i_pipeline_pkg.sv} --- \code{fence_i} field added to
    \code{id_ex_reg_t}
  \item
    \ok{} \code{rv32i_pipeline_if.sv} --- cache interface (replaces AXI
    state machine)
  \item
    \ok{} \code{rv32i_pipeline_mem.sv} --- cache interface + FENCE.I
    (replaces AXI state machine)
  \item
    \ok{} \code{rv32i_hazard_unit.sv} --- rename stall signals
    (\code{if_cache_stall}, \code{mem_cache_stall})
  \item
    \ok{} \code{rv32i_core.sv} --- cache + arbiter integration
    (instantiated at lines 742/\allowbreak{}762/\allowbreak{}794)
  \end{itemize}
\end{itemize}

\textbf{Verification} \ok{} COMPLETE

\begin{itemize}
\tightlist
\item
  \ok{} Python cache reference model (\src{tb/models/cache_model.py})
\item
  \ok{} I-cache unit tests: \src{tb/cocotb/mem/test_icache.py} ---
  \textbf{7/7 pass}
\item
  \ok{} D-cache unit tests: \src{tb/cocotb/mem/test_dcache.py} ---
  \textbf{8/8 pass}
\item
  \ok{} Integration: \src{tb/cocotb/cpu/test_cache_integration.py} ---
  \textbf{5/5 pass}
\item
  \ok{} Full regression: \textbf{139/139 tests pass} (119 Phase 2 + 20
  Phase 3 cache tests)
\item
  \ok{} FENCE.I ordering verified; no livelock or deadlock scenarios
  found
\end{itemize}

\textbf{Physical Design \& Power (OpenROAD back-end)} \ok{} COMPLETE
(ASAP7 sign-off)

\begin{itemize}
\tightlist
\item
  \ok{} Phase 2+3 RTL signed off on ASAP7 7nm --- \textbf{Run 43
  (2026-05-20)}:

  \begin{itemize}
  \tightlist
  \item
    \textbf{1418 MHz fmax}, 27.27 mW power, 3,844 µm² stdcell area
  \item
    \textbf{0 DRC, 0 antenna violations}, +5.97 ps setup slack, +22.54
    ps hold slack
  \item
    \warn{} Timing/\allowbreak{}power unvalidated post-bead-\code{8f3}
    (STA-zero-wire defect, 2026-09-16); Run 43 artifacts wiped, not
    re-checkable --- bead \bead{0p6}
  \item
    Full run history: \cref{cr:run:asap7}
  \end{itemize}
\item
  Sky130 realistic ceiling: 75--120 MHz (PDK-limited; Phase 2 achieved
  75 MHz)
\end{itemize}

\notemark{} AI contribution \textasciitilde40\%

\ok{} Exit criteria --- ALL MET

\begin{itemize}
\tightlist
\item
  \ok{} No deadlocks in cache FSMs (verified by 139/139 regression +
  integration tests)
\item
  \ok{} Correct behavior under randomized memory latency
\item
  \ok{} FENCE.I invalidation semantics implemented and verified
\item
  \ok{} Timing closure on ASAP7 (1418 MHz, 0 violations)
\end{itemize}

\end{record}

\subsubsection{Phase 4 --- "GPU-lite" compute engine (research-level)}\label{cr:status:roadmap:phase-4-gpu-lite-compute-engine-research-level}

\begin{record}

\warn{} This is not a CPU extension.

\textbf{Prerequisites}: Phase 3 complete (pipelined CPU with caches and
interrupts)

\textbf{Scope (keep it sane)}

\begin{itemize}
\tightlist
\item
  Single compute unit

  \begin{itemize}
  \tightlist
  \item
    SIMD/SIMT lanes
  \item
    No graphics
  \item
    No preemption
  \item
    \textbf{GPU interrupt output} to notify CPU of kernel completion
    (requires Phase 2+ CPU)
  \end{itemize}
\end{itemize}

\textbf{RTL}

\begin{itemize}
\tightlist
\item
  AI SHOULD

  \begin{itemize}
  \tightlist
  \item
    Vector ALU blocks
  \item
    Lane control logic
  \item
    Simple warp scheduler templates
  \end{itemize}
\item
  Human MUST

  \begin{itemize}
  \tightlist
  \item
    Define execution model (SIMD vs SIMT)
  \item
    Handle divergence
  \item
    Define memory coalescing rules
  \end{itemize}
\end{itemize}

\textbf{Verification}

\begin{itemize}
\tightlist
\item
  Python kernel model

  \begin{itemize}
  \tightlist
  \item
    Vector reference execution
  \item
    Divergence checker
  \end{itemize}
\item
  AI SHOULD

  \begin{itemize}
  \tightlist
  \item
    Generate kernel tests
  \item
    Build comparison harness
  \end{itemize}
\item
  Human MUST

  \begin{itemize}
  \tightlist
  \item
    Debug race conditions
  \item
    Validate numerical correctness
  \end{itemize}
\end{itemize}

\textbf{Physical Design \& Power (OpenROAD back-end)}

\begin{itemize}
\tightlist
\item
  GPU compute unit as separate power domain
\item
  UPF for GPU power gating (idle power savings)
\item
  Vector ALU datapath optimization
\item
  Register file (8-lane SIMT) physical design
\item
  Warp scheduler timing constraints
\item
  Power analysis for compute kernels
\item
  Dynamic voltage/\allowbreak{}frequency scaling (DVFS) infrastructure
\item
  Thermal analysis for sustained compute workloads
\item
  GPU-CPU interface timing constraints
\end{itemize}

\notemark{} AI contribution \textasciitilde30--40\%

\ok{} Exit criteria (updated)

\begin{itemize}
\tightlist
\item
  \textbf{GPU power domain functional}
\item
  \textbf{Timing closure for compute datapath}
\item
  \textbf{Power consumption validated per warp}
\end{itemize}

\end{record}

\subsubsection{\texorpdfstring{Phase 5 --- SoC integration (AI shines
again) --- \ok{} COMPLETE 2026-06-24 (ASAP7 SoC 571 MHz / 62.9 mW
sign-off \warn{} timing/\allowbreak{}power unvalidated
post-bead-\code{8f3}; current RTL re-measured on the fixed flow
2026-09-18, bead \code{0d0}: −727 ps / 781 violators / hold clean /
283.6 mW, still not closed --- see
\cref{cr:run:phase5})}{Phase 5 --- SoC integration (AI shines again) ---  COMPLETE 2026-06-24 (ASAP7 SoC 571 MHz / 62.9 mW sign-off  timing/power unvalidated post-bead-; current RTL re-measured on the fixed flow 2026-09-18, bead : −727 ps / 781 violators / hold clean / 283.6 mW, still not closed --- see )}}\label{cr:status:roadmap:phase-5-soc-integration-ai-shines-again-white-check-mark-complete-2026-06-24-asap7-soc-571-mhz-629-mw-sign-off-warning-timingpower-unvalidated-post-bead-8f3-current-rtl-re-measured-on-the-fixed-flow-2026-09-18-bead-0d0-727-ps-781-violators-hold-clean-2836-mw-still-not-closed-see-docsphase5-run-historymd}

\begin{record}

\begin{itemize}
\tightlist
\item
  Components

  \begin{itemize}
  \tightlist
  \item
    CPU + GPU-lite
  \item
    Interconnect
  \item
    DMA
  \item
    UART / SPI / Timer
  \item
    Boot ROM
  \end{itemize}
\end{itemize}

\textbf{RTL}

\begin{itemize}
\tightlist
\item
  AI SHOULD

  \begin{itemize}
  \tightlist
  \item
    Address maps
  \item
    Interconnect glue
  \item
    Peripheral controllers
  \end{itemize}
\item
  Human MUST

  \begin{itemize}
  \tightlist
  \item
    Define memory map
  \item
    Decide arbitration rules
  \item
    Own integration decisions
  \end{itemize}
\end{itemize}

\textbf{Verification}

\begin{itemize}
\tightlist
\item
  cocotb system tests

  \begin{itemize}
  \tightlist
  \item
    Boot tests
  \item
    Firmware execution
  \item
    Peripheral stress tests
  \end{itemize}
\item
  AI SHOULD

  \begin{itemize}
  \tightlist
  \item
    Write firmware tests
  \item
    Automate regressions
  \item
    Generate fault injection
  \end{itemize}
\end{itemize}

\textbf{Physical Design \& Power (OpenROAD back-end - Full SoC)}

\begin{itemize}
\tightlist
\item
  Top-level SoC floorplan
\item
  Multiple power domains (CPU, GPU, peripherals, always-on)
\item
  UPF for system-level power management

  \begin{itemize}
  \tightlist
  \item
    Domain isolation cells
  \item
    Level shifters for voltage domains
  \item
    Retention strategies for low-power modes
  \end{itemize}
\item
  Hierarchical P\&R (block-level then top-level)
\item
  Clock distribution network (multiple clock domains)
\item
  Power delivery network (PDN) analysis
\item
  IR drop analysis and mitigation
\item
  Electromigration (EM) analysis
\item
  Sign-off timing analysis (multi-corner, multi-mode)
\item
  Power estimation for system workloads
\item
  Tape-out quality checks (DRC, LVS, antenna)
\end{itemize}

\notemark{} AI contribution \textasciitilde70\%

\ok{} Exit criteria

\begin{itemize}
\tightlist
\item
  Boots software
\item
  Passes long random regressions
\item
  \textbf{Full SoC timing closure (all corners)}
\item
  \textbf{Power delivery verified (IR drop \textless{} 5\%)}
\item
  \textbf{All power modes functional (active, idle, sleep)}
\item
  \textbf{DRC/LVS clean}
\item
  \textbf{Tape-out ready (if target technology selected)}
\end{itemize}

\end{record}

\subsubsection{\texorpdfstring{Phase 6+ --- IP Expansion \& Technology
Node Exploration (\ok{} 6a--6c complete; Gate B \ok{}
2026-10-05)}{Phase 6+ --- IP Expansion \& Technology Node Exploration ( 6a--6c complete; Gate B  2026-10-05)}}\label{cr:status:roadmap:phase-6-ip-expansion-technology-node-exploration-white-check-mark-6a6c-complete-gate-b-white-check-mark-2026-10-05}

\begin{record}

\textbf{Prerequisites}: Phase 5 complete (full SoC validated)

\textbf{Status (2026-10-15)}: Phase 6 complete. 6a (GPIO, PWM, WDT,
TRNG, I2C), 6b (CRYPTO) and 6c (NPU) are built; the batched Sky130
harden (Gate B) closed 2026-10-05 (bead \bead{f7vs.15}, \prref{214}). ASAP7
is out of scope for Phase 6 (beads \code{ma7}, \code{lxv}, \code{2kn}).

\begin{quote}
\textbf{Golden spec:} \src{docs/PHASE6_IP_EXPANSION_PLAN.md} --- full
register maps, address/\allowbreak{}IRQ allocation, per-item acceptance
criteria and non-goals. This section is the roadmap-level summary; the
plan is authoritative where they differ.
\end{quote}

\end{record}

\paragraph{6a. Additional Peripherals (APB4 slaves on the Phase 5 APB
sub-tree)}\label{cr:status:roadmap:6a-additional-peripherals-apb4-slaves-on-the-phase-5-apb-sub-tree}

\begin{record}

\begin{quote}
\textbf{Bus correction (2026-09-20):} this heading previously read
"AXI4-Lite slaves". \src{docs/PERIPHERAL_BUS_EVALUATION.md} settled the
peripheral-bus standard as \textbf{APB4} behind
\code{apb4_register_bank} --- "Phase 6 peripherals become trivial APB
drop-ins". The per-node frequency columns below are \emph{projections},
not measured results.
\end{quote}

{\small\setlength{\tabcolsep}{4pt}

{
\begin{longtable}[]{@{}
  >{\raggedright\arraybackslash}p{\dimexpr(\linewidth-10\tabcolsep)*127/1000\relax}
  >{\raggedright\arraybackslash}p{\dimexpr(\linewidth-10\tabcolsep)*292/1000\relax}
  >{\raggedright\arraybackslash}p{\dimexpr(\linewidth-10\tabcolsep)*112/1000\relax}
  >{\raggedright\arraybackslash}p{\dimexpr(\linewidth-10\tabcolsep)*112/1000\relax}
  >{\raggedright\arraybackslash}p{\dimexpr(\linewidth-10\tabcolsep)*112/1000\relax}
  >{\raggedright\arraybackslash}p{\dimexpr(\linewidth-10\tabcolsep)*224/1000\relax}@{}}
\toprule\noalign{}
\begin{minipage}[b]{\linewidth}\raggedright
IP
\end{minipage} & \begin{minipage}[b]{\linewidth}\raggedright
RTL file
\end{minipage} & \begin{minipage}[b]{\linewidth}\raggedright
Sky130
\end{minipage} & \begin{minipage}[b]{\linewidth}\raggedright
FreePDK45
\end{minipage} & \begin{minipage}[b]{\linewidth}\raggedright
ASAP7
\end{minipage} & \begin{minipage}[b]{\linewidth}\raggedright
Note
\end{minipage} \\
\midrule\noalign{}
\endhead
\bottomrule\noalign{}
\endlastfoot
GPIO controller & \src{rtl/periph/gpio_controller.sv} & \ok{} 75 MHz &
\ok{} 400 MHz & \ok{} 1 GHz & ← \textbf{RTL landed} (bead \bead{ckc}) \\
I2C controller & \src{rtl/periph/i2c_controller.sv} & \ok{} 75 MHz &
\ok{} 400 MHz & \ok{} 1 GHz & \\
PWM controller & \src{rtl/periph/pwm_controller.sv} & \ok{} 75 MHz &
\ok{} 400 MHz & \ok{} 1 GHz & ← \textbf{RTL landed} (bead
\bead{f7vs.6}) \\
Watchdog timer & \src{rtl/periph/watchdog_timer.sv} & \ok{} 75 MHz &
\ok{} 400 MHz & \ok{} 1 GHz & ← \textbf{RTL landed} (bead
\bead{f7vs.7}) \\
TRNG & \src{rtl/periph/trng.sv} & \ok{} RTL portable & \ok{} RTL
portable & \ok{} RTL portable & ← \textbf{RTL landed} (bead
\bead{f7vs.8}) \\
\end{longtable}
}

}

Note (corrected 2026-09-28, bead \bead{f7vs.1}): the TRNG \textbf{RTL is
fully portable}. Sky130-exclusivity is confined to a swappable entropy
sub-module selected by \texttt{`ifdef} --- \code{trng_ro_sky130.sv}
(ring oscillator) vs. the default \code{trng_lfsr_entropy.sv}, which is
deterministic and \textbf{explicitly not cryptographic} (it forces
\texttt{TRNG\_\allowbreak{}STATUS.\allowbreak{}INSECURE\ =\allowbreak{}\ 1}).
What is Sky130-exclusive is the \textbf{entropy-quality claim}, not the
RTL, so \code{trng.sv} has zero top-level ports and sits in every file
list and every PD flow. See \src{docs/PHASE6_IP_EXPANSION_PLAN.md} §7.

\textbf{GPIO controller --- status (Phase 6a, bead \bead{ckc},
2026-09-20).} RTL + SoC integration + cocotb suite landed:
\code{test_gpio} \textbf{16/16}, \code{soc_all} \textbf{217/217} across
27 suites (was 201 pre-GPIO; CI \code{PASS_FLOOR} raised 170 → 210).
APB4 slave at \texttt{0x2000\_\allowbreak{}A000–AFFF} (APB index 7,
\code{APB_N_SLAVES} 7 → 8); \code{AXIL_APB_LIMIT} and
\code{PERIPH_LIMIT} extended to \code{0x2000_AFFF}. 32 pins exposed as
the unidirectional triplet \code{gpio_out_o} / \code{gpio_oe_o} /
\code{gpio_in_i} --- this RTL tree contains no tristate and the Sky130
SoC hardens as a core macro with no pad ring, so the bidirectional merge
is left to pad-ring integration. Per-pin level/\allowbreak{}edge
interrupt with polarity; new \code{interrupt_controller} source bit 5
(\code{N_SOURCES} 5 → 6). Wired into the Sky130 and ASAP7 sv2v file
lists. \textbf{Hardened on Sky130 2026-09-26}
(\code{RUN_2026-09-26_00-07-59}, bead \bead{00ef}): Netgen LVS PASSED,
routing DRC 0, setup AND hold both clean at all 9 corners. \warn{} The
SoC clock is \textbf{25 ns / 40 MHz} (\src{pnr/sky130/soc/config.json}),
so this measures GPIO \emph{inside a 40 MHz SoC} --- the 75 MHz column
above is a per-peripheral projection and is still NOT a measurement.
Detail: \src{docs/SKY130_REAL_DRC_LVS_EVALUATION.md}.

\end{record}

\paragraph{6b. AES-128 + SHA-256 Accelerator}\label{cr:status:roadmap:6b-aes-128-sha-256-accelerator}

\begin{record}

\begin{quote}
\textbf{Bus correction (2026-09-28, bead \bead{f7vs.1}):} this section
previously said "AXI4-Lite slave". The peripheral-bus standard is
\textbf{APB4} behind \code{apb4_register_bank}
(\code{docs/PERIPHERAL_BUS_EVALUATION.md:5,43} names AES-SHA
explicitly). The 2026-09-20 correction reached the 6a heading only and
never propagated here. Golden spec:
\src{docs/PHASE6_IP_EXPANSION_PLAN.md} §2.
\end{quote}

\ok{} \textbf{LANDED 2026-10-03 (bead \bead{f7vs.10})} ---
\src{rtl/periph/crypto_accel.sv} + \code{aes128_core.sv} +
\code{sha256_core.sv}, APB4 slave index 12 at \code{0x2000_F000-FFFF},
IRQ bit 10. \code{test_crypto} 41/41, \code{test_soc_crypto} 1/1.

\begin{itemize}
\tightlist
\item
  \textbf{APB4 slave} (control/\allowbreak{}status,
  key/\allowbreak{}IV/\allowbreak{}digest registers), one 4 KB slot at
  \code{0x2000_F000}
\item
  Bulk data by \textbf{programmed I/O over APB}. An AXI4
  master/\allowbreak{}slave path is an explicit \textbf{non-goal of
  Phase 6b} --- it would need \code{SOC_N_SLAVES} 3 → 4 and a change to
  \code{axi4_crossbar}. Tracked as Phase 6b-2, not scheduled.
\item
  \textbf{AES-128 encrypt-only, exposing ECB \emph{and} CTR.} CTR makes
  an encrypt-only core a complete cipher in both directions, so there is
  no inverse S-box and no inverse key schedule. Iterative 128-bit
  datapath, on-the-fly key schedule, \code{SBOX_PARALLEL=16} →
  \textbf{11 cycles/\allowbreak{}block} (the documented fold-to-4
  fallback is 41 and bit-identical; Gate A measured the default at 277
  875 µm² --- \textasciitilde28 \% of the SoC\textquotesingle s stdcell
  area, not the 3.59 \% first recorded, which divided by stdcell
  \textbf{plus} macro area; core utilisation only moves 39.0 \% →
  \textasciitilde43.7 \% with CRYPTO + NPU, so the fallback was still
  not taken. See plan §8\textquotesingle s correction).
\item
  \textbf{SHA-256} single block-compress, 64 rounds at 1/cycle →
  \textbf{66 cycles/\allowbreak{}block}. Padding, length and multi-block
  chaining are software\textquotesingle s job per FIPS 180-4;
  \texttt{CTRL[4]\ SHA\_\allowbreak{}CONT} chains from the RO
  \code{DIGEST} words.
\item
  \textbf{Keys are not bus-readable} (shadow register outside the bank;
  \code{KEY0-3} read 0 forever), but any master can substitute its own
  key and use the block as an oracle --- see the security notes in
  \src{docs/design/MEMORY_MAP.md}. \textbf{Not production cryptography}:
  no DPA/\allowbreak{}side-channel resistance, no fault hardening, not
  certified.
\end{itemize}

\warn{} \textbf{The throughput table below describes the DEFERRED
AXI4-master version, not what Phase 6b lands.} An APB single-beat
register feed is far slower. Do not quote these figures against the
delivered RTL. (Bead \code{f7vs.1}.)

Note the delivered AES is \textbf{11 cycles/\allowbreak{}block}, not the
100 this table assumed, so the per-block latency is \textasciitilde9×
better than projected --- but the APB register feed, not the round
count, is what bounds throughput, and nothing in this repo measures it.

{
\begin{longtable}[]{@{}
  >{\raggedright\arraybackslash}p{\dimexpr(\linewidth-4\tabcolsep)*245/1000\relax}
  >{\raggedright\arraybackslash}p{\dimexpr(\linewidth-4\tabcolsep)*187/1000\relax}
  >{\raggedright\arraybackslash}p{\dimexpr(\linewidth-4\tabcolsep)*548/1000\relax}@{}}
\toprule\noalign{}
\begin{minipage}[b]{\linewidth}\raggedright
Node
\end{minipage} & \begin{minipage}[b]{\linewidth}\raggedright
AES throughput
\end{minipage} & \begin{minipage}[b]{\linewidth}\raggedright
Notes
\end{minipage} \\
\midrule\noalign{}
\endhead
\bottomrule\noalign{}
\endlastfoot
Sky130 (75 MHz) & \textasciitilde75 MB/s & assumed 100
cycles/\allowbreak{}block; delivered RTL is 11 \\
FreePDK45 (400 MHz) & \textasciitilde400 MB/s & Same pipeline, 5×
faster \\
ASAP7 (1 GHz) & \textasciitilde1 GB/s & High-throughput secure
element \\
\end{longtable}
}

\end{record}

\paragraph{6c. NPU --- Minimal INT8 Inference Engine}\label{cr:status:roadmap:6c-npu-minimal-int8-inference-engine}

\begin{record}

\ok{} \textbf{LANDED 2026-10-04 (bead \bead{f7vs.11})} ---
\src{rtl/npu/npu_top.sv} + \code{npu_mac_array.sv} +
\code{npu_weight_mem.sv}, APB4 slave index 13 at
\code{0x2001_0000-0FFF}, IRQ bit 11. \code{test_npu} 39/39,
\code{test_soc_npu} 1/1, mutation campaign 65 mutants: 55 killed, 10
provably equivalent, 0 gaps. Phase 6 is complete.

\begin{itemize}
\tightlist
\item
  \textbf{4×4 INT8 weight-stationary array} (16 MACs), INT8 × INT8 →
  INT32 accumulate, plus a shared requantizer (×16-bit scale, arithmetic
  floor shift, INT8 saturation, optional ReLU, one lane per cycle). PIO
  over APB4 only --- weights via \code{WADDR}/\code{WDATA}, activations
  via a depth-4 \code{AIN} FIFO, results via a depth-4 \code{AOUT} FIFO.
\item
  \textbf{4 KB weight SRAM, not 16 KB} --- exactly one
  \code{sky130_sram_4kbyte_1rw1r_32x1024_8}, a macro already in the
  Sky130 SoC flow. Every \code{ifdef} arm is a hard macro (4×
  \code{sram_1rw_256x32_*} on ASAP7/\allowbreak{}FreePDK45), so the
  weight memory can never be synthesised as the 32 768 flops that OOM
  this host. 16 KB stays a parameter bump plus three more macros.
\item
  \textbf{No new runtime-indexed mux} (the \code{ma7} Synlig
  \code{OPT_MUXTREE} defect class): proven on the coarse-grain netlist,
  zero \texttt{\$shiftx}/\texttt{\$pmux}/\texttt{\$mem*} in the three
  \src{rtl/npu} modules.
\item
  \textbf{Gate A:} 203 179 µm² stdcell (≈20 \% of the
  SoC\textquotesingle s stdcell area) + the 472 439 µm² macro,
  \textbf{+1.84 ns} setup at ss on 25 ns. Thin: the limiter is the MAC
  array\textquotesingle s single-cycle multiply + adder tree +
  accumulate, and the pre-identified Gate B fallback (registering the
  column sums) was not needed: the Gate B NPU MAC path measured +3.56 ns
  at max\_ss in run 1 (\code{f7vs.15} notes).
\item
  Target: keyword spotting / gesture detection class workloads ---
  \textbf{demo firmware is a non-goal}.
\end{itemize}

{
\begin{longtable}[]{@{}
  >{\raggedright\arraybackslash}p{\dimexpr(\linewidth-6\tabcolsep)*291/1000\relax}
  >{\raggedright\arraybackslash}p{\dimexpr(\linewidth-6\tabcolsep)*140/1000\relax}
  >{\raggedright\arraybackslash}p{\dimexpr(\linewidth-6\tabcolsep)*388/1000\relax}
  >{\raggedright\arraybackslash}p{\dimexpr(\linewidth-6\tabcolsep)*162/1000\relax}@{}}
\toprule\noalign{}
\begin{minipage}[b]{\linewidth}\raggedright
Node
\end{minipage} & \begin{minipage}[b]{\linewidth}\raggedright
NPU clock
\end{minipage} & \begin{minipage}[b]{\linewidth}\raggedright
Peak throughput
\end{minipage} & \begin{minipage}[b]{\linewidth}\raggedright
INT4/\allowbreak{}sparsity
\end{minipage} \\
\midrule\noalign{}
\endhead
\bottomrule\noalign{}
\endlastfoot
Sky130 (40 MHz, as built) & 40 MHz & 1.28 GOPS (16 MAC × 2 op × 40 MHz)
& \bad{} Too slow \\
FreePDK45 (300--400 MHz) & 300--400 MHz & \textasciitilde20 GOPS & \ok{}
Feasible \\
ASAP7 (800 MHz--1 GHz) & 800 MHz+ & \textasciitilde50 GOPS & \ok{}
Preferred \\
\end{longtable}
}

INT8 MAC critical path: \textasciitilde10--13 logic levels. Full NPU
(INT4/\allowbreak{}tiling/\allowbreak{}sparsity) requires FreePDK45 or
ASAP7.

\end{record}

\paragraph{Technology Node Cross-Comparison (PD Results + Estimates)}\label{cr:status:roadmap:technology-node-cross-comparison-pd-results-estimates}

\begin{record}

{
\begin{longtable}[]{@{}
  >{\raggedright\arraybackslash}p{\dimexpr(\linewidth-6\tabcolsep)*188/1000\relax}
  >{\raggedright\arraybackslash}p{\dimexpr(\linewidth-6\tabcolsep)*188/1000\relax}
  >{\raggedright\arraybackslash}p{\dimexpr(\linewidth-6\tabcolsep)*344/1000\relax}
  >{\raggedright\arraybackslash}p{\dimexpr(\linewidth-6\tabcolsep)*261/1000\relax}@{}}
\toprule\noalign{}
\begin{minipage}[b]{\linewidth}\raggedright
Phase
\end{minipage} & \begin{minipage}[b]{\linewidth}\raggedright
Sky130 (130nm)
\end{minipage} & \begin{minipage}[b]{\linewidth}\raggedright
FreePDK45 (45nm)
\end{minipage} & \begin{minipage}[b]{\linewidth}\raggedright
ASAP7 (7nm)
\end{minipage} \\
\midrule\noalign{}
\endhead
\bottomrule\noalign{}
\endlastfoot
Ph2 CPU & 75 MHz \ok{} & 400 MHz \ok{} (run 2026-03-23) & 1 GHz
(planned) \\
Ph3 CPU+Cache & 75 MHz \ok{} & \textbf{400 MHz \ok{} WNS=0} (run
2026-03-23) & 500 MHz--1 GHz (planned) \\
Ph4 GPU-lite & 50--75 MHz (est.) & 300 MHz (est.) & 500 MHz (est.) \\
Ph5 SoC & 50--75 MHz (est.) & 250 MHz (est.) & 400 MHz (est.) \\
Ph6+ peripherals & 75 MHz & 400 MHz & 1 GHz \\
Ph6+ AES/\allowbreak{}SHA-256 & 75 MHz & 400 MHz & 1 GHz \\
Ph6+ NPU (min.) & 50 MHz & 300 MHz & 800 MHz \\
Ph6+ NPU (full) & \bad{} & 300 MHz & 800 MHz \\
\end{longtable}
}

\textbf{FreePDK45 Phase 3 PD results (run 2026-03-23)}:

\begin{itemize}
\tightlist
\item
  WNS = 0 ns, TNS = 0 ns @ 400 MHz (timing closed)
\item
  Die: 760 × 480 µm = 364,800 µm² (15.8× smaller than Sky130)
\item
  Power: 22.59 mW total (vs 59.86 mW Sky130 --- 62\% lower despite 5.3×
  faster clock)
\item
  Cells: 26,038 standard cells, SRAM macros: 164,286 µm²
\item
  PDN: SRAM macros not connected to grid
  (PDN\_\allowbreak{}CONNECT\_\allowbreak{}MACROS\_\allowbreak{}TO\_\allowbreak{}GRID=false)
  --- known limitation for predictive PDK
\end{itemize}

\textbf{ASAP7 sign-off} \ok{} (Run 43, 2026-05-20): 1418 MHz / 27.27 mW
/ 3 844 µm² stdcell, 0 DRC / 0 antenna / 0 setup-hold violations. RVT TT
@ 0.7 V / 25 °C. SRAM via FF-array stub
(\code{sram_1rw_256x32_asap7_stub.v} + Liberty/\allowbreak{}LEF). Config
at \src{pnr/asap7/}. Full campaign history:
\cref{cr:run:asap7}. \warn{} Timing/\allowbreak{}power
unvalidated post-bead-\code{8f3} (STA-zero-wire, 2026-09-16); Run 43
artifacts wiped, not re-checkable --- bead \bead{0p6}.

\textbf{Sky130 real DRC/LVS sign-off} (GH epic \#102) --- the only PDK
in this project with genuine physical verification; ASAP7 skips
Magic/\allowbreak{}KLayout/\allowbreak{}Netgen entirely, and FreePDK45
has them permanently blocked. Stages 1--2 \ok{} complete 2026-07-31,
Stages 3--4 \hold{} host-gated.

\begin{itemize}
\tightlist
\item
  \textbf{Stage 1 --- CPU standalone} (\ghref{103}): KLayout DRC 0, Netgen
  LVS MATCH, setup +0.366 ns @ 75 MHz, native Sky130 macro views. Magic
  DRC 27.7 M waived (100 \% inside SRAM macro footprints, PDK artifact;
  bead \bead{45a}).
\item
  \textbf{Stage 2 --- CPU macro + peripherals SoC, no GPU} (\ghref{104}):
  LVS PASSED, routing DRC 0, PDN 0, hold clean, KLayout DRC 4 (waived),
  40 MHz, nom\_tt 37.63 mW, 226,952 stdcells, die 6700 × 3100 µm. Two
  documented residual gates, neither open work: setup fails at ss (bead
  \bead{ujv}, \textbf{deferred host-blocked}) and antenna 140 pin / 120
  net (bead \bead{58q}, \textbf{closed} as a quantified accepted
  trade-off against hold).
\item
  \warn{} \textbf{Stage-2 timing is a typical-corner (nom\_tt)
  sign-off}, not a validated multi-corner one --- the SRAM macro is
  characterized at TT only, so the nine reported corners are nine labels
  backed by a single macro model. The physical gates
  (LVS/\allowbreak{}DRC/\allowbreak{}PDN) carry no such caveat.
  Per-corner SRAM SPICE characterization measured at
  \textasciitilde80--95 h against a host that reboots every 2--8 h →
  bead \bead{o1i}, deferred host-blocked.
\item
  \textbf{Stages 3--4 --- GPU} (\ghref{105}/\allowbreak{}\#106): not
  attempted; need ≥32 GB RAM + \textasciitilde500 GB scratch vs the
  current \textasciitilde15 GB host. GPU stays ASAP7-indicative-only.
\item
  Full record: \src{docs/SKY130_REAL_DRC_LVS_EVALUATION.md} §7.
\end{itemize}

\end{record}

\subsubsection{Post-Phase-6 tracks: DFT and FPGA (not numbered phases)}\label{cr:status:roadmap:post-phase-6-tracks-dft-and-fpga-not-numbered-phases}

\begin{record}

\begin{itemize}
\tightlist
\item
  \textbf{DFT: scan, MBIST and JTAG for the Sky130 SoC} (epic
  \code{j41m}): Stage 0 architecture \ok{}
  (\src{docs/design/DFT_ARCHITECTURE.md}); Stage 1a fabric RTL merged
  (\prref{258}, bead \bead{j41m.2} still open); Stage 1b JTAG TAP \ok{} (\prref{270}). Stage 2 scan insertion is not started, and MBIST has no RTL
  yet.
\item
  \textbf{FPGA prototyping} (epic \code{hlrm}): Stage 0 \ok{} without
  hardware (\src{docs/FPGA_STAGE0_EVALUATION.md}, \prref{279}). Board
  recommendation Digilent Nexys Video XC7A200T
  (\src{docs/FPGA_BOARD_SELECTION.md}, \prref{283}). Open-tool ECP5 fmax
  lower bound \textasciitilde48--52 MHz CPU-only. Later stages wait for
  a board.
\end{itemize}

\end{record}

\subsubsection{Phase 7 --- Mixed-Signal PLL Clock Generator}\label{cr:status:roadmap:phase-7-mixed-signal-pll-clock-generator}

\begin{record}

\textbf{Status}:
M-a/\allowbreak{}M-b1/\allowbreak{}M-b2/\allowbreak{}M-c complete
(2026-06-20); M-d = documentation (this update). Full plan →
\src{docs/PHASE7_MIXED_SIGNAL_PLL_PLAN.md} (golden spec).

A dual-PDK charge-pump integer-N PLL (ring VCO) designed end-to-end via
the \textbf{analog-design orchestrator agents}, then integrated into the
SoC as the clock source via a PDK-agnostic real-number model (RNM). This
phase also doubled as a validation run of the installed analog-design
agent suite (real ngspice / magic / netgen all ran; the closed-loop meta
fix loop and adversarial physical verification both demonstrably added
value).

\begin{itemize}
\tightlist
\item
  \textbf{ASAP7 variant} (\src{analog/pll_clkgen/}, indicative): 100 MHz
  → 1.282 GHz, N=13, 0.7 V --- matches the ASAP7 SoC 780 ps clock.
  Electrical is \textbf{indicative only}: open ngspice has no BSIM-CMG
  FinFET models, so transistor sims used a calibrated planar BSIM4
  substitute.
\item
  \textbf{Sky130 variant} (\src{analog/pll_clkgen/sky130/}, \textbf{real
  sign-off}): 10 MHz → 100 MHz, N=10, 1.8 V, real \code{sky130_fd_pr}
  models, real magic layout on \code{sky130A.tech}. Charge-pump block:
  real magic DRC = 0 + netgen LVS MATCH.
  VCO/\allowbreak{}LF/\allowbreak{}BIAS block regen + full-chip DRC/LVS
  are follow-ups.
\item
  \textbf{SoC integration} (\src{rtl/soc/pll/}): \code{clk_i} → PLL ref
  → \code{core_clk}; AXI-Lite config/\allowbreak{}status slave at
  \code{0x2000_7000} (\code{AXIL_N_SLAVES} 6→7). Two RTL bugs
  found+fixed (PLL bootstrap deadlock on gated domain;
  \code{PERIPH_LIMIT} decode hole). Cosim 7/7 (\code{test_pll_lock} 3/3
  + \code{test_pll_regs} 4/4) + regression clean. RNM-mode AXI CDC
  synchroniser deferred.
\end{itemize}

\end{record}

% Moved verbatim from CLAUDE.md by tools/docs/md2tex.py; edit this file, not the original.
\chapter{Project Overview}\label{cr:status:claude-md-overview}

\begin{record}

Multi-phase RV32I RISC-V microprocessor + GPU-lite SoC project.

\textbf{Current Phase}: Phase 5 (SoC Integration) --- \ok{} COMPLETE
(2026-06-24). M1--M12 done; M11 ASAP7 SoC P\&R signed off at \textbf{571
MHz / 62.9 mW \warn{} fabric-only, excludes both hard macros / 520×520
µm / 65.6 \% util / 0 DRC / 0 antenna \warn{} routing claims
unsupported, see below} (run 14, sv2v frontend; PDN \warn{} withdrawn
--- the grid is EMPTY, see below). Power-attribution root cause +
bounds: \cref{cr:run:phase5} §"Power figure caveat" (bead
\bead{ew3}). \textbf{The "0 DRC / 0 antenna" figure is vacuous --- run
14\textquotesingle s detailed routing committed zero wires (bead
\bead{xy6}); timing/\allowbreak{}power/\allowbreak{}area numbers are
unaffected (GRT-estimated parasitics, as always) but every routing-clean
claim is withdrawn.} \cref{cr:run:phase5} §"Routing /
physical-closure caveat". \warn{} \textbf{The 571 MHz / 62.9 mW
timing+power figure itself is now also unvalidated (bead
\bead{8f3}/\code{0p6}, 2026-09-16): it came from
\code{OpenROAD.STAMidPNR-3} run WITHOUT an in-session
\code{global_route}, which a CPU-block proof showed reports
falsely-clean 0/0/0 violations on \textasciitilde94\% zero-wire
parasitics. Run 14\textquotesingle s artifacts are wiped (bead
\bead{alq}) and cannot be re-measured with the now-fixed
\code{STA_POSTGRT_INSESSION_GRT=1}; the honest current-RTL replacement
is 258.5 mW / setup WNS −2084 ps / 41,864 violators
(\cref{cr:run:phase5} §"Power figure caveat", bead
\bead{w3a}/\code{86a}/\code{0p6}) --- a different, unfinished design
point, not a re-run of run 14.} \ok{} \textbf{UPDATE 2026-09-18 (bead
\bead{0d0}): the SoC has since been re-measured three times on the fixed
flow; the current honest current-RTL point is setup WNS −727.37 ps / TNS
−60,665 ps / 781 violators, hold CLEAN, 283.6 mW total / 89.6 mW
fabric-only (\code{RUN_2026-09-17_20-41-11}, pinned OpenROAD 26Q2,
annotation-gated, zero real signal drivers unannotated) --- this
supersedes the −2084 ps figure above as the current basis. Timing is
still NOT closed at 1750 ps. Run 14 itself remains unvalidated and
unrecoverable.} See \cref{cr:run:phase5}.

\textbf{Status}: Phase 4 complete (2026-05-27, all GPU tests green;
ASAP7 GPU sign-off 571 MHz / 262 mW \warn{} STA-zero-wire, unvalidated
post-\code{8f3}, re-run deferred --- see
\cref{cr:run:asap7-gpu}); Phase 3 complete (2026-05-21,
139/139 total, ASAP7 1418 MHz sign-off \warn{} same STA caveat, bead
\bead{8f3}/\code{0p6}); Phase 2 complete (2026-03-08, 75 MHz on Sky130,
real post-route SPEF --- unaffected).

\bad{} \textbf{ASAP7 CPU/GPU macro netlists: PROVEN frontend miscompile
(bead \bead{u99} → \code{ma7}, 2026-09-19).} Both macros are built with
\code{USE_SYNLIG:true}. A differential (same
testbench/\allowbreak{}ROM/\allowbreak{}sequencing, Synlig gate netlist
vs RTL sim) showed \code{rv32i_cpu_top}\textquotesingle s netlist
returns \textbf{0x30} for an APB debug read of \code{DBG_GPR[4]} where
RTL returns \textbf{0x0c} --- harness error, timing races and
debug-semantics explanations all excluded. Mechanism: \code{OPT_MUXTREE}
mis-elaborating a \textbf{runtime-indexed mux}, the same class proven in
beads \code{gcd} and \code{b0t}. \src{rtl/cpu/core/rv32i_regfile.sv}
uses that idiom on \textbf{all three} read ports, so the two feeding
ID/\allowbreak{}execute are \textbf{UNVERIFIED, not proven clean}
(\code{u99} checked insn/PC retire order, not register values);
\code{rtl/gpu/warp_scheduler.sv:152} is the same class, untested. Remedy
(switch both configs to \code{USE_SYNLIG:false} + sv2v → re-harden both
macros → SoC re-run, expect PPA to worsen as \code{valu_hls} did) is
\textbf{not applied} and needs a human decision: bead \bead{ma7} (P1).
Treat every ASAP7 CPU/GPU macro PPA figure below as resting on a netlist
from a frontend with a proven miscompile in this repo. \textbf{The
Sky130 SoC fabric is unaffected by the frontend question (sv2v), but the
Sky130 CPU hard macro is NOT: it is configured \code{USE_SYNLIG:true}
and its committed netlist
\src{pnr/sky130/cpu/macro/rv32i_cpu_top.nl.v.gz} is itself functionally
WRONG (bead \bead{dud4}, 2026-10-10,
\cref{cr:defect:sky130-cpu-synlig-dud4})} --- the u99 program leaves
\code{x4=0x30} (RTL \code{0x0c}), the ma7
\texttt{BEQ\ x2,\allowbreak{}x3} is evaluated NOT-TAKEN, confirmed in
two simulators; the Sky130 SoC consumes this macro. Root cause
(measured, simulation-confirmed): Synlig elaborates
\code{.if_id_rs1_addr(if_id_reg.instruction[19:15])} /
\code{.if_id_rs2_addr(...[24:20])} in \code{rv32i_core.sv}
\code{u_hazard} as \texttt{5'x} (log:
\texttt{Range\ select\ [639:\allowbreak{}608]\ out\ of\ bounds\ .\allowbreak{}.\allowbreak{}.\allowbreak{}\ Setting\ all\ 32\ result\ bits\ to\ undef}),
so forwarding selects are wrong and stored values are wrong; an sv2v
synthesis of the same RTL is cycle-exact identical to RTL on all test
programs (31 regs x 3 ports). It is NOT the regfile read mux. INFERRED,
untested: ASAP7\textquotesingle s CPU (same two lines, Synlig-built
until \code{ma7}\textquotesingle s config switch) has the same defect,
so the u99/ma7 "OPT\_\allowbreak{}MUXTREE regfile read mux" diagnosis is
in doubt. \bad{} \textbf{UPDATE 2026-09-19 (bead \bead{ma7} step 1): the
corruption is PROVEN to reach the execute datapath, not just the debug
port.} A branch-dependent differential (same \code{tb_cpu_macro_check.v}
harness/\allowbreak{}method as \code{u99}, new program whose control
flow depends on computed register values so an ID-stage regfile read
corruption becomes observable via \code{commit_pc_o}, not the known-bad
debug port) shows the two arms retire IDENTICALLY through a
\texttt{BEQ\ x2,\allowbreak{}x3} instruction\textquotesingle s own fetch
(x2=x3=9, must be TAKEN) and then diverge exactly at that
instruction\textquotesingle s commit: the frontend-independent RTL
reference (Verilator) takes the branch correctly (matches the Python
golden model bit-exact); the same Synlig gate netlist u99 used evaluates
it as NOT-TAKEN and parks at a dead marker instead. Exact coverage:
proven corrupt for register indices x2/x3 (value 9) as read by the
branch comparator\textquotesingle s rs1/rs2 ports feeding EX; not proven
for other indices, instruction types, or access patterns. Full method,
evidence paths, and the read-port cone-sharing analysis: bead \bead{ma7}
notes and \src{docs/CPP_TO_RTL_HLS_EVALUATION.md}. \bad{} \textbf{UPDATE
2026-09-19 (bead \bead{ma7} remedy, PARTIALLY APPLIED --- CPU re-harden
BLOCKED, GPU still deferred).} User-approved remedy applied to configs
only so far: \src{pnr/asap7/cpu/config.json} + \code{config_3014.json}
and \src{pnr/asap7/gpu/config.json} + \code{config_3014.json} are now
\code{USE_SYNLIG:false} with a new sv2v-preprocessing Makefile pipeline
(\code{asap7-cpu-sv2v}, \code{asap7-gpu-sv2v}), mirroring the
already-applied \code{valu_hls} fix. \textbf{The CPU re-harden PD run
has NOT completed cleanly}: 3 attempts (lint fix needed and applied; a
known unbounded-resizer-loop defect, bead bpp, needed and got
\code{PL_RESIZER_TIMING_MAX_PASSES=2}; the 3rd attempt,
\code{RUN_2026-09-19_15-42-11}, reached \code{OpenROAD.DetailedRouting}
cleanly through CTS/\allowbreak{}GRT/\allowbreak{}resizer but then hit a
\textbf{new, unresolved routing-congestion regression} --- 108716 DRC
violations, \textasciitilde53x this block\textquotesingle s own
best-ever prior routed baseline of 2045, concentrated as \code{DRT-0255}
failures on SRAM macro pins --- that leaves \code{OpenROAD.RCX}
extracting zero parasitics and \code{OpenROAD.STAPostPNR} failing
outright). Root cause NOT pinned down (ruled out: wrong OpenROAD binary,
simple netlist-size growth); tracked as bead \bead{lxv}.
\textbf{\code{pnr/asap7/soc/macro/rv32i_cpu_top.*} are UNCHANGED and
remain the old, Synlig-built, proven-corrupt views} --- no macro view
was regenerated from an incomplete run. GPU re-harden itself remains
fully deferred (host-gated, unchanged from the note above);
\code{pnr/asap7/soc/macro/gpu_top.*} also remain Synlig-built and
suspect. See bead \bead{ma7} notes for the full attempt history and the
methodology-deviation discussion (disabling SPEF/\allowbreak{}MCSTA to
force flow completion was considered and explicitly rejected as not
matching this block\textquotesingle s own established reference
methodology). \bad{} \textbf{UPDATE 2026-09-19 (bead \bead{lxv}, root
cause narrowed further --- CPU re-harden STILL BLOCKED).} Five
single-variable experiments across two sessions have now excluded:
\code{PL_MACRO_HALO}, \code{GRT_LAYER_ADJUSTMENTS}, a structural netlist
difference at the failing \code{clk0}/\code{csb0}/\code{din0} pins
(byte-identical gate types/\allowbreak{}fanin between the Synlig and
sv2v builds), a placement-clearance difference around the SRAM macros
(\textless1\% delta), a genuine global-routing capacity problem (the
failing run\textquotesingle s own GRT reports 0/0/0 overflow at 11.86\%
usage), a relaxed \code{CLOCK_PERIOD} (780→1000 ps: DRC 109687 vs.
108716 baseline, despite −9.8\% fewer resizer-inserted buffers), and a
raised \code{PL_RESIZER_TIMING_MAX_PASSES} (2→4: DRC 109562 vs. 108716,
despite +2.1\% more buffers). DRC count is invariant to resizer-driven
cell-count changes in either direction. The only unexcluded explanation
is the SRAM hard-macro\textquotesingle s own pin-access geometry
(\code{sram_1rw_256x32_asap7.lef}: all \textasciitilde74 signal pins
\texttt{SHAPE\ ABUTMENT} on one edge) --- pre-existing in both
frontends, consistent with \code{e69}\textquotesingle s own non-zero
2045-violation baseline on identical geometry. A real fix likely
requires a hard-macro LEF/\allowbreak{}pin-layout change, which
\textbf{needs explicit human approval} before any macro view is touched;
none was attempted. Full trail: bead \bead{lxv},
\cref{cr:run:asap7}. \hold{} \textbf{DECISION 2026-09-20
(user): SHELVED until a ≥32 GB host is available} --- beads \code{lxv}
and \code{ma7} are now formally blocked on \code{2kn}. The user declined
both alternatives (authorising a hard-macro LEF edit, or re-baselining
against the pre-existing 2 045-violation routing baseline) in favour of
deferring. \textbf{Consequence to carry forward: the ASAP7 CPU and GPU
macro views remain the Synlig-built netlists proven to mis-execute a
branch, so every ASAP7 CPU/\allowbreak{}GPU/\allowbreak{}SoC PPA figure
in this repo rests on a netlist with a proven functional miscompile.}
The correctness fix \emph{is} in the configs (\code{USE_SYNLIG:false} +
sv2v); what is missing is a completed PD run to regenerate the views. Do
not attempt further config-level knob sweeps on this host --- the
exclusion list is exhaustive for config variables (10 hypotheses refuted
across 5 single-variable experiments). On obtaining the host, resume by
re-running the sv2v CPU re-harden \textbf{with} the real power grid
before touching macro geometry: a PDN-bearing route has never been tried
and may change the DRT picture on its own. \bad{} \textbf{UPDATE
2026-10-10 (bead \bead{ainf}, \cref{cr:defect:synlig-struct-select-ainf}):
"Synlig + \prref{259}" is supported for the CPU struct-member selects but
NOT for the GPU.} Synlig has three silent struct-select defects (D1:
\code{s.m[..]} on a struct type with exactly two packed-range members;
D2: LHS part-select of a member; D3: element select of a
multi-dimensional packed-array member). \src{rtl/cpu}: the 15 remaining
member selects are all on structs with R\textgreater=7 ranged members,
there is no LHS member select and no multi-dim array member, and
Synlig-vs-sv2v miter+SAT proves \code{rv32i_pipeline_ex/_mem/_wb/_ex1c}
EQUIV (not proven: whole core, regfile, caches). \src{rtl/gpu}:
\code{gpu_compute_unit} \code{id_ex_q.rs1_data[l]}/\code{rs2_data[l]} is
read as ONE bit by Synlig (D3, measured NOT-EQUIV), so
VLD/\allowbreak{}VST/\allowbreak{}VLDS/\allowbreak{}VSTS addresses and
data are wrong in any Synlig-built GPU netlist, including the committed
\code{pnr/asap7/soc/macro/gpu_top.*}. The CPU conclusion is a property
of today\textquotesingle s source (\code{if_id_reg_t} is a latent R=2
type; \code{check_struct_member_select.py} guards new selects). \ok{}
\textbf{UPDATE 2026-10-10 (beads \code{ma7}/\code{dud4},
\cref{cr:defect:asap7-cpu-synlig-ma7}): the ASAP7 miscompile is the same
defect as Sky130\textquotesingle s, NOT a regfile read-mux /
\code{OPT_MUXTREE} fault --- this SUPERSEDES the
"OPT\_\allowbreak{}MUXTREE runtime-indexed mux" mechanism and the "two
read ports feeding ID/\allowbreak{}execute are UNVERIFIED" statement
above (the dated entries above are kept as history).} Gate-vs-RTL
differential on four ASAP7 \code{rv32i_cpu_top} netlists, two simulators
(yosys \code{sim} with the original u99/ma7 harness, and Verilator with
Liberty-derived cell models, cycle-exact on every output port): the
\textbf{committed macro} and a fresh \textbf{Synlig} synthesis of main
RTL (arm S) both give \code{DBG_GPR[4]}=\code{0x30} (RTL \code{0x0c})
and \texttt{BEQ\ x2,\allowbreak{}x3} NOT taken; \textbf{Synlig + the two
hazard-unit part-selects routed through a named wire (arm W) and the
sv2v netlist (arm V) are IDENTICAL to RTL on all 6 programs and pass the
full 31-register x (rs1, rs2, both, debug) port matrix (124/124 cells)},
while S and the committed macro fail 96/124. Mechanism (MEASURED
isolation, INFERRED internals): Synlig elaborates
\code{rv32i_core.sv:436-437}
(\code{.if_id_rs1_addr(if_id_reg.instruction[19:15])},
\code{...[24:20]}) as \texttt{5'x}, so the forwarding selects ignore the
consumer\textquotesingle s source register and every ALU operand becomes
the previous instruction\textquotesingle s result (the failing values
reproduce exactly: \code{x2=5+7}, \code{x3=12+12},
\code{x4=24+24=0x30}); the wrong stored values then make the BEQ
operands unequal. The regfile read path is correct in W and V. Coverage
limits: 6 programs, not exhaustive; W/V correctness is by simulation,
not formal equivalence; the committed macro is judged on its own netlist
plus RTL \code{999e44a}. \textbf{Consequence for \code{lxv}/\code{ma7}
(an option, not a decision):} the blocked sv2v CPU re-harden is not the
only route to a correct ASAP7 CPU macro --- a Synlig build with the RTL
fix (\prref{259}: the two part-selects through a full-width wire) gave a
functionally correct netlist (arm W) and keeps the Synlig
netlist/\allowbreak{}P\&R baseline whose routing was healthier (2 045
DRC violations in \code{e69} vs 108 716 in the sv2v attempts), but it
still needs a full re-harden to regenerate the macro views, and the GPU
macro (\code{warp_scheduler.sv:152}, other \code{struct.member[a:b]}
port connections) was NOT checked for the same class. \bad{}
\textbf{UPDATE 2026-10-11 (bead \bead{dud4}, Sky130 CPU macro
re-hardened on sv2v and proven vs RTL; supersedes the earlier "Sky130
and the SoC fabric are unaffected (both use the sv2v frontend)" wording,
which was FALSE for the CPU hard macro).} The July Sky130 CPU macro was
Synlig-built (\code{USE_SYNLIG:true}) and mis-executed (u99
\code{x4=0x30}, BEQ not taken). Root cause: Synlig elaborated the two
hazard-unit port connections
\code{.if_id_rs1_addr(if_id_reg.instruction[19:15])}/\code{rs2} as
\texttt{5'x} (struct-member part-select), deleting the \code{fwd_b_*}
forwarding-select flops --- \textbf{not} \code{OPT_MUXTREE} on the
regfile read mux (that \code{u99}/\code{ma7} attribution is in doubt,
ASAP7 shows the same two lines). Fixed by the RTL rewrite (\prref{259})
\textbf{plus} this re-harden: \src{pnr/sky130/cpu/config.json} is now
\code{USE_SYNLIG:false} on the new
\texttt{make\ -\allowbreak{}C\ pnr\ sky130-\allowbreak{}cpu-\allowbreak{}sv2v}
output; run \code{RUN_2026-10-11_06-00-12} (75 MHz, config otherwise
unchanged); the committed post-fill netlist matches RTL on all 6
gate-vs-RTL programs and 124/124 port-matrix cells, census 4868 flops
with \code{fwd_b_*} present, negative control detected
(\src{tools/verif/gls/sky130/run_netlist_proof.sh}). New views are in
\src{pnr/sky130/cpu/macro/} (all 403 LEF pins at identical positions).
CPU standalone: LVS match, KLayout DRC 0, routing DRC 0, Magic
27,733,913 (same count as before, 99.9869 \% inside SRAM footprints,
remaining 3,626 = known \code{nwell.4} abstract artifact); setup clean
at all tt/ff corners (+3.44 ns nom\_tt) but \textbf{ss reg-to-reg setup
FAILS (max\_ss −3.70 ns, 996 violators; icache tag-SRAM half-cycle paths
--- reported, not tuned)}; hold fails only on I/O ports. Sky130 SoC
re-run on the new macro (\code{RUN_2026-10-11_08-39-41}, rccal flow):
LVS 0, routing DRC 0, setup and hold clean at all 9 corners (worst +2.48
/ +0.287 ns); \code{CVT.MaxCapViolations} FAILS (3 flat nets, ss
corners) and the \code{axi_rdata_i[12]} waiver is now stale --- owner
decisions open (see bead \bead{dud4}, \cref{cr:defect:sky130-slew-cap-e45j}
§8). \textbf{This is a NEW baseline, not a single-variable comparison}
(frontend + RTL since July + registered-APB-output macro). Every earlier
Sky130 CPU/SoC result wrapped the functionally wrong macro.
\textbf{ASAP7 CPU/GPU macros are still the corrupt Synlig views
(shelved, blocked on bead \bead{2kn}/\code{lxv}).} Note: the gc0y
synthesis gate\textquotesingle s structural \code{x-driven-net} check
false-positives on 509 sv2v function-inlining temporaries when run on
the SoC (it was validated on the CPU only); the SoC run used
\code{CVT_SYNTH_GATE=warn} (shell export, not in any tracked file) ---
owner decision open.

\textbf{Current status (2026-10-15)}: Phase 6 IP expansion \ok{}
complete (6a--6c RTL; batched Sky130 Gate B harden \ok{} 2026-10-05,
items 13--15 below). Two tracks are in progress: \textbf{DFT} (epic
\code{j41m}: scan + JTAG + MBIST for the Sky130 SoC; Stage 1b JTAG TAP
merged, Stage 1a fabric RTL merged with bead \bead{j41m.2} still open,
item 16) and \textbf{FPGA prototyping} (epic \code{hlrm}: Stage 0 done,
no hardware, item 17).

\textbf{Target}: Build a complete SoC with CPU, GPU, caches, and
peripherals through incremental phases.

Per-phase completion records → \cref{cr:status:phase-history} and
\cref{cr:status:phase-status}.

\end{record}

\section{Project Phases}\label{cr:status:claude-md-overview:project-phases}

\begin{record}

See \cref{cr:status:roadmap} for the complete phase plan and
\cref{cr:status:phase-status} for current status.

{\footnotesize\setlength{\tabcolsep}{3pt}

{
\begin{longtable}[]{@{}
  >{\raggedright\arraybackslash}p{\dimexpr(\linewidth-4\tabcolsep)*159/1000\relax}
  >{\raggedright\arraybackslash}p{\dimexpr(\linewidth-4\tabcolsep)*410/1000\relax}
  >{\raggedright\arraybackslash}p{\dimexpr(\linewidth-4\tabcolsep)*410/1000\relax}@{}}
\toprule\noalign{}
\begin{minipage}[b]{\linewidth}\raggedright
Phase
\end{minipage} & \begin{minipage}[b]{\linewidth}\raggedright
Scope
\end{minipage} & \begin{minipage}[b]{\linewidth}\raggedright
Status
\end{minipage} \\
\midrule\noalign{}
\endhead
\bottomrule\noalign{}
\endlastfoot
0 --- Foundations & Specs + Python reference models (no RTL) & \ok{}
2026-01-18 \\
1 --- Minimal RV32I core & Single-cycle; AXI4-Lite master; APB3 debug &
\ok{} 2026-02-13 \\
2 --- Pipelined CPU & 5-stage in-order; RV32I+Zicsr; M-mode interrupts;
hazard/\allowbreak{}forwarding & \ok{} 2026-03-08 (75 MHz Sky130) \\
3 --- Memory system & L1 I\$ + D\$ (4 KB each, direct-mapped,
write-back); FENCE.I & \ok{} 2026-05-21 (ASAP7 1418 MHz) \\
4 --- GPU-Lite SIMT & 8-lane warps, single CU, round-robin, 1-level
divergence, 16 KB shared mem & \ok{} 2026-05-27 (ASAP7 571 MHz) \\
5 --- SoC integration & CPU + GPU + DMA + AXI4 crossbar +
UART/\allowbreak{}SPI/\allowbreak{}timer/\allowbreak{}IRQ + behavioral
SRAM & \ok{} 2026-06-24 (ASAP7 SoC 571 MHz / 62.9 mW \warn{} fabric-only
sign-off) \\
6+ --- IP expansion &
GPIO/\allowbreak{}I2C/\allowbreak{}PWM/\allowbreak{}WDT/\allowbreak{}TRNG/\allowbreak{}AES-SHA
peripherals; INT8 NPU; tech-node exploration & \ok{} \textbf{Phase 6
complete 2026-10-05} --- \textbf{6a \ok{}} (GPIO \code{ckc}/\code{00ef}
+ Sky130-hardened, PWM \code{f7vs.6}, WDT \code{f7vs.7}, TRNG
\code{f7vs.8}, I2C \code{f7vs.9}); \textbf{6b CRYPTO \ok{} done}
(AES-128 + SHA-256, bead \bead{f7vs.10}); \textbf{6c NPU \ok{} done}
(INT8 4×4, bead \bead{f7vs.11}) --- Phase 6 RTL complete; Gate B batched
Sky130 harden \ok{} 2026-10-05 (\code{f7vs.15}, \prref{214}); golden spec
\src{docs/PHASE6_IP_EXPANSION_PLAN.md}, epic \code{f7vs} (closed) \\
7 --- Mixed-Signal PLL & Dual-PDK charge-pump PLL (ASAP7 indicative +
Sky130 real DRC/LVS) via analog-design agents; AMS RNM integrated as SoC
clock source & done (M-a..M-c) \\
\end{longtable}
}

}

\end{record}


% Moved verbatim from CLAUDE.md by tools/docs/md2tex.py; edit this file, not the original.
\chapter{Current Workflow (Phase 5 --- SoC Integration)}\label{cr:status:phase5-workflow}

\begin{record}

Legend: \ok{} done · \wip{} in progress · \hold{} not started. Milestone
status (M1--M12) tracked in \src{docs/PHASE5_SOC_INTEGRATION_PLAN.md}.

\begin{enumerate}
\def\labelenumi{\arabic{enumi}.}
\item
  \ok{} \textbf{AXI4 interconnect} (M1/M3): \code{axi4_crossbar.sv},
  \code{axi_lite_interconnect.sv}, \code{axi_lite_register_bank.sv};
  Phase 3 refill FSMs upgraded to AXI4 burst (M2)
\item
  \ok{} \textbf{Peripherals + DMA} (M4/M5): UART, SPI, timer, interrupt
  controller, DMA engine in \src{rtl/periph/}
\item
  \ok{} \textbf{SRAM controller} (M6): behavioral AXI4-slave in
  \src{rtl/soc/sram_controller.sv}
\item
  \ok{} \textbf{SoC top integration} (M8): CPU + GPU + DMA + peripherals
  wired in \src{rtl/soc/soc_top.sv} (+ \code{boot_rom.sv})
\item
  \ok{} \textbf{Performance counters} (M7): CSR-mapped + AXI-Lite GPU
  stats (CPU re-sign-off 1282 MHz)
\item
  \ok{} \textbf{SoC verification} (M9): boot 100/100; DMA+UART+SPI
  loopback; SW coherency (D\$ flush→GPU→D\$ inval); CPU-GPU IRQ-driven
  integration; DUT-side boot SRAM check. Found+fixed 2 RTL bugs (D-cache
  MMIO caching \code{go9}; axi4\_\allowbreak{}crossbar AR/AW
  handshake+arbitration \code{7fs}). \code{soc_all} 73/73 at M9 (142/142
  as of 2026-08-01, after the PMU epic + \ghref{91}
  async\_\allowbreak{}axi\_\allowbreak{}fifo).
\item
  \ok{} \textbf{Full regression}: CPU rollup + GPU + cache green;
  \textbf{1M+ cycle SoC stress (1,079,867 cyc, 0 fail)} + CPU/GPU
  benchmarks (M9 \code{pgf})
\item
  \ok{} \textbf{L2 cache evaluation} (M10): benchmarked → \textbf{NO-GO}
  (no L2, keep direct-mapped L1; 2-way also evaluated + rejected).
  \src{docs/M10_L2_DECISION_ANALYSIS.md}
\item
  \ok{} \textbf{Backend flow} (M11): full SoC synth + P\&R + STA via
  \texttt{make\ librelane-\allowbreak{}asap7-\allowbreak{}soc} (sv2v
  frontend --- Synlig can\textquotesingle t synth this SoC;
  \src{pnr/asap7/soc/}, \code{phase5_soc.sdc} single-clock,
  \code{phase5_soc.upf} 4-domain). Signed off \textbf{571 MHz / 62.9 mW
  \warn{} fabric-only, excludes both hard macros / 520×520 µm / 65.6 \%
  util / 0 DRC / 0 antenna \warn{} unsupported, see bead \bead{xy6}}
  (run 14). PDN \warn{} withdrawn --- every ASAP7 run committed ZERO
  power-grid shapes (bead \bead{4l8}); the "benign tap-cell artifact"
  reading was wrong. \warn{} Timing/\allowbreak{}power also unvalidated
  post-\code{8f3} (STA zero-wire, run 14 artifacts wiped) --- see
  \cref{cr:run:phase5} §"Power figure caveat" bead
  \bead{0p6} note. \ok{} Current RTL HAS been re-measured on the fixed
  flow (bead \bead{0d0}, 2026-09-18): −727.37 ps / 781 violators / hold
  clean / 283.6 mW at \code{RUN_2026-09-17_20-41-11}; run
  14\textquotesingle s own figures stay unvalidated.
  \cref{cr:run:phase5}
\item
  \ok{} \textbf{Sign-off} (M12): M11 PPA recorded; docs updated
  (CLAUDE/\allowbreak{}ROADMAP/\allowbreak{}PHASE\_\allowbreak{}STATUS →
  Phase 5 complete). Indicative ASAP7 (predictive PDK). Follow-ups:
  sv2v↔RTL equivalence --- largely done 2026-09-13 (bead \bead{q7n}:
  flat \code{soc_top} 99.65 \% proven / 0 disproven; 25 of 27 modules
  proven outright via \src{tools/verif/equiv_sv2v_module.sh}, bounded
  post-reset BMC passes on \code{pmu} and \code{dma_engine}; the SoC now
  elaborates in strict LRM mode after fixing \code{soc_top.sv}
  use-before-declare and two \code{__pnr__} shim
  \texttt{parameter\ string} mismatches --- \code{soc_all} 183/183,
  re-sweep on the regenerated netlist unchanged) + CPU-macro burst-port
  LEF (\code{g0o})
\item
  \ok{} \textbf{Behavioral PMU} (Pre-Phase-6 \#5, GH epic \#98 →
  \#99/\allowbreak{}\#100/\allowbreak{}\#101, 2026-08-01):
  \src{rtl/soc/pmu.sv} --- APB4 slave + per-domain sequencer FSM
  encoding the 4 PST states from \code{phase5_soc.upf}. Sequencing order
  is \textbf{fixed by the UPF} (down:
  \texttt{ret\_\allowbreak{}save\ →\ iso\_\allowbreak{}en\ →\ gate\ clk\ →\ reset};
  up: the reverse), walked one action per cycle through an explicit
  9-state chain. PMU sits at \texttt{0x2000\_\allowbreak{}8000–8FFF} ---
  no free APB slot existed, so \code{PERIPH_LIMIT} +
  \code{AXIL_APB_LIMIT} were extended to \code{0x2000_8FFF}. Two
  \code{rv32i_clock_gate} cells gate \code{core_clk} into
  \code{u_cpu}/\code{u_gpu} (first SoC-level clock gates); isolation
  clamps live in \code{soc_top} per
  \texttt{-\allowbreak{}location\ parent}. \textbf{Outputs are
  registered (bead \bead{ktwt}, 2026-10-18)}:
  \code{clk_en}/\code{iso_en}/\code{rst_n}/\code{ret_*} are flop Q pins
  loaded from the NEXT-state decode (zero latency change), because the
  earlier combinational decode of the binary \code{dom_state_q} could
  glitch into \code{cdc_sync} D pins and the async clear of
  \code{u_cpu_pmu_rst_sync};
  \src{tools/verif/check_registered_outputs.py} proves it structurally
  (a zero-delay sim cannot). A Sky130 SoC re-harden is a follow-up.
  \textbf{Isolation handshake (bead \bead{3ogd})}: the CPU clamp control
  (\code{pmu_cpu_iso_en_sync}, \code{cpu_core_clk}) is synchronised back
  to \code{core_clk} (\code{u_cpu_iso_ack_sync}) as
  \code{pmu.cpu_iso_ack_i}; the CPU \code{DOM_PD_ISO} waits for it
  before gating the clock and asserting reset ("isolate first" now holds
  at any clock ratio; it failed at a CPU clock 3x slower than the
  fabric), and \code{DOM_PU_ISO} waits for it to fall (4-phase RZ, so a
  stale ack cannot satisfy the next power-down). \code{ISO_ACK_TIMEOUT}
  (255) bounds the wait and sets sticky \code{STATUS[3]} (dead CPU clock
  must not hang the PMU). The GPU needs none (same clock). UPF step
  order unchanged. PD timing of the new crossing is NOT validated
  (budgets added to the check SDCs, no routed run). \textbf{Sequencing
  only} --- retention save/\allowbreak{}restore are emitted in order but
  are a functional no-op (the UPF, rewritten 2026-10-17 bead
  \bead{h3lm}, declares no retention and no power switches), and no
  \code{pg_ctrl}. \code{soc_all} \textbf{120/120}, \code{test_pmu} 10/10
  with per-cycle ordering assertions.
\item
  \ok{} \textbf{2-domain multi-clock SoC + CDC} --- COMPLETE 2026-08-09,
  both domains close (Pre-Phase-6 \#4, GH epic \#90 → \#91--\#96, bead
  \bead{d8v} decision in \src{docs/3PLL_CDC_EVALUATION.md}): split the
  SoC into CPU @ \textasciitilde1282 MHz and GPU+bus+periph @ 571 MHz,
  giving exactly \textbf{one} CDC boundary (CPU ↔ crossbar M0).
  \textbf{\#91 done (2026-08-01)}: \src{rtl/soc/async_axi_fifo.sv} ---
  dual-clock AXI4 bridge, 5 gray-code FIFOs
  (AW/\allowbreak{}W/\allowbreak{}B/\allowbreak{}AR/\allowbreak{}R,
  id-less to match the crossbar\textquotesingle s master face), built on
  new reusable primitives in \src{rtl/soc/cdc/} (\code{cdc_sync},
  \code{cdc_reset_sync}, \code{cdc_gray_fifo}). Cummings style-1
  pointers: only the \emph{registered} gray pointer crosses, through one
  2-FF sync each way; \code{full_q}/\code{empty_q} registered in their
  own domains, so \textbf{no combinational path crosses the boundary}.
  Reset is cross-coupled
  (\texttt{rst\_\allowbreak{}both\_\allowbreak{}n\ =\allowbreak{}\ s\_\allowbreak{}rst\_\allowbreak{}n\_\allowbreak{}i\ \&\ m\_\allowbreak{}rst\_\allowbreak{}n\_\allowbreak{}i}
  into both \code{cdc_reset_sync}) --- a one-sided reset would otherwise
  inject fabricated AXI beats. Reset is a \textbf{hard flush, not a
  graceful abort} (a burst can be left half-transferred; both resets
  must share a common root). \code{DEPTH} must be a power of two ≥ 2
  (elaboration guard; the DEPTH=2 variant was added by bead \bead{3snb}
  2026-10-13 --- the full flag at depth 2 is
  \texttt{wr\_\allowbreak{}gray\_\allowbreak{}d\ =\allowbreak{}=\allowbreak{}\ \textasciitilde{}rd\_\allowbreak{}gray\_\allowbreak{}sync\_\allowbreak{}w}
  because a 2-bit gray pointer has no low bits for the general "invert
  the top two bits" compare; \code{async_axi_fifo} itself remains
  validated at ≥ 4 only). \code{test_async_axi_fifo} \textbf{22/22},
  \code{soc_all} \textbf{142/142}. \warn{} The \code{mem_q} →
  read-domain combinational read is a real async path --- \textbf{\#94
  must add
  \texttt{set\_\allowbreak{}max\_\allowbreak{}delay\ -\allowbreak{}datapath\_\allowbreak{}only}}
  on it and on the two gray-pointer crossings. \textbf{\#92 done
  (2026-08-01)}: second \code{pll_subsystem} instance
  (\code{u_cpu_pll_sub}) + second top-level reference port
  \code{cpu_clk_i}/\code{cpu_rst_n_i} + \code{APB_PLL2} slot at
  \texttt{0x2000\_\allowbreak{}9000–9FFF}
  (\code{PERIPH_LIMIT}/\code{AXIL_APB_LIMIT} extended).
  \code{pll_clkgen_stub} is a pure ref passthrough with no divide, so
  two PLLs off \textbf{one} reference would be bit-identical --- hence a
  second \emph{port}, not a divided clock (also avoids the Run-13 CTS
  split-tree lesson in \code{phase5_soc.sdc}). \ghref{89} \code{pll_enable}
  self-brick was already fixed in \code{110d7ac}; \#92 only replicates
  it to instance 2. \textbf{\#93 done (2026-08-02)}: CPU domain made
  real. \code{u_cpu_cg} re-sourced \code{core_clk}→\code{cpu_core_clk};
  \texttt{async\_\allowbreak{}axi\_\allowbreak{}fifo\ u\_\allowbreak{}cpu\_\allowbreak{}axi\_\allowbreak{}cdc}
  inserted on CPU↔crossbar M0; \code{pmu_cpu_clk_en}/\code{iso_en} via
  \code{cdc_sync} and \code{pmu_cpu_rst_n} via \code{cdc_reset_sync}
  into the CPU domain; \code{ext_irq}/\code{timer_irq} synchronised
  (both level-held, so 2-FF is safe); ISO\_CPU clamps moved to the
  bridge\textquotesingle s CPU-domain face. New
  \textbf{\src{rtl/soc/apb_cdc_bridge.sv}} (originally a 2-phase toggle
  handshake at \#93; converted to a 4-phase return-to-zero handshake by
  bead \bead{rfz}, 2026-08-12 --- see the file header) closes the
  APB\_PLL2 crossing that went live once
  \texttt{cpu\_\allowbreak{}clk\_\allowbreak{}i\ !=\allowbreak{}\ clk\_\allowbreak{}i}
  (\code{axil_to_apb} confirmed tolerant of arbitrary \code{pready} wait
  states). TB wrappers now expose \code{cpu_clk_i}/\code{cpu_rst_n_i} as
  real ports with a shared \code{soc_clocks.py} helper; existing suites
  stay at 1:1. \code{soc_all} \textbf{159/159} (21 suites),
  \code{test_apb_cdc_bridge} 15/15, \code{soc_multiclock} 4/4 +
  \code{soc_pll_multiclock} 1/1 at a \textbf{7 ns / 3 ns coprime ratio}
  --- the first real CDC coverage in this repo. \warn{} \textbf{Two RTL
  bugs found only at a divergent ratio}, both invisible to every 1:1
  suite \emph{and} to \code{async_axi_fifo}\textquotesingle s own 22
  unit tests, because both reset their domains together: (1)
  \code{4398c2e} --- the CPU clock-enable synchroniser reset to
  "disabled", so \code{cpu_gated_clk} was stopped for the whole reset
  window and the CPU\textquotesingle s \emph{synchronous} reset never
  loaded \code{RESET_PC} (fix: synchronise the inverted enable so
  reset-to-0 means "not disabled"); (2) \code{7a9f9d4} ---
  \code{cdc_gray_fifo}\textquotesingle s
  \texttt{wr\_\allowbreak{}ready\_\allowbreak{}o\ =\allowbreak{}\ !full\_\allowbreak{}q}
  was not gated by \code{wr_rst_n_i}, so the FIFO advertised ready while
  held in reset and silently swallowed beats (fix:
  \texttt{wr\_\allowbreak{}ready\_\allowbreak{}o\ =\allowbreak{}\ wr\_\allowbreak{}rst\_\allowbreak{}n\_\allowbreak{}i\ \&\&\ !full\_\allowbreak{}q}).
  Shared signature: \textbf{a reset value that advertises availability}.
  \#95 should gate on that class explicitly rather than trusting
  unit-suite green. \textbf{\#94 done (2026-08-02)}:
  \src{pnr/constraints/phase5_soc_multiclock.sdc} --- evolved from the
  run-14 signed-off \code{phase5_soc.sdc}, not rewritten. Two
  \code{create_clock} on the \textbf{ports} (\code{sys_clk} 1750 ps on
  \code{clk_i}, \code{cpu_clk} 780 ps on \code{cpu_clk_i}) --- the
  Run-13 CTS split-tree lesson is preserved and extended to the second
  PLL stub, so still \textbf{no \code{create_generated_clock}} for
  \code{core_clk}/\code{cpu_core_clk}. Adds
  \texttt{set\_\allowbreak{}clock\_\allowbreak{}groups\ -\allowbreak{}asynchronous}
  between the two, then
  \texttt{set\_\allowbreak{}max\_\allowbreak{}delay\ -\allowbreak{}datapath\_\allowbreak{}only}
  (bounded to one \emph{destination}-domain period) on: all 5
  \code{async_axi_fifo} channel FIFOs (\code{mem_q} combinational read +
  both gray-pointer crossings each),
  \code{apb_cdc_bridge}\textquotesingle s 2 toggle crossings + 4
  \code{cmd_*}/2 \code{resp_*} payload buses, and the PMU/IRQ single-bit
  crossings into \code{cpu_core_clk}. \code{u_cpu_pmu_rst_sync} gets
  \code{set_false_path} instead --- it is an async-clear reset path, not
  a setup-checked datapath. Also fixed the sv2v build lists
  (\code{SKY130_SOC_SV_FILES}/\code{SOC_SV_FILES} were missing all 5 CDC
  modules \textbf{and} \code{rv32i_clock_gate.sv}); both
  \texttt{make\ <pdk>-\allowbreak{}soc-\allowbreak{}sv2v} runs now
  regenerate clean. \warn{} \textbf{sv2v has no Verilator-style
  \code{-I} auto-discovery}, so a missing module there is silently
  blackboxed --- the cocotb regression can never catch it. SDC is
  Tcl-syntax-validated only, not yet run through OpenSTA and not yet
  wired into any \code{config.json} (\#96). Residual, non-blocking:
  metastability injection is not modelable in Verilator; PMU-driven live
  power-cycling across domains is untested (bead \bead{2k8}) --- the
  underlying ordering hazard (\code{pmu.sv} releases \code{rst_n_w} one
  cycle \emph{before} ungating \code{clk_en_w} on power-up) was
  \textbf{closed in \code{soc_top.sv}}, not \code{pmu.sv}, by gating
  \code{u_cpu_cg} with the CPU-domain-native
  \code{pmu_cpu_rst_n_cpu_sync}
  (\texttt{cpu\_\allowbreak{}gated\_\allowbreak{}clk\_\allowbreak{}en\ =\allowbreak{}\ \textasciitilde{}pmu\_\allowbreak{}cpu\_\allowbreak{}clk\_\allowbreak{}dis\_\allowbreak{}sync\ \&\ pmu\_\allowbreak{}cpu\_\allowbreak{}rst\_\allowbreak{}n\_\allowbreak{}cpu\_\allowbreak{}sync},
  commit \code{30beab3}): \code{pmu.sv}\textquotesingle s sequence order
  is UPF-mandated and the sync-reset CPU domain depends on it to keep
  its flop state across a PMU cycle
  (\code{docs/POWER_DOMAIN_EVALUATION.md:22,25}), so it stays unchanged.
  \warn{} That state survival is a behavioural-model artifact of the
  \textbf{CPU only} (bead \bead{4a7i}): the GPU domain uses async reset,
  so \code{pmu_gpu_rst_n} clears GPU state (IRQ latch,
  \code{irq_enable}, \code{done}) despite the gated clock (repro
  \code{tb/cocotb/soc/test_soc_integration.py::test_soc_gpu_isolation_via_pmu}).
  Neither domain models real retention (none in Phase 5); software must
  re-initialise the GPU after power-up. \code{timer_irq} is never
  exercised by any reused firmware. Also unfixed by design: the PMU
  power-down path has \textbf{no drain/\allowbreak{}quiesce step}, so a
  mid-burst CPU reset can wedge a crossbar slave engine
  (\code{soc_top.sv:84-93}) --- pre-existing, widened by the
  bridge\textquotesingle s hard-flush reset. \textbf{\#96 done
  (2026-08-09)}: multi-clock PD re-closure, 8 runs (16--23). Driven by a
  new \src{pnr/asap7/soc/config_multiclock.json} +
  \code{librelane-asap7-soc-multiclock} target, leaving
  \code{config.json} and the run-14 sign-off byte-reproducible.
  \textbf{Best and accepted: run 23} (\code{RUN_2026-08-09_14-36-16})
  --- \textbf{BOTH domains close}: \code{sys_clk} \textbf{+74.8 ps / TNS
  0 / 0 violators → 571 MHz MET}, \code{cpu_clk} \textbf{+9.8 ps / TNS 0
  / 0 violators → 1282 MHz MET}, the first time the CPU domain has ever
  met its 780 ps target. 0 routing DRC, 0 antenna, 0 slew/cap \warn{}
  all unsupported, run 23 committed zero routed wires --- bead
  \bead{xy6}, 51.9 mW \warn{} fabric-only (excludes both hard macros ---
  see bead \bead{ew3} / \cref{cr:run:phase5}; the 62.9→51.9
  mW drop vs. run 14 must NOT be read as a multi-clock power win, it
  coincides with the unrelated first-ever \code{USE_ICG_CELL}
  insertion), 66.9 \% util, and the \textbf{post-route CDC budget check
  PASS}. \warn{} Both domains\textquotesingle{} timing figures here are
  also \code{8f3}-affected STA-zero-wire results, unvalidated and
  run-23-artifacts-gone (bead \bead{0p6}); the honest current-RTL
  post-GRT point (different RTL vintage, not a run-23 re-run) is
  \textbf{setup WNS −727.37 ps / TNS −60,665 ps / 781 violators, hold
  CLEAN, 283.6 mW total / 89.6 mW fabric-only} at
  \code{RUN_2026-09-17_20-41-11} (bead \bead{0d0}, 2026-09-18;
  supersedes the earlier −2084 ps / 41,864-violator \code{w3a-sta-7}
  snapshot) --- \cref{cr:run:phase5} §"Power figure caveat".
  The CPU domain runs \textbf{2.24× the fabric} instead of being
  de-rated onto it. Getting there needed two independent fixes: (a) the
  macro exported \textasciitilde600--800 ps combinational in→out arcs on
  its APB outputs ---
  \code{apb_prdata_o}/\code{apb_pslverr_o}/\code{apb_pready_o} are now
  PURE registered outputs (two earlier attempts failed and are
  documented in \code{rv32i_cpu_top.sv}: registering \code{prdata} alone
  merely RELOCATED the delay into an 876.85 ps input setup requirement,
  and a SETUP-phase decode for \code{pslverr} still left a 617.9 ps arc
  re-sourced from \code{apb_psel_i}); and (b) the CPU macro had
  \textbf{no \code{FP_PIN_ORDER_CFG}}, so every regeneration reshuffled
  \textbf{393 of 401} boundary pins --- that, not the RTL, caused run
  22\textquotesingle s −596 ps collapse (post-CTS \code{cpu_clk} skew
  −485.6 → −922.1 ps, hold violations 805 → 4702, 4814 hold buffers
  inserted). \src{pnr/asap7/cpu/pin_order.cfg} now pins it: two runs
  produce a byte-identical LEF, which made tuning measurable for the
  first time (block setup −25.87 → −10.48 → \textbf{−0.87 ps} across
  three revisions). Full table + caveats →
  \cref{cr:run:phase5}. \warn{} \textbf{Three findings that
  outlive \#96, all measured:} (1) bead \bead{s9f} --- librelane
  declares the resizer slack margins \texttt{units=\allowbreak{}"ns"}
  but passes them raw to OpenROAD, which reads the Liberty unit; ASAP7
  is \code{1ps}, so \code{0.075}/\code{0.1} meant 0.075 ps / 0.1 ps.
  \textbf{Every ASAP7 sign-off in this repo (CPU 1282 MHz, GPU, SoC run
  14) had \textasciitilde0 optimisation guard band, not 75/100 ps.}
  Reported slacks remain valid; the guard band does not. (2) bead
  \bead{7l5} (fixed + closed) ---
  \code{check_cdc_timing.tcl}\textquotesingle s exceptions never bound
  on a flattened netlist: hierarchy survives only as dot-joined escaped
  \textbf{net} names while cells become anonymous,
  OpenSTA\textquotesingle s
  \texttt{-\allowbreak{}filter\ "name\ =\allowbreak{}\textasciitilde{}"}
  mis-parses a literal \code{.}, and the old exceptions targeted the
  synchroniser \textbf{output} (\code{q_o}), a destination-domain-only
  path rather than the crossing. Now re-pointed at the
  launch→first-capture segment and proven non-vacuous by negative
  control (all budgets forced to 1 ps → 7 violated groups). (3) beads
  \code{cyb}/\code{g0o} --- \textbf{both now FIXED in run 23}: the
  800.71 ps combinational APB-debug arc is gone (all three APB outputs
  registered; no output pin of the macro carries a combinational arc),
  and the burst-era ports (\code{axi_rlast_i}, \code{awlen},
  \code{wlast}, \code{arlen}, \ldots) are now constrained in the block
  SDC and present in both the LEF and the Liberty. A new finding
  replaces them: the macro boundary was never pinned, so any
  regeneration could silently wreck the integrating
  top\textquotesingle s timing --- fixed by \code{pin_order.cfg}, and
  the GPU macro has the same untreated exposure. IR drop was
  unobtainable until 2026-09-12 (now fixed, see below):
  \code{analyze_power_grid} failed \code{PSM-0069} --- \warn{} not from
  a tap-cell artifact, but because \textbf{no ASAP7 run ever committed a
  power grid at all} (bead \bead{4l8}): \code{pdngen} raises
  \code{PDN-0179} inside \code{build_grids}, before \code{write_to_db},
  and the local warn-and-continue patch swallowed it. Measured: 0
  \code{dbSWire} shapes on VDD/VSS in run 23 and in the 3.0.14 CPU run,
  against 65 610 in a sky130 control. Consequence: every ASAP7
  routing/\allowbreak{}DRC number was produced with M1/M2/M5 free of PDN
  metal and is a lower bound. See \cref{cr:run:phase5} §PDN
  caveat. \textbf{Fixed 2026-09-12 (bead \bead{4l8}):} the grid topology
  was wrong --- M1 rails were connected straight to M5 with M2 as a
  parallel sparse strap; the ORFS ASAP7 reference puts followpins on
  both M1 and M2, then M2→M5→M6. With \src{pnr/asap7/pdn_asap7_orfs.tcl}
  the CPU block builds 27 999/27 213 PG shapes, \code{check_power_grid}
  PASSes both nets, and IR drop reports 0.18 \%/0.19 \% worst case
  (indicative --- harness on a floorplan-stage ODB, not a full run).
  ASAP7 also shipped no via resistances, so \code{VIAS_R} had to be
  added or PSM aborts with \code{PSM-0021}.
\end{enumerate}

\textbf{Phase 6+} (after Phase 5 sign-off; Phase 6 complete 2026-10-05,
see items 13--15): additional \textbf{APB4} peripherals
(GPIO/\allowbreak{}I2C/\allowbreak{}PWM/\allowbreak{}WDT/\allowbreak{}TRNG/\allowbreak{}AES-SHA),
minimal INT8 NPU (\src{rtl/npu/}), and FreePDK45 tech-node exploration.
See \cref{cr:status:roadmap}. \warn{} Peripherals are \textbf{APB4}, not
AXI4-Lite --- \src{docs/PERIPHERAL_BUS_EVALUATION.md} settled that
("Phase 6 peripherals become trivial APB drop-ins"); the AXI4-Lite
wording that stood here and in \code{docs/ROADMAP.md:398} (\cref{cr:status:roadmap}) until
2026-09-20 was stale.

\begin{enumerate}
\def\labelenumi{\arabic{enumi}.}
\setcounter{enumi}{12}
\item
  \ok{} \textbf{Phase 6a --- GPIO controller} (bead \bead{ckc}, RTL
  2026-09-20; Sky130-hardened 2026-09-26 via bead \bead{00ef}):
  \src{rtl/periph/gpio_controller.sv}, APB4 slave on
  \code{apb4_register_bank}, \textbf{32 pins}, at
  \texttt{0x2000\_\allowbreak{}A000–AFFF} as APB index 7.
  \code{APB_N_SLAVES} 7 → 8 and both \code{AXIL_APB_LIMIT} /
  \code{PERIPH_LIMIT} extended to \code{0x2000_AFFF} --- the same
  two-edit pattern PLL/\allowbreak{}PMU/\allowbreak{}PLL2 each used.
  Pins are the unidirectional triplet \code{gpio_out_o} /
  \code{gpio_oe_o} / \code{gpio_in_i}: \textbf{this RTL tree contains no
  \code{inout} or tristate at all}, and the Sky130 SoC hardens as a core
  macro with no pad ring, so the bidirectional merge is deliberately
  left to pad-ring integration. \code{gpio_in_i} is asynchronous and is
  synchronised by \textbf{\code{N_PINS} separate
  \texttt{cdc\_\allowbreak{}sync\ \#(.\allowbreak{}WIDTH(1))}
  instances}, not one wide instance --- that primitive\textquotesingle s
  own scope warning forbids multi-bit binary buses, and GPIO pins carry
  no cross-bit word coherency. Per-pin level/\allowbreak{}edge
  interrupts with polarity. \code{GPIO_IRQ_STAT} is \textbf{edge-sticky
  / level-live}, ARM-PL061-style: an edge-mode bit latches and is
  cleared only by a \code{GPIO_IRQ_CLR} write (a \textbf{set beats a
  same-cycle clear}, so an edge is never lost to a racing clear), while
  a level-mode bit simply follows its condition and \code{GPIO_IRQ_CLR}
  has no effect on it. Making level events sticky too was rejected ---
  at the reset defaults (level, active-low, pins undriven low) every
  pin\textquotesingle s condition is true, so a fully sticky raw
  register would latch to all-ones on the first cycle out of reset.
  \code{GPIO_IRQ_EN} masks the output only, and \code{irq_o} is
  level-held as \code{soc_top.sv:623-627} requires. New
  \code{interrupt_controller} source bit 5 (\code{N_SOURCES} 5 → 6).
  \code{test_gpio} \textbf{16/16}, \code{soc_all} \textbf{217/217} over
  27 suites (201 pre-GPIO; CI \code{PASS_FLOOR} 170 → 210). Growing the
  APB window to \code{0x2000_AFFF} broke
  \code{test_axil_interconnect.test_decerr_unmapped}, which probed
  \code{0x2000_A000} expecting DECERR --- \code{BAD_HIGH} bumped to
  \code{0x2000_B000}, which is the per-growth maintenance step that
  test\textquotesingle s own comment already prescribes; expect the same
  edit for every future APB slot. Also fixed in passing:
  \code{apb_interconnect.sv}\textquotesingle s response-mux index width,
  which tripped \code{WIDTHTRUNC} once \code{N_SLAVES} hit the
  power-of-two 8 (now separate \code{SEL_W}/\code{IDXW} localparams,
  with \code{IDXW} guarded so the module\textquotesingle s own
  \code{N_SLAVES=1} default stays legal). \ok{} \textbf{Hardened on
  Sky130 2026-09-26} (\code{RUN_2026-09-26_00-07-59}, bead \bead{00ef}):
  \textbf{Netgen LVS PASSED}, routing DRC 0, \textbf{setup AND hold both
  clean at all 9 corners} (+7.08 ns / +0.1413 ns), 48.25 mW, 38.98 \%
  util, 40 MHz. Magic DRC 9,081 is the only deferred error (bead
  \bead{45a} waiver). \warn{} First Sky130 SoC harden since 2026-07-30
  and \textbf{NOT} a single-variable comparison with it --- 33 RTL
  commits intervened, 13 touching \code{soc_top.sv}. Notably \textbf{the
  ss setup failure bead \bead{ujv} was deferred on no longer
  reproduces}. KLayout DRC was \textbf{skipped, not waived} (8.5 GB peak
  OOM-killed the prior attempt on this 15 GB host, and
  \code{KLayout.DRC} precedes \code{Netgen.LVS} in
  librelane\textquotesingle s Classic flow, so the kill cost LVS for a
  check not on the path to it);
  \texttt{make\ librelane-\allowbreak{}sky130-\allowbreak{}soc-\allowbreak{}noklayout}
  does this via \code{--skip} so \code{config.json} keeps
  \texttt{RUN\_\allowbreak{}KLAYOUT\_\allowbreak{}DRC:\allowbreak{}\ true}.
  Hold went \textbf{71 → 3 → 0} by raising
  \code{GRT_RESIZER_HOLD_SLACK_MARGIN} 0.05 → 0.3 and adding
  \texttt{PL\_\allowbreak{}RESIZER\_\allowbreak{}HOLD\_\allowbreak{}SLACK\_\allowbreak{}MARGIN:\allowbreak{}\ 0.\allowbreak{}3},
  then closing the last 3 --- root cause was that all design repair runs
  on GRT-\textbf{estimated} parasitics with no repair stage after RCX,
  so real-route erosion lands where nothing can react (every violation
  was a boundary-port path; \code{timing__hold_r2r_vio__count} = 0 at
  all 9 corners). The 3 survivors were one path, \code{gpio_in_i[31]} →
  a \code{cdc_sync} stage-1 flop, closed by
  \texttt{set\_\allowbreak{}false\_\allowbreak{}path\ -\allowbreak{}from\ [get\_\allowbreak{}ports\ gpio\_\allowbreak{}in\_\allowbreak{}i]}
  (plus \code{cpu_rst_n_i}, a latent gap of the same class dating to
  2026-08-02). Open follow-ups: bead \bead{e45j} ---
  max-slew/\allowbreak{}max-cap, root-caused to the fact that
  \textbf{every} librelane repair step is pre-route and runs on
  GRT-\emph{estimated} parasitics (\code{repair_design_postgrt.tcl} uses
  \texttt{estimate\_\allowbreak{}parasitics\ -\allowbreak{}global\_\allowbreak{}routing}),
  with no post-RCX repair stage existing at all, so real-parasitic
  erosion lands where nothing can act on it (0/0 at steps 33/40/41 →
  8623/346 at step 51). Adding a real post-RCX step would mean editing
  the \textbf{shared external librelane install} every other flow here
  uses, and the \code{Odb.InsertECOBuffers} alternative targets pins by
  synthesis-generated name so any committed list goes stale --- both
  rejected. Instead, over-repairing \emph{before} routing
  (\texttt{GRT\_\allowbreak{}DESIGN\_\allowbreak{}REPAIR\_\allowbreak{}MAX\_\allowbreak{}\{SLEW,\allowbreak{}CAP\}\_\allowbreak{}PCT}
  10 → 30,
  \texttt{DESIGN\_\allowbreak{}REPAIR\_\allowbreak{}MAX\_\allowbreak{}\{SLEW,\allowbreak{}CAP\}\_\allowbreak{}PCT}
  20 → 40) took \textbf{slew 8623 → 2380 and cap 346 → 70} at
  \code{RUN_2026-09-26_07-34-03} while setup/\allowbreak{}hold slacks
  both \emph{improved} --- at the cost of \textbf{antenna 132 → 182
  nets}, the trade \code{sky130_soc.sdc:144-149} already documented. Not
  zero, so still open. Bead \code{58q} (reopened 2026-09-26) then
  recovered antenna by raising \code{GRT_ANTENNA_ITERS} 8 → 20 /
  \code{GRT_ANTENNA_MARGIN} 25 → 50 --- \textbf{antenna 263 → 182 pins,
  182 → 114 nets} (nets now better than the 120 \code{58q} originally
  accepted), hold/\allowbreak{}setup/\allowbreak{}LVS/\allowbreak{}DRC
  all still clean, at the cost of \textbf{max-slew 2380 → 3199}. \warn{}
  \textbf{These two knobs pull against each other} --- repair buffers
  cut slew and raise antenna; antenna diodes cut antenna and raise slew,
  each diode being another load pin. The Pareto frontier is demonstrated
  but uncharacterised; both current points are single-axis optima. Also
  note \code{58q}\textquotesingle s old close reason contains a
  \textbf{false premise} --- it claims GRT antenna repair skips the
  legalization sub-steps, but \code{antenna_repair.tcl:23-27} sources
  \code{dpl.tcl} and \code{grt.tcl} after \code{repair_antennas}; the
  lever was safe only because \code{OpenROAD.RepairAntennas} already
  runs in the adopted flow with hold clean. Both checkers remain
  \code{--skip}-ped, and "max fanout = 0" is \textbf{vacuous} because
  the SDC has no \code{set_max_fanout} at all. \ok{} \textbf{UPDATE
  2026-10-07 (bead \bead{e45j}, \cref{cr:defect:sky130-slew-cap-e45j}):
  root cause corrected and largely fixed.} The "0/0/0 at steps 33/40/41"
  figures were \textbf{nom\_tt only} (mid-PnR STA runs one corner; other
  corners\textquotesingle{} \code{state_out} values are stale step-9
  copies), and the resizer \emph{is} corner-aware: an in-session GRT
  estimate on the step-40 ODB already gave max\_ss slew 307 / cap 0. The
  gap is the \textbf{estimator} --- sky130\textquotesingle s tech-LEF
  layer C is ground-only, and OpenRCX extracts \textbf{1.69 / 1.88 /
  2.04× the estimated C (min/\allowbreak{}nom/\allowbreak{}max) and 0.62
  / 1.02 / 1.91× the R} (17 525 nets). A \textbf{project-local LibreLane
  plugin} (\src{pnr/sky130/soc/plugin/librelane_plugin_cvt_sky130}, no
  shared-install edit) rescales the estimate per corner in the four
  repair/\allowbreak{}resizer steps;
  \texttt{make\ -\allowbreak{}C\ pnr\ librelane-\allowbreak{}sky130-\allowbreak{}soc-\allowbreak{}rccal-\allowbreak{}noklayout}
  (\code{config_rccal.json}). Measured vs Gate B baseline
  \code{RUN_2026-10-05_06-50-11} (one run,
  \code{RUN_2026-10-06_22-43-34}, \code{dma_engine.sv} pinned to its
  Gate B version so only one variable changed): \textbf{max\_ss slew
  7622 → 295, cap 140 → 0 (cap 0 at all 9 corners)}, antenna
  \textbf{210/305 → 184/259 nets/pins (better, not worse)}, setup worst
  slack +0.579 → \textbf{+3.276 ns}, hold +0.279 → +0.275 ns, LVS
  PASSED, routing DRC 0, Magic 8984 unchanged, util 44.8 → 45.2 \%. This
  broke the slew-vs-antenna trade-off. \code{e45j} stays open (295
  residual: 16 SRAM-input pins = \code{o1i}, \textasciitilde33 buffer
  nets, diode pins); \code{Checker.MaxCapViolations} could now be
  un-skipped but was not executed this run. \ok{} \textbf{UPDATE
  2026-10-09 (bead \bead{e45j}, \cref{cr:defect:sky130-slew-cap-e45j:7-adoption-2026-10-09-rccal-is-the-default-max-cap-gated-with-a-named-waiver}):
  the RC-calibrated flow is now the DEFAULT Sky130 SoC flow}
  (\texttt{make\ -\allowbreak{}C\ pnr\ librelane-\allowbreak{}sky130-\allowbreak{}soc[-\allowbreak{}noklayout]}
  runs \code{config_rccal.json}; \code{-norccal[-noklayout]} runs the
  old \code{config.json} for A/B; \code{SKY130_SOC_CONFIG=} is honoured
  by every target). Max-cap is gated at all 9 corners by the plugin step
  \code{CVT.MaxCapViolations} with \textbf{one named, user-approved
  single-net waiver} in
  \src{pnr/sky130/soc/constraints/max_cap_waivers.json}: the load is the
  CPU macro pin \code{u_cpu*/axi_rdata_i[12]} (Liberty pin cap 0.2464 pF
  vs 0.0083 pF median; net 0.56-0.68 pF vs 0.5301 pF limit at the six
  tt/ss corners), keyed on the macro pin (not the stale-prone driver
  name \code{_078722_}), with a 0.70 pF ceiling and a stale-waiver
  warning. \warn{} Un-skipping the stock \code{Checker.MaxCapViolations}
  alone gates NOTHING (its
  \texttt{corner\_\allowbreak{}override=\allowbreak{}[""]} matches no
  corner; measured exit 0 + a WARNING), which is why the step is custom;
  the stock checker cannot waive by net. Validated by 50 unit tests
  (\src{tb/tests/test_cvt_maxcap_waiver.py}) plus the check applied to
  the finished runs: PASS with the waiver, FAIL with an empty list; the
  substituted step was exercised in LibreLane (\code{--only}) but
  \textbf{no full flow has run with the adopted config yet}. Re-check
  any run with
  \texttt{make\ -\allowbreak{}C\ pnr\ sky130-\allowbreak{}soc-\allowbreak{}maxcap-\allowbreak{}check}.
  Max-slew stays skipped (102-423 pins at max\_ss residual);
  \code{GRT_DESIGN_REPAIR_MAX_WIRE_LENGTH=1500} is NOT adopted (needs
  another RTL-vintage arm); the source fix (an input buffer inside the
  CPU macro) still needs hard-macro approval. \warn{} \textbf{UPDATE
  2026-10-08 (bead \bead{e45j}, \cref{cr:defect:sky130-slew-cap-e45j:5-confirmation-run-on-current-rtl-2026-10-08-not-adopted}):
  confirmation on CURRENT RTL did NOT hold, so rccal is NOT adopted as
  default and \code{Checker.MaxCapViolations} stays skipped.}
  \code{RUN_2026-10-08_05-21-39} (rccal, no DMA pin): LVS PASSED,
  routing DRC 0, setup+hold 0 violators (+3.329 / +0.282 ns), but
  \textbf{max-cap = 1 at six tt/ss corners} (one driver, \code{_078722_}
  o22a\_4 on \code{cpu_bridge_s_rdata[12]}, 0.614 pF vs 0.530 pF at
  nom\_tt), max\_ss slew 295 → 423, antenna 184/259 → 205/273 nets/pins
  (different netlist, one run per arm, noise uncharacterised; still
  \textasciitilde18x better than Gate B on max\_ss slew). Un-skipping
  would fail the flow. \ok{} \textbf{UPDATE 2026-10-09 (bead
  \bead{e45j}, §6): still NOT adopted; noise is zero.} A same-config
  repeat (\code{RUN_2026-10-08_21-07-20}) is \textbf{bit-identical} to
  the confirmation run through step 51 (deterministic flow, so
  run-to-run noise = 0). The one cap violation is a long wire (2 196 um,
  rank 54/117 008) into a CPU-macro input pin whose Liberty cap is
  \textbf{0.2464 pF vs a 0.0083 pF median} (\code{axi_rdata_i[12]}) ---
  not an RC-estimator error. Single-variable run
  \code{GRT_DESIGN_REPAIR_MAX_WIRE_LENGTH=1500}
  (\code{config_rccal_wl1500.json}, \code{RUN_2026-10-08_23-12-19}):
  max\_ss slew \textbf{423 → 102}, antenna 205/273 → 212/253,
  setup/\allowbreak{}hold/\allowbreak{}LVS/\allowbreak{}DRC unchanged,
  but \textbf{cap still 1 at the same six corners} (0.646 pF), so no
  default changed and \code{Checker.MaxCapViolations} stays skipped.
  Remaining fix is macro-side (CPU macro input buffering, needs
  approval) or a named single-net waiver (user\textquotesingle s call).
  \ok{} \textbf{UPDATE 2026-10-28 (bead \bead{e45j},
  \cref{cr:defect:sky130-slew-cap-e45j:9-wire-length-1500-adopted-as-default-2026-10-28-bead-e45j}):
  \code{GRT_DESIGN_REPAIR_MAX_WIRE_LENGTH=1500} IS now adopted (in
  \src{pnr/sky130/soc/config_rccal.json}, the default), which SUPERSEDES
  the "NOT adopted (needs another RTL-vintage arm)" above (kept as
  history).} Single-variable arm \code{RUN_2026-10-28_05-14-30} vs
  \code{RUN_2026-10-27_05-40-27} (identical netlist through floorplan):
  max\_ss max-slew \textbf{1067 -\textgreater{} 129} (max\_tt 695
  -\textgreater{} 1, max\_ff 596 -\textgreater{} 0), max\_ss setup WNS
  -0.153 -\textgreater{} -0.055 ns (the 5 accepted NPU SRAM paths, bead
  \bead{7dhp}), hold 0 violators, LVS 0, routing DRC 0, Magic DRC 10 962
  vs 10 959 (same rules), CDC check PASS at 9 corners.
  \code{config_rccal_wl1500.json} removed; \code{hold-dlygate-scd-link}
  waiver ceiling 0.115 -\textgreater{} 0.120 pF, the two run-pinned
  waivers deleted (\code{--strict-stale} on the arm: 9 waived, 0
  unwaived, 0 stale). \code{check_scan_chains.py} needs
  \texttt{-\allowbreak{}-\allowbreak{}en-\allowbreak{}net\ dft\_\allowbreak{}scan\_\allowbreak{}en}
  or it false-fails on the resizer clone of the enable gate.
\item
  \ok{} \textbf{Phase 6b --- CRYPTO: AES-128 + SHA-256 accelerator}
  (bead \bead{f7vs.10}, 2026-10-03): \src{rtl/periph/crypto_accel.sv} +
  \code{aes128_core.sv} + \code{sha256_core.sv} --- two cores,
  \textbf{one} APB4 slot, \textbf{one} IRQ, at
  \texttt{0x2000\_\allowbreak{}F000–FFFF} as APB index 12 / IRQ bit 10,
  with
  \texttt{parameter\ bit\ EN\_\allowbreak{}AES\ =\allowbreak{}\ 1,\allowbreak{}\ EN\_\allowbreak{}SHA\ =\allowbreak{}\ 1}
  generate guards so a PD run can drop either
  (\texttt{EN\_\allowbreak{}AES=\allowbreak{}0\ \&\&\ EN\_\allowbreak{}SHA=\allowbreak{}0}
  is an elaboration error, not a config). \code{APB_N_SLAVES} 12 → 13;
  the window limits were already final at \code{0x2001_0FFF} from
  groundwork \code{f7vs.2}, so unlike every 6a item \textbf{neither
  \code{AXIL_APB_LIMIT} nor \code{PERIPH_LIMIT} nor \code{BAD_HIGH}
  moved} --- \code{f7vs.2}/\code{f7vs.3} had pre-declared
  \code{crypto_irq} and pre-placed it at \code{irq_src_i} bit 10,
  reducing integration to one tie deletion plus the
  array/\allowbreak{}count edits. \textbf{AES is encrypt-only, exposing
  ECB \emph{and} CTR}: CTR makes an encrypt-only core a complete cipher
  in both directions, so there is no inverse S-box and no inverse key
  schedule --- nearly a second core avoided. Iterative 128-bit datapath,
  \textbf{on-the-fly key schedule in a single 128-bit \code{rk_q}}; the
  precomputed-11-round-key form was rejected twice over, at 1408 flops
  \emph{and} because fetching round key \code{r} is exactly the
  runtime-indexed read the \code{ma7} defect class forbids.
  \code{SBOX_PARALLEL=16} → \textbf{11 cycles/\allowbreak{}block}; the
  documented fold-to-4 fallback is \textbf{41 cycles and bit-identical},
  folding via a byte-lane carousel that rotates the state through a
  fixed substitution window (\code{s_q[31:0]} is always the substituted
  slice) so it introduces \textbf{no} runtime-indexed mux and stays a
  parameter change. SHA-256 is a single 512-bit block compress, 64
  rounds at 1/cycle → \textbf{66 cycles/\allowbreak{}block}; padding,
  length and multi-block iteration are
  \textbf{software\textquotesingle s job} per FIPS 180-4, and
  \texttt{CTRL[4]\ SHA\_\allowbreak{}CONT} chains from the RO
  \code{DIGEST} words so chaining costs no 256-bit register. \textbf{No
  new runtime-indexed mux} (the \code{ma7} acceptance criterion): S-box
  and SHA \code{K} are constant-label \code{case} functions, Rcon is an
  \code{xtime()} chain, the \code{W} schedule a shift register ---
  verified by \texttt{synth\ -\allowbreak{}flatten} showing \textbf{no
  \texttt{\$shiftx} and no \texttt{\$pmux}}. \warn{} \textbf{That
  verification was vacuous} (found 2026-10-05, bead \bead{f7vs.15}): a
  full \code{synth} runs \code{techmap}, which lowers every
  \texttt{\$shiftx}/\texttt{\$pmux}, so the post-synth count reads zero
  for any design. The coarse-grain netlist
  (\texttt{synth\ -\allowbreak{}run\ :\allowbreak{}fine}) of
  \code{crypto_accel} actually carries \textbf{4 \texttt{\$shiftx} + 2
  \texttt{\$shift}} (all inside \code{apb4_register_bank}, present in
  every netlist), \textbf{26 \texttt{\$pmux}} (yosys \code{proc} lowers
  every \code{case}, including constant-label ones) and \textbf{21
  \texttt{\$mem\_\allowbreak{}v2}} (the S-boxes, turned into ROMs by
  \code{proc_rom}). None of this is new runtime indexing in the CRYPTO
  datapath, and Sky130 uses the sv2v frontend, not Synlig, so the
  \code{ma7} \code{OPT_MUXTREE} defect is not exposed --- but the stated
  proof did not prove it. Use the per-module, coarse-grain method item
  15 established. \code{N_REGS=32}. \textbf{Keys live in the
  repo\textquotesingle s first shadow register outside
  \code{apb4_register_bank}}, loaded by APB write-snoop with the
  bank\textquotesingle s KEY words at \code{WMASK=0} and never
  HW-written, so they read 0 forever --- a true write-only register is
  otherwise impossible because \code{prdata} is combinational off
  \code{regs[addr_word]} (\code{apb4_register_bank.sv:120-129}). Four
  register-map entries from the plan\textquotesingle s own table proved
  \textbf{unimplementable as written} and were corrected in the golden
  spec: \code{CTRL[2]} start cannot be a stored RW bit (a same-cycle
  \code{hw_wen} clear is the §3 one-cycle window, and the
  bead-\code{6o8w} collision rule makes the SW \code{1} win anyway) so
  it is \textbf{W1P, reads 0}, decoded as a snoop pulse; \code{DIN0-3}
  cannot hold a SHA block (4 words vs 512 bits) so they are a
  \textbf{4-word aperture onto one 512-bit shift register} where
  \emph{write order}, not address, is the contract; multi-block SHA
  chaining was impossible with an RO \code{DIGEST} and no \code{H}
  input, hence \code{SHA_CONT}; and
  \texttt{STATUS[3]\ key\_\allowbreak{}write\_\allowbreak{}rejected} had
  no clear, now reset / \code{IRQ_CLR[1]} / an accepted key write.
  \textbf{DOUT and DIGEST have no result register at all} --- the bank
  word \emph{is} the result register, written from the
  cores\textquotesingle{} combinational next-state outputs on the
  completion edge, so they are either current or one operation stale,
  never a partial round state. An illegal start (\code{mode==3}, no key,
  or a compiled-out core) takes a \textbf{2-cycle zero-length path} that
  sets \code{done} with no writeback, so a \code{while(!done);} driver
  cannot hang --- but \textbf{\code{done} does not imply valid data}.
  \code{test_crypto} \textbf{41/41} (FIPS-197 App. B/C + SP 800-38A
  F.5.1 vectors, double-pinned in model \emph{and} suite) plus a
  \textbf{25-fault mutation campaign} leaving only 2 provably-equivalent
  survivors; \code{test_soc_crypto} \textbf{1/1} driving the block from
  CPU firmware over AXI-Lite→APB and proving
  \texttt{irq\_\allowbreak{}src\_\allowbreak{}i\ =\allowbreak{}=\allowbreak{}\ crypto\_\allowbreak{}irq\ <<\ 10}
  exclusively, per cycle. \ok{} \textbf{Gate A ran} (first of the phase,
  plan §8): 24,227 cells / 277,875 µm², \textbf{+14.04 ns setup slack at
  ss on the real 25 ns period} → \code{SBOX_PARALLEL} stays 16. \warn{}
  The originally recorded "\textbf{3.59 \%} of SoC stdcell area" was
  \textbf{wrong by \textasciitilde10×} (corrected 2026-10-04, bead
  \bead{f7vs.11}): its 7,744,600 µm² denominator is
  \code{design__instance__area}, which is stdcell \textbf{plus}
  hard-macro area --- macros are 6,952,440 µm² of it (88 \%). Against
  stdcell alone CRYPTO is \textbf{\textasciitilde28 \%} of the
  SoC\textquotesingle s 992,638 µm². The decision survives on a
  different reason --- core utilisation goes 39.0 \% →
  \textasciitilde43.7 \% with CRYPTO + NPU, and the fold-to-4 fallback
  would save only 0.2 pp of it for 3.7× the latency --- but "an order of
  magnitude away from 40 \%" was false. Plan §8 holds the full
  correction. \warn{} \textbf{NOT production cryptography} and the
  header says so outright: no DPA/\allowbreak{}side-channel resistance
  (duration is fixed \emph{by construction}, but that is a cycle-count
  statement only --- no claim about power, EM or glitch), no fault
  hardening, not certified. The key is \textbf{not} bus-readable, but
  that buys little: no privilege split exists
  (M0/\allowbreak{}M1/\allowbreak{}M2/\allowbreak{}M3 all decode to the
  APB ring and \code{axil_to_apb} ignores \code{AxPROT}), so \textbf{any
  master can substitute its own key between a victim\textquotesingle s
  load and start and read \code{DOUT}} --- in CTR that recovers the
  victim\textquotesingle s plaintext, in SHA it yields
  \code{H(victim_msg)}. Treat it as an unprotected shared oracle, not
  key storage. \code{key_valid} drops only on a whole-SoC reset or on
  the DFT scan wipe (bead \bead{j41m.6}, 2026-10-17: every
  \code{scan_mode} edge, entry and exit, hardware-zeroizes the key,
  round/SHA state, message shadow and
  DOUT/\allowbreak{}DIGEST/\allowbreak{}IV in one cycle, 3 edges after
  the change; the clear is functional D logic, so it needs 3
  \code{scan_en=0} edges first and protects against software and a
  cooperating tester, not an adversarial one: \code{DFT_ARCHITECTURE.md}
  14.10), \code{rk_q} retains the round-10 key (which inverts to the
  master key) between operations, and \textbf{everything else is exposed
  through scan}, so DFT access must be treated as key access.
\item
  \ok{} \textbf{Phase 6c --- INT8 NPU} (bead \bead{f7vs.11},
  2026-10-04): \src{rtl/npu/npu_top.sv} + \code{npu_mac_array.sv} +
  \code{npu_weight_mem.sv} at \texttt{0x2001\_\allowbreak{}0000–0FFF} as
  \textbf{APB index 13 / IRQ bit 11} --- the last Phase 6 item; Phase 6
  RTL is complete. Weight-stationary 4×4 INT8 grid, INT8×INT8→INT32,
  shared requantizer (unsigned ×\code{SCALE[15:0]}, \textbf{arithmetic}
  \texttt{>>>\ SCALE[20:\allowbreak{}16]} so rounding is \textbf{floor},
  49-bit product never truncated, saturate then ReLU) draining one
  lane/\allowbreak{}cycle; PIO only (\code{WADDR} auto-increment +
  \code{WDATA}, depth-4 \code{AIN}/\code{AOUT}), no SRAM aperture.
  Integration needed \textbf{no} window move ---
  \code{AXIL_APB_LIMIT}/\code{PERIPH_LIMIT}/\code{BAD_HIGH} were final
  since \code{f7vs.2}, and \code{npu_irq} was pre-declared and
  pre-placed in the \code{irq_src_i} MSB, so the IRQ edit was one tie
  deletion. \textbf{The weight-SRAM / OOM answer: every \code{ifdef} arm
  is a hard macro and there is no behavioural array anywhere in the RTL}
  --- 1× \code{sky130_sram_4kbyte_1rw1r_32x1024_8} (\code{SRAM_SKY130}),
  4× \code{sram_1rw_256x32_asap7} + \code{rv32i_clock_gate}
  (\code{SRAM_ASAP7}), 4× \code{sram_1rw_256x32_freepdk45}
  (default/\allowbreak{}cocotb). The \code{sram_controller.sv} 2-way
  shape (macro / \code{else} flat array) was \textbf{rejected}: its
  \code{else} arm is the 32 768-flop structure bead \bead{rvb} proved
  makes the ASAP7 post-CTS resizer non-terminating. It cost no new PD
  asset --- \code{MACROS} is keyed by \textbf{module, not instance}, and
  all four 4 KB views are committed at \src{pnr/sky130/soc/macro/} (the
  stub/\allowbreak{}model comments pointing at
  \code{/nobackup/openram_sky130_4kb} were wrong; that dir was wiped
  2026-09-15 and nothing reads it). Two traps it did cost: (1) the macro
  instance \textbf{must be named \code{u_sram_macro}} so the existing
  \code{PDN_MACRO_CONNECTIONS} regex matches --- any other name silently
  floats vccd1/\allowbreak{}vssd1; (2)
  \textbf{\code{--define=SRAM_ASAP7} had to be added to
  \code{SOC_SV2V_DEFINES}} because sv2v resolves \code{ifdef} before
  Yosys, so \code{config.json}\textquotesingle s define is too late ---
  without it the ASAP7 netlist instantiates the FreePDK45 module and
  silently blackboxes it (proven fixed: regenerated ASAP7 netlist has 0
  \code{sram_1rw_256x32_freepdk45}). Read latency is normalised to
  \textbf{2 cycles on every arm} in \code{npu_weight_mem.sv}, because
  the sim models disagree (Sky130 negedge-launched vs FreePDK45
  posedge). \textbf{\code{ma7} gate proven on the coarse-grain netlist}
  (\texttt{synth\ -\allowbreak{}run\ :\allowbreak{}fine}, before techmap
  hides the cells): zero
  \texttt{\$shiftx}/\texttt{\$pmux}/\texttt{\$mem*} in the three
  \src{rtl/npu} modules. Two lessons: the flat count can never be zero
  because \code{apb4_register_bank} carries 4 \texttt{\$shiftx} + 2
  \texttt{\$shift} of its own, and \textbf{a constant-label \code{case}
  is NOT sufficient} --- yosys \code{proc} lowers every \code{case} to
  \texttt{\$pmux} --- so runtime selects are 2:1 ternary trees / one-hot
  AND-OR, and AIN/AOUT are \textbf{positional shift registers}, never
  \code{array[runtime_ptr]}. \code{KLEN} is \textbf{1..63} (the
  plan\textquotesingle s "1..64" was unreachable in 6 bits; caught by
  tests-first). \code{test_npu} \textbf{39/39}, \code{test_soc_npu}
  \textbf{1/1} (CPU firmware over AXI-Lite→APB, interrupt-driven,
  \texttt{irq\_\allowbreak{}src\_\allowbreak{}i\ =\allowbreak{}=\allowbreak{}\ npu\_\allowbreak{}irq\ <<\ 11}
  proven every cycle, six negative controls), mutation campaign 65
  mutants: 55 killed, 10 provably equivalent, 0 gaps; \code{soc_all_ci}
  from clean \textbf{476 pass / 0 fail / 16 skip} over 53 suites (CI
  \code{PASS_FLOOR} 433 → 473). \textbf{Gate A}: 19 603 cells /
  \textbf{203 179 µm²} stdcell ≈ 20 \% of the SoC\textquotesingle s
  stdcell area, plus the 472 439 µm² macro; \textbf{+1.84 ns setup at ss
  on 25 ns} --- thin, limiter is the MAC array\textquotesingle s
  single-cycle multiply + column tree + accumulate from a fanout-129 AIN
  flop; Gate B fallback is registering the column sums (4×18 flops, +1
  cycle). \warn{} Same slack is \textbf{−5.0 ns} under plain
  \texttt{abc\ -\allowbreak{}liberty} --- Gate A is recipe-sensitive,
  use the flow-faithful one. \warn{} \textbf{This probe exposed that
  CRYPTO\textquotesingle s Gate A "3.59 \%" was wrong
  \textasciitilde10×} (stdcell+macro denominator) --- corrected in item
  14 and plan §8. Gate B watch items: the MAC path, the second
  macro\textquotesingle s \code{macro_placement.cfg} line (6.9 µm GCell
  multiples, path read from a synth probe not guessed), and \code{e45j}
  on \textasciitilde48 \% more stdcell logic. Non-goals: INT4, sparsity,
  tiling, 16 KB, demo firmware, L3 driver.
\item
  \wip{} \textbf{DFT --- scan + MBIST + JTAG for the Sky130 SoC} (epic
  \code{j41m}, \ghref{244}; architecture and decisions 1--9 in
  \src{docs/design/DFT_ARCHITECTURE.md}). Stage 0 \ok{} (\prref{250}).
  Stage 1a: scan-ready SoC fabric merged (\prref{258}: test ports on
  \code{soc_top}, \code{test_en} on the clock gates, reset and clock
  scan muxes, crypto key mask; seam module
  \src{rtl/soc/dft/dft_ctrl_ports.sv}), but bead \bead{j41m.2} is still
  open (◐) and the GPU scan-reset test gap \code{j41m.9} is open. Stage
  1b \ok{}: IEEE 1149.1 JTAG TAP (\src{rtl/soc/dft/jtag_tap.sv}, \prref{270}, \code{j41m.3}), plus the TRST reset-sync structural check (\prref{280}, \code{j41m.11}) and \code{dft-scan-check} for both PDKs in CI
  (\prref{281}, \code{j41m.12}). Stage 3 ATPG: the patched Quaigh completes
  the fabric (99.95 \% test / 99.16 \% fault coverage on the
  generic-gate collapsed universe, \code{j41m.18}). \textbf{99 \% / 98
  \% is adopted as the target, provisionally} (\code{j41m.23},
  2026-10-30): the fabric and the CPU macro are separate rungs, and
  sign-off waits on the cell-pin re-grade \code{j41m.20}
  (\code{DFT_ARCHITECTURE.md} 2.1, 18.8). Open: Q-TRST (\code{trst_ni}
  power-up tie-off at the chip top), Q-STUB (1-bit SCAN stub until Stage
  2 chains exist), Stage 2 scan insertion and scan-mode SDC
  (\code{j41m.8}, \code{j41m.5}), the CPU-macro scan boundary and
  re-harden (\code{j41m.4}, which depends on the \code{dud4} outcome),
  the crypto scan-entry scrub protocol (\code{j41m.6}), and the MBIST
  control-word crossing (\code{j41m.10}). Stage 4 MBIST (beads
  \code{j41m.14}/\code{j41m.10}, \prref{306}, spec
  \src{docs/design/DFT_ARCHITECTURE.md} section 16): \code{mbist_ctrl}
  (ROM-microcode March C- / retention / checkerboard / byte-lane, 6
  backgrounds), \code{mem_collar} on the two SoC-level SRAM macros,
  \code{MBIST_CTRL}/\code{MBIST_STATUS} TAP registers (opcodes 6/7) and
  the TCK -\textgreater{} core crossing through a depth-2
  \code{cdc_gray_fifo}; fault-injection suites prove detection per fault
  class against an independent golden model. The ten CPU-macro cache
  SRAMs are designed (section 16.7) but NOT wired: they need the CPU
  macro boundary change and a re-harden (proposed with \code{j41m.4}).
\item
  \wip{} \textbf{FPGA prototyping} (epic \code{hlrm}, \ghref{245}). Stage 0
  \ok{} (no board and no vendor toolchain on the host): one \code{FPGA}
  define selects block-RAM-inferable memories (\src{rtl/fpga/};
  \texttt{make\ -\allowbreak{}C\ tb/\allowbreak{}cocotb/\allowbreak{}soc\ FPGA=\allowbreak{}1\ soc\_\allowbreak{}all\_\allowbreak{}ci}
  reports 762 passed, 0 failed, with the SRAM backdoor helper, \prref{282}); yosys resource counts for Xilinx 7-series and ECP5, with and
  without GPU, are in \src{docs/FPGA_STAGE0_EVALUATION.md} (\prref{279}).
  Board recommendation, not a purchase: Digilent Nexys Video XC7A200T
  (\src{docs/FPGA_BOARD_SELECTION.md}, \prref{283}). The only measured fmax
  is an open-tool ECP5 \textbf{lower bound}, CPU-only on LFE5U-85F:
  \textasciitilde48--52 MHz (CPU) and \textasciitilde47--55 MHz
  (fabric); the with-GPU build does not place on that device. A Vivado
  Artix-7 run is in progress on a separate branch and has no results in
  this repo yet. Board-gated stages (bring-up, firmware validation,
  dual-oscillator CDC) and the clock-gating choice are open human
  decisions.
\end{enumerate}

\end{record}


% Moved verbatim from CLAUDE.md by tools/docs/md2tex.py; edit this file, not the original.
\chapter{AI/Human Boundaries}\label{cr:status:ai-human-boundaries}

\section{AI MAY assist with}\label{cr:status:ai-human-boundaries:ai-may-assist-with}

\begin{record}

\begin{itemize}
\tightlist
\item
  Python boilerplate (class structure, imports)
\item
  Simple instruction implementations (after human verification)
\item
  cocotb driver scaffolding; test case generation; documentation
  formatting
\end{itemize}

\end{record}

\section{Human MUST}\label{cr:status:ai-human-boundaries:human-must}

\begin{record}

\begin{itemize}
\tightlist
\item
  Write and approve all specifications
\item
  Implement complex instructions (branches, loads, stores); design
  control FSMs
\item
  Define verification strategy; review all AI-generated code; make
  architectural decisions
\end{itemize}

See each phase in \src{docs/verification/VERIFICATION_PLAN.md} for
detailed AI/Human responsibilities.

\end{record}
